\section{Fourier coefficient fits}\label{sec:pf:fit}

\providecommand{\fitFUN}{\sfname{FUN}}
\providecommand{\fitmlog}{\operatorname{meanLog}}
\providecommand{\fitmi}[1]{\langle #1\rangle_{\infty}}
\providecommand{\fitht}{\operatorname{ht}}
\providecommand{\fitwt}{\operatorname{wt}}
\providecommand{\fitden}{\mathfrak{a}}
\providecommand{\fitClosed}{\sfname{Closed}}
\providecommand{\fitEv}{\mathrm{Ev}}
\providecommand{\fitPC}{\mathrm{PC}}
\providecommand{\fitRR}{\mathit{RR}}
\providecommand{\fitmumax}{\mu_{\max}}
\providecommand{\fitmumin}{\mu_{\min}}
\providecommand{\fitbudget}{\mathcal{G}}
\providecommand{\fitpol}{\operatorname{pol}}
\providecommand{\fitSq}{\mathrm{Sq}}
\providecommand{\fitGBig}{\sfname{GBig}}
\providecommand{\fitBigC}{\sfname{BigC}}
\providecommand{\fitFit}{\mathrm{Fit}}

This section proves the second half of \cref{prop:families:points}: the bundle $F=\sfname{ThetaLocalInputs}$ of
classical results, two components of the bundle $P$ (the statement $\sfname{JoinOut}$ and the functional equation of
Hecke $L$-functions of CM fields), the two statements of \cref{thm:theorem-K}, the separation statement
$\sfname{FunctionalsSeparate}$ of \cref{prop:families:fs} and the low-weight statement $\sfname{LambdaLowerBound}_2$
imply $\sfname{ConstantFunctionFits}$: closedly on the class $\sfname{FamHyp}$, at the points of the families which
carry a grid at a large level $C_F$, the glued lattice $\fitRR$ of the first side $S_{\mathrm f}$ of the point
satisfies the slope inequality $\fitFUN$,
\begin{equation}\label{eq:pf:fit:FUN}
\fitmumax(\fitRR)\ \le\ x^{C_H}b-\fitmlog(\kappa),
\end{equation}
for every value of the constant $C$. (The name of the statement refers to the existence of a ``fitted function''
attached to a grid, see \cref{sec:pf:fit:assembly}; the conclusion is an Arakelov upper bound for a maximal
slope.)

\paragraph{The idea.}
The bound \eqref{eq:pf:fit:FUN} is obtained from the slope inequality of Bost \cite{Bost1996,Bost2001}. One
constructs an injective $K$-linear map $\Psi$ from the space $K^D$ of $\fitRR$ into a space $K^n$ which carries a
``test lattice'' $M$ of controlled maximal slope, such that $\Psi(\fitRR)$ is contained in $\fitden M$ for an
explicit fractional ideal $\fitden=\prod_{w\in T}w^{-a(w)}$ (the \emph{denominators}) and such that $\Psi$ has
controlled height at the infinite places. The slope inequality then gives
\[
\fitmumax(\fitRR)\le-\fitmlog(\kappa)+\fitmumax(M)+h(\Psi)+2\sum_{w\in T}\ell_wa(w),
\]
and the three terms on the right are bounded by explicit expressions whose sum is at most $x^{C_H}b$. The map $\Psi$
is a \emph{fitted function}: its coordinates on the base type are class sums $v_K(t,x,\cdot)$, finite sums over the
classes of a definite unitary group, which by the row formula $\sfname{FormulaK}$ are, up to a universal constant
and roots of unity, the values at the theta vectors of the point of separating linear functionals (pairings with
column theta functions, which play the role of the Fourier coefficients of the name). On the other types $t$ the map is related to the base type through
ratios $r_t(i,j)$ of constants of matched ``square'' theta integrals; the denominators $a(w)$ are bounded by the
valuations of these ratios, which are controlled through Hecke $L$-values at integers, toric values (CM periods) and
norms of bottom theta pairs (Rallis' inner product formula), and the heights are bounded by size estimates for class
sums at all places.

\paragraph{Organisation.}
\Cref{sec:pf:fit:setting} recalls slopes, indefinite data, grids and the statement, and states the main theorem.
\Cref{sec:pf:fit:slopes} proves a slope inequality with denominators. \Cref{sec:pf:fit:grid} proves the
``grid estimate'', which sums local bounds for the valuations of the ratios $r_t(i,j)$ over the finite places.
\Cref{sec:pf:fit:cond} introduces functional data and proves the inequality $\fitFUN$ under seventeen ``leaf
conditions'' on a functional datum (\cref{pf:fit:flb}); this is the numerical heart of the section.
\Cref{sec:pf:fit:assembly} lifts this to families and reduces $\sfname{ConstantFunctionFits}$ to sixteen family
statements (\cref{pf:fit:leaves}). \Cref{sec:pf:fit:inputs} collects the inputs (estimates for the ``square'' terms,
size bounds for class sums, the row properties), \cref{sec:pf:fit:leafproofs} proves the sixteen statements from them and
\cref{sec:pf:fit:main} proves the main theorem. The remaining subsections prove the principal inputs:
\cref{sec:pf:fit:normbig,sec:pf:fit:botconst,sec:pf:fit:betaarch,sec:pf:fit:archnorm} the square estimates (norms of bottom
pairs through the Rallis inner product formula, bottom constants, the archimedean bound for normalised toric values through
complex multiplication and the Chowla--Selberg formula, and archimedean bounds for norms of pairs),
\cref{sec:pf:fit:sizes,sec:pf:fit:factors} the size bounds for class sums, and \cref{sec:pf:fit:testlat} the slope bound for
test lattices. \Cref{sec:pf:fit:remarks} contains remarks on the formal statements.

Throughout, ``closedly'', ``eventually'' and the class $\sfname{FamHyp}$ have the meaning of
\cref{def:families:closed} and \cref{sec:families:nonempty}; as noted after \cref{def:families:closed}, closed
statements say nothing about classes without large members.

% ---------------------------------------------------------------------------
\subsection{Slopes, indefinite data and the statement}\label{sec:pf:fit:setting}

\paragraph{Places, valuations and heights.}
Let $K$ be a totally complex number field and $r_2=[K:\Q]/2$ the number of its infinite places $v$; for each $v$
fix an embedding $\sigma_v\colon K\to\C$ inducing it. A finite place $w$ is a non-zero prime ideal of $\OO_K$; we
write $v_w$ for its normalised valuation on $K^\times$, $\Nm w$ for its absolute norm and
\[
\ell_w=\frac{\log\Nm w}{[K:\Q]}>0 .
\]
For a function $g$ on the infinite places put $\fitmi{g}=\frac1{r_2}\sum_vg(v)$, and for $\kappa_v>0$ put
$\fitmlog(\kappa)=\fitmi{\log\kappa}$. For $z\in K^\times$, a function $g$ from finite places to $\Z$ and $B\in\R$ we
say that $z$ has \emph{$g$-height at most $B$}, written $\fitht_g(z)\le B$, if $v_w(z)+g(w)\ne0$ for only finitely
many $w$ and
\[
\fitht_g(z)=\sum_w\ell_w\,\bigl|v_w(z)+g(w)\bigr|\le B .
\]
For a finite set $S$ of primes, a function $c$ on $\N$ and a finite place $w$ put
$\fitwt_S(c)(w)=\sum_{l\in S}c(l)\,v_w(l)$. For a finite set $T$ of finite places and $a\colon T\to\Z$ let
$\fitden_{T,a}=\prod_{w\in T}w^{-a(w)}$, a fractional ideal with $\log\Nm\fitden_{T,a}=-\sum_{w\in T}a(w)\log\Nm w$.
For a fractional ideal $\fitden$ and an $\OO_K$-submodule $M$ of a $K$-vector space, $\fitden M$ is the
$\OO_K$-module generated by the products $ax$, $a\in\fitden$, $x\in M$. For an $\OO_K$-submodule $A$ and a finite place
$w$, $A_{(w)}=\{x:\ sx\in A\text{ for some }s\in\OO_K\smallsetminus w\}$ is the localisation of $A$ at $w$ inside the
ambient space.

\paragraph{Slopes.}
We use the slopes of \cref{def:covolume:slopes}: a $K$-vector space $V$ with a basis and a \emph{frame}
$X=(X(v))_v$ of invertible complex matrices carries the norms $\|x\|_{X,v}=|X(v)^{-1}\sigma_v(x)|$ (coordinates in the
basis), and a non-zero finitely generated $\OO_K$-submodule $L\subseteq V$ has the slope $\mu(L)=\mu_X(L)$, the maximal
slope $\fitmumax(L)$ (supremum of $\mu(N)$ over the non-zero submodules $N\subseteq L$) and the minimal slope
$\fitmumin(L)$. The \emph{height} of a $K$-linear map $\varphi$ between two such spaces is
$h(\varphi)=\fitmi{\log\|\varphi\|_v^2}$, where $\|\varphi\|_v$ is the operator norm for the norms at $v$. For positive
real numbers $\kappa=(\kappa_v)_v$ and a frame $X$ we write $X_\kappa$ for the frame $X_\kappa(v)=\kappa_v^{1/2}X(v)$,
whose norms are $\kappa_v^{-1/2}\|\cdot\|_{X,v}$. The vector spaces below are $K^d$ with their standard bases.

\paragraph{Indefinite data and $\fitFUN$.}
We use the indefinite data of \cref{sec:families:nonempty} in the formal version \cite{LeanSource2026}. A \emph{base}
consists of a subfield $K_0\subseteq K$, an infinite place $v_0$, a set of \emph{types} with a base type $\mu_0$, and
finite sets $\mathrm{Ix}$ and $J$. A \emph{definite side} over a base provides, for every type $t$: an integer
$n_V(t)$, a hermitian matrix $H_V(t)$, an $\OO_K$-submodule $M_t\subseteq K^{n_V(t)}$, an integer $d_Q(t)$, a matrix
$P_t\in M_{d_Q(t)\times n_V(t)}(K)$, a frame $X_Q(t)$ on $K^{d_Q(t)}$ and subspaces $Q_p(t)\subseteq K^{d_Q(t)}$
($p\in\mathrm{Ix}$), together with data indexed by $J$ which play no role here. Put
\[
\Lambda_t=P_t(M_t)\subseteq K^{d_Q(t)},\qquad U_t=\ker P_t .
\]
The side is \emph{well formed} if, among other conditions, $H_V(t)$ is positive definite at every infinite place,
$M_t$ and $\Lambda_t$ are lattices, $P_t$ is surjective, $X_Q(t)$ is the quotient metric of $H_V(t)$ along $P_t$
(that is, $(X_Q(t)(v)^{-1})^*X_Q(t)(v)^{-1}=(P_tH_V(t)^{-1}P_t^*)^{-1}$ at $v$, the matrices being mapped by
$\sigma_v$), and $K^{d_Q(\mu_0)}=\bigoplus_pQ_p(\mu_0)$.

An \emph{indefinite datum} $S$ over a definite side provides: a map $\mu$ from the finite places of $K$ to types, a
finite set $\mathrm{Tw}$ of \emph{twisted} places, an integer $D\ge0$, a frame $X_E$ on $K^D$, subspaces
$E_p\subseteq K^D$, matrices $\iota_t\in M_{D\times d_Q(t)}(K)$, positive numbers $\kappa_v$, numbers
$\lambda_{j,t,p,v}$, and further data (lattices $M_j$ and theta-type maps, error terms) used in the other point
inequalities. The \emph{glued lattice} of $S$ is
\begin{equation}\label{eq:pf:fit:RR}
\fitRR=\fitRR(S)=\bigcap_w\bigl(\iota_{\mu(w)}(\Lambda_{\mu(w)})\bigr)_{(w)}\subseteq K^D,
\end{equation}
the intersection over all finite places $w$. For numbers $N$ (\cref{sec:families:nonempty}) with entry $hh$, $S$
satisfies $\fitFUN(N)$ if $\fitmumax(\fitRR)\le hh-\fitmlog(\kappa)$, the slope being taken for the frame $X_E$. For the
numbers $N_p(C,C_H)$ of a point $p$ one has $hh=x^{C_H}b$, where $x$ is the scale and $b=e^{x^\eps}$.

We shall use the following \emph{row conditions} on $S$, which are part of the conditions satisfied at the points of the
families: $D>0$; $\mu(\bar w)=\mu(w)$ for all $w$, and $w\in\mathrm{Tw}$ if and only if $\mu(w)\ne\mu_0$;
$\iota_{\mu_0}(Q_p(\mu_0))=E_p$ for all $p$; $\det X_E(v)\ne0$ for all $v$; and the \emph{scalar conditions}: every
$\iota_t$ is bijective, $\kappa_v>0$, $\lambda_{j,t,p,v}>0$, and for $t=\mu_0$ or $v=v_0$
\begin{equation}\label{eq:pf:fit:scalar}
\langle\iota_tx,\iota_ty\rangle_{X_E,v}=\kappa_v\prod_{j\in J}\lambda_{j,t,p,v}\cdot\langle x,y\rangle_{X_Q(t),v}
\qquad(x\in\iota_t^{-1}(E_p),\ y\in K^{d_Q(t)}),
\end{equation}
where $\langle x,y\rangle_{X,v}=(X(v)^{-1}\sigma_v x)^*(X(v)^{-1}\sigma_vy)$ is the hermitian form of the frame. Under
these conditions $\fitRR$ is a lattice and $\fitRR_{(w)}=(\iota_{\mu(w)}(\Lambda_{\mu(w)}))_{(w)}$ for every $w$
(\cref{pf:fit:glue}).

\paragraph{Grids.}
At a point $p$ of a family over $(x,i)$ one has a real datum $R$ (with a level character) and, for every $C_F$, the
\emph{grid data} $G$ at level $(1,C_F,b)$ for it \cite{LeanSource2026}. A grid datum consists of a finite set $S_0$ of
places of the totally real field $F$ of the host, with a base place $v_0$ above the first modulus $q_0$ and chosen
split places $w(v)$ of $E$ above them, a family of Hecke characters $\chi_H(i)$ ($1\le i\le m$) of $E$ whose infinity
types are prescribed by the gap constant, and hermitian spaces with lattices at the nodes of the grid. A \emph{Weil
datum} $W$ over $G$ consists of a standard additive character $\psi$ of $\A_F$, theta families, Cartan involutions and
orthonormal bases at the nodes, Weil representations of the dual pairs of the grid, vacuum vectors and coset
vectors; $W$ is \emph{over} $R$ if it is compatible with the real datum. The grid \emph{fits} $p$ ($\sfname{GridFit}$)
if: the values of the characters $\chi_H(i)$ and suitable roots of unity of order $\Nm v^{\mathrm{depth}(v)}$
($v\in S_0$) lie in the coefficient field; the twisted places of the point lie over the chosen pairs of places;
reduction conditions of Bloch--Kato type and a power condition hold at the twisted places; the twist characters are
compatible; and the $\lambda$-tie of the grid is full.

For a property $\mathrm{Fld}(p,C_F,G,W)$ let $\sfname{HasGrid}(\mathrm{Fld},p,C_F)$ be the statement that there are a grid
datum $G$ at level $(1,C_F,b)$ and a Weil datum $W$ over $G$, over $R$, such that $G$ fits $p$ and
$\mathrm{Fld}(p,C_F,G,W)$. The property used is
\[
\sfname{Fld}(p,C_F,G,W):\iff\text{there are }n\in\N\text{ and a choice }a\text{ of the remaining data with }\mathrm{OK}(a),
\]
where a \emph{choice of the remaining data} consists of a volume datum, real data $\rho,\xi$, a chain vector, a
real match, a sum choice, class representatives, an enumeration $e\colon\{1,\dots,n\}\to(\text{test indices})\times
(\text{weight indices})$ and constants $c(i,j,t)\in\C$, and the condition $\mathrm{OK}(a)$ asserts that the associated choice
$\mathrm{ch}(a)$ is admissible, the Eisenstein tie $\sfname{ESerTie}$ of the datum holds, the field condition
$\sfname{FieldOK}$ holds, and the square values at the base type lie in the coefficient field.

\paragraph{Families and the statement.}
For a family $\mathcal F$ and a property $Q$ of points,
\[
\fitEv(\mathcal F,Q):\iff\exists x_0\ \forall x>x_1\text{ with }x\ge x_0\ \forall i\in I_x\
\forall p\in\mathrm{Pt}(\mathcal F,x,i):\ \fitPC(p)\Rightarrow Q(x,i,p),
\]
where $\fitPC(p)$ are the point conditions of \cref{sec:families:nonempty}. At a point we write $S_{\mathrm f}$ and
$S_{\mathrm t}$ for the two indefinite data (the first and the second side), $Z_p$ for its scale (with entries $x$ and
$b$) and $N_p(C,C_H)$ for its numbers.

\begin{definition}\label{pf:fit:defCFF}
For a property $\mathrm{Fld}$ as above, $\sfname{ConstantFunctionFits}(\mathrm{Fld})$ is the statement: there is
$C_{F0}$ such that for every $C_F\ge C_{F0}$
\begin{multline*}
\fitClosed\Bigl(\sfname{FamHyp},\ (\mathcal F,C_H)\mapsto\Bigl(\sfname{DIVinN}(\mathcal F)\Rightarrow\fitEv\bigl(\mathcal F,\\
p\mapsto\sfname{HasGrid}(\mathrm{Fld},p,C_F)\Rightarrow\forall C\in\R:\ S_{\mathrm f}\text{ satisfies }
\fitFUN(N_p(C,C_H))\bigr)\Bigr)\Bigr).
\end{multline*}
\end{definition}

The hypotheses of the main theorem are statements of \cref{sec:families}: $F=\sfname{ThetaLocalInputs}$, the bundle of
$53$ classical results of \cref{sec:families:theta} (among them Howe duality at the real place, the Siegel sets of
unitary groups, the Rallis inner product formula in the convergent range and the doubling method); $\sfname{JoinOut}$,
the conjunction of five statements on CM fields: a bound for the denominators of toric values in terms of class zeta
values ($\sfname{ToricDenominator}$), the transport of CM lattices under $\Aut(\C)$ in the form of the main theorem of
complex multiplication ($\sfname{Transport}$), the algebraicity of twisted CM sums normalised by a period
($\sfname{TwistedTransport}$), bounds for the denominators of the partial zeta values of the totally real subfield
($\sfname{ZetaDenominators}$), and a reflection formula, up to units, for the Hecke $L$-functions attached to a grid
($\sfname{Refl}$); $\sfname{HeckeTateFunctionalEquation}$, the functional equation of completed Hecke $L$-functions of CM
fields with non-trivial infinity type \cite{Hecke1920,Tate1967thesis}; $\sfname{ThetaIntegralsAgree}$ and
$\sfname{LocalMatching}$ (\cref{def:families:K}); $\sfname{FunctionalsSeparate}$ (\cref{prop:families:fs}); and
$\sfname{LambdaLowerBound}_2$ (\cref{sec:families:nonempty}: the second side satisfies $\sfname{L6}$).

\begin{theorem}[\cref{prop:families:points}, second part]\label{pf:fit:main}
$\sfname{ThetaLocalInputs}$, $\sfname{JoinOut}$, $\sfname{HeckeTateFunctionalEquation}$, $\sfname{ThetaIntegralsAgree}$,
$\sfname{LocalMatching}$, $\sfname{FunctionalsSeparate}$ and $\sfname{LambdaLowerBound}_2$ imply
$\sfname{ConstantFunctionFits}(\sfname{Fld})$.
\formallabel{GlobalTheory.BigTop.Fun.Close.fourier_coefficient_fit_of_join_alt}
\end{theorem}

The proof is given in \cref{sec:pf:fit:main}.

% ---------------------------------------------------------------------------
\subsection{A slope inequality with denominators}\label{sec:pf:fit:slopes}

We begin with three facts of commutative algebra and the effect of scalings on slopes.

\begin{lemma}[Gluing of lattices]\label{pf:fit:glue}
Let $L_0$ and $L_w$ (for every finite place $w$) be lattices in $K^D$ such that $(L_w)_{(w)}=(L_0)_{(w)}$ for all but
finitely many $w$. Then $L'=\bigcap_w(L_w)_{(w)}$ is a lattice and $L'_{(w)}=(L_w)_{(w)}$ for every $w$.
\formallabel{Formalism.DivisionValues.lattice_gluing}
\end{lemma}

\begin{proof}
Let $S$ be the finite set of $w$ with $(L_w)_{(w)}\ne(L_0)_{(w)}$ and choose $0\ne c\in\OO_K$ with $cL_0\subseteq L_w$
and $cL_w\subseteq L_0$ for $w\in S$. Then $cL_0\subseteq(L_w)_{(w)}$ for all $w$, so $cL_0\subseteq L'$, and
$cL'\subseteq\bigcap_w(L_0)_{(w)}=L_0$; the last equality holds because for $x$ in the intersection the ideal
$\{s\in\OO_K:sx\in L_0\}$ lies in no maximal ideal. Hence $cL_0\subseteq L'\subseteq c^{-1}L_0$ and $L'$ is a lattice.
Clearly $L'_{(w)}\subseteq(L_w)_{(w)}$. Conversely let $y\in(L_w)_{(w)}$ and, for every $v$, let
$I_v=\{s\in\OO_K:sy\in(L_v)_{(v)}\}$. Choose $0\ne d\in\OO_K$ with $dy\in cL_0$; then $d\in I_v$ for all $v$, and
$I_v=\OO_K$ for $v=w$ and for all $v\notin S$ not dividing $d$. Each $I_v$ is non-zero and saturated with respect to
$\OO_K\smallsetminus v$, hence a power of $v$. Thus $I=\bigcap_vI_v=\{s:sy\in L'\}$ is a product of powers of primes
different from $w$, and any $s\in I\smallsetminus w$ shows $y\in L'_{(w)}$.
\end{proof}

\begin{lemma}[Local--global inclusion]\label{pf:fit:locglob}
Let $M,B$ be $\OO_K$-submodules of a $K$-vector space, $T$ a finite set of finite places and $a\colon T\to\Z$. If
$M_{(w)}\subseteq(w^{-a(w)}B)_{(w)}$ for $w\in T$ and $M_{(w)}\subseteq B_{(w)}$ for $w\notin T$, then
$M\subseteq\fitden_{T,a}B$.
\formallabel{PI.local_global_inclusion'}
\formallabel{PI.LocGlobInclWork.main}
\end{lemma}

\begin{proof}
An element $x$ belongs to an $\OO_K$-module $N$ as soon as for every $w$ there is $s\in\OO_K\smallsetminus w$ with
$sx\in N$, since then $\{s:sx\in N\}$ lies in no maximal ideal. Let $x\in M$ and fix $w$. The fractional ideal
$\prod_{v\in T\smallsetminus\{w\}}v^{-a(v)}$ contains the integral ideal $\prod_{v\in T,\,a(v)<0,\,v\ne w}v^{-a(v)}$, which
is not contained in $w$; choose $t\in\OO_K\smallsetminus w$ in it. If $w\in T$, choose $s\notin w$ with
$sx\in w^{-a(w)}B$; then $tsx\in\fitden_{T,a}B$. If $w\notin T$, choose $s\notin w$ with $sx\in B$ and $t\notin w$ in
$\fitden_{T,a}\cap\OO_K$; again $tsx\in\fitden_{T,a}B$. Hence $x\in\fitden_{T,a}B$.
\end{proof}

\begin{lemma}[Scalings]\label{pf:fit:scale}
Let $V$ be a $K$-vector space with a basis and a frame $X$.
\begin{enumerate}[label=(\alph*)]
\item For a non-zero fractional ideal $\fitden$ and a non-zero finitely generated $\OO_K$-submodule $M\subseteq V$,
$\fitden M$ is finitely generated and $\mu(\fitden M)=\mu(M)-2\log\Nm\fitden/[K:\Q]$; the same holds for $\fitmumax$
and $\fitmumin$.
\item For $\kappa_v>0$ and a lattice $R\subseteq K^D$ with a frame $X$ on $K^D$,
$\fitmumax(R;X)=\fitmumax(R;X_\kappa)-\fitmlog(\kappa)$.
\end{enumerate}
\formallabel{Formalism.LFunctionChainParameters.LLcv.ideal_scaling_holds}
\formallabel{Formalism.DivisionValues.kappa_normalise}
\end{lemma}

\begin{proof}
Let $\rho_X\colon V\to\bigoplus_v\C^I$ be the map of \cref{def:covolume:slopes}, so that
$\mu(N)=-2\log\covol\rho_X(N)/([K:\Q]\rank N)+c_K$ with $c_K$ depending only on $K$. (a) If $N$ is torsion free of rank
$r$, then $N\cong\OO_K^{r-1}\oplus\mathfrak b$ and $[N:\fitden N]=\Nm\fitden^{\,r}$ for integral $\fitden$; by
multiplicativity $\covol\rho_X(\fitden N)=\Nm\fitden^{\,r}\covol\rho_X(N)$ for all $\fitden$, which gives the formula for
$\mu$. Since $N\mapsto\fitden N$ is a bijection from the submodules of $M$ to those of $\fitden M$ (with inverse
$N'\mapsto\fitden^{-1}N'$) preserving saturation and quotients up to the same factor, the formulas for $\fitmumax$ and
$\fitmumin$ follow. (b) The norms of $X_\kappa$ at $v$ are $\kappa_v^{-1/2}$ times those of $X$. For a submodule $N$ of
rank $r$ the real span of $\rho_X(N)$ is the sum over $v$ of complex subspaces of dimension $r$ of the $v$-components,
so $\covol\rho_{X_\kappa}(N)=\prod_v\kappa_v^{-r}\covol\rho_X(N)$ and
$\mu_{X_\kappa}(N)=\mu_X(N)+\frac2{[K:\Q]}\sum_v\log\kappa_v=\mu_X(N)+\fitmlog(\kappa)$. Taking the supremum over $N$
gives (b).
\end{proof}

\begin{remark}\label{pf:fit:rem:kappa}
The formal statement of \cref{pf:fit:scale}(b) also contains the corresponding identities for the slope of the dual
lattice for the dual frame, for heights of maps into the dual space and for ratios of determinants; they are not used
here. Its formal proof passes through the dual lattice and the identity
$\fitmumin(\text{dual})=-\fitmumax$; the direct argument above computes the same quantities.
\end{remark}

The following inequality is Bost's slope inequality \cite{Bost1996,Bost2001} in the form proved in the development:
by \lref{Formalism.HN.slope_inequality_mean}, if $\varphi\colon V\to W$ is $K$-linear between spaces with frames,
$A\subseteq V$ and $B\subseteq W$ are finitely generated, $\varphi(A)\ne0$, $\varphi(A)\subseteq B$ and $\varphi$ is
injective on $A$, then
\begin{equation}\label{eq:pf:fit:bost}
\fitmumax(A)\le\fitmumax(B)+h(\varphi).
\end{equation}
(For a non-zero submodule $N\subseteq A$, $\varphi$ multiplies covolumes at $v$ by at most $\|\varphi\|_v^{2\rank N}$, so
$\mu(N)\le\mu(\varphi(N))+h(\varphi)\le\fitmumax(B)+h(\varphi)$; compare \cref{lem:covolume:slopes}.)

\begin{proposition}[Slope inequality with denominators]\label{pf:fit:slopeden}
Let $\varphi\colon V\to W$ be an injective $K$-linear map between spaces with frames, $A\subseteq V$ and $B\subseteq W$
non-zero finitely generated $\OO_K$-submodules, $T$ a finite set of finite places and $a\colon T\to\Z$ with
$\varphi(A)\subseteq\fitden_{T,a}B$. Then
\[
\fitmumax(A)\le\fitmumax(B)+h(\varphi)+2\sum_{w\in T}\ell_wa(w).
\]
\formallabel{Formalism.AtG.grid_lemma_A}
\end{proposition}

\begin{proof}
$\fitden_{T,a}B$ is finitely generated and $\varphi(A)\ne0$ because $\varphi$ is injective and $A\ne0$. By
\eqref{eq:pf:fit:bost} and \cref{pf:fit:scale}(a),
$\fitmumax(A)\le\fitmumax(\fitden_{T,a}B)+h(\varphi)=\fitmumax(B)-2\log\Nm\fitden_{T,a}/[K:\Q]+h(\varphi)$, and
$-\log\Nm\fitden_{T,a}/[K:\Q]=\sum_{w\in T}a(w)\log\Nm w/[K:\Q]=\sum_{w\in T}\ell_wa(w)$.
\end{proof}

\begin{proposition}[Largest slope of a lattice with denominators]\label{pf:fit:largest}
Let $R\subseteq K^D$ ($D>0$) be a lattice with a frame $X_E$, let $\kappa_v>0$, let $M\subseteq W$ be a non-zero
finitely generated submodule of a space with a frame $X'$, and let $\Psi\colon K^D\to W$ be injective with
$\Psi(R)\subseteq\fitden_{T,a}M$. Let $r\colon T\to K^\times$ and $d_F\colon T\to\R$ with $a(w)\le v_w(r(w))+d_F(w)$
for $w\in T$, and let $E_M,H,G,B_F$ be real numbers with
\[
\fitmumax(M)\le E_M,\quad h(\Psi)\le H,\quad\sum_{w\in T}\ell_wv_w(r(w))\le G,\quad\sum_{w\in T}\ell_wd_F(w)\le B_F,
\]
where $h(\Psi)$ is taken for the frames $X_{E,\kappa}$ and $X'$. Then
$\fitmumax(R;X_E)\le-\fitmlog(\kappa)+E_M+H+2(G+B_F)$.
\formallabel{Formalism.AtG.grid_function_of_lemma_A}
\formallabel{Formalism.AtG.grid_function_bound}
\end{proposition}

\begin{proof}
$R$ is finitely generated and non-zero since $D>0$. By \cref{pf:fit:scale}(b) and \cref{pf:fit:slopeden} for the frame
$X_{E,\kappa}$,
$\fitmumax(R;X_E)=\fitmumax(R;X_{E,\kappa})-\fitmlog(\kappa)\le-\fitmlog(\kappa)+E_M+H+2\sum_{w\in T}\ell_wa(w)$, and
$\sum_{w\in T}\ell_wa(w)\le\sum_{w\in T}\ell_w(v_w(r(w))+d_F(w))\le G+B_F$ because $\ell_w\ge0$.
\end{proof}

% ---------------------------------------------------------------------------
\subsection{The grid estimate}\label{sec:pf:fit:grid}

In the application the elements $r(w)$ of \cref{pf:fit:largest} are products $\prod_{i,j}r_{\mu(w)}(i,j)$ of
\emph{square ratios}, and their valuations are bounded place by place by an expression built from a few algebraic
numbers of controlled height. This section sums such local bounds over the finite places.

\begin{lemma}[Product formula]\label{pf:fit:prodformula}
\begin{enumerate}[label=(\alph*)]
\item For $z\in K^\times$, $\sum_w\ell_wv_w(z)=\fitmi{\log|\sigma_vz|}$.
\item For an integer $n\ge1$ and every finite set $F$ of finite places, $\sum_{w\in F}\ell_wv_w(n)\le\log n$.
\item Let $z\in K^\times$, $d_l\ge0$ a function on the finite places and $B_\delta\in\R$ such that
$v_w(z)\ge-3d_l(w)$ for all $w$ and $\sum_{w\in F}\ell_wd_l(w)\le B_\delta$ for every finite set $F$. Then
$\fitht_0(z)\le\fitmi{\log^+|\sigma_vz|}+6B_\delta$.
\end{enumerate}
\formallabel{Formalism.PairingDivisionLocalBalance.sum_valuation_nat}
\formallabel{Formalism.AtG.partial_nat}
\formallabel{Formalism.AtG.norm_height_bound}
\formallabel{Formalism.ValuationGridEstimates.height_bound_alt}
\end{lemma}

\begin{proof}
(a) is the product formula (\lref{Formalism.DivisionValues.prod_formula'}: the sum is finite and equals
$\log|\Nm_{K/\Q}z|/[K:\Q]$): factor the principal ideal $(z)$ and use $|\Nm_{K/\Q}z|=\prod_v|\sigma_vz|^2$, $K$ being
totally complex. (b) All terms $\ell_wv_w(n)$ are non-negative and their total sum is $\log n$ by (a). (c) Writing
$|t|=t+2\max(0,-t)$ and using (a),
$\sum_w\ell_w|v_w(z)|=\fitmi{\log|\sigma_vz|}+2\sum_w\ell_w\max(0,-v_w(z))\le\fitmi{\log^+|\sigma_vz|}
+2\sum_{w}3\ell_wd_l(w)$, the last sum running over the finite support of $v_w(z)$.
\end{proof}

\begin{lemma}\label{pf:fit:primes}
For every finite set $P$ of primes, $\sum_{l\in P}\log l/l^2<1$. Consequently, if $N_0\colon\N\to\R$ satisfies
$\sum_{l\in P}N_0(l)\log l\le B$ for all finite sets of primes $P$, then
$\sum_{l\in P}(N_0(l)+l^{-2})\log l\le B+1$ for all such $P$.
\formallabel{Formalism.ValuationGridEstimates.prime_sum_bound}
\formallabel{Formalism.ValuationGridEstimates.sum_log_div_sq_telescoping}
\formallabel{Formalism.ValuationGridEstimates.prime_budget_bound}
\end{lemma}

\begin{proof}
Let $f(t)=(1+\log t)/t$, so $f'(t)=-\log t/t^2$; since $\log t/t^2$ decreases for $t\ge4$, we have
$\log n/n^2\le f(n-1)-f(n)$ for $n\ge5$, and summing, $\sum_{n=5}^M\log n/n^2\le f(4)-f(M)<(1+\log4)/4$ for $M\ge4$.
Taking $M\ge\max P$, the set $P$ is contained in $\{2,3\}\cup[5,M]$, and
$\sum_{l\in P}\log l/l^2<\frac{\log2}4+\frac{\log3}9+\frac{1+2\log2}4<0.174+0.123+0.597<1$, using
$\log3\le2\log2$ and $\log2<0.6932$. The second statement follows by adding.
\end{proof}

We now fix the data of the grid estimate. Let $m,N\in\N$ and let $\fitSq=\{1,\dots,m\}^2$ be the set of index pairs
$\eta=(i,j)$. Let:
\begin{itemize}
\item $d,h,y\colon\fitSq\to\N$ with $d(\eta)=2h(\eta)$, $d(\eta)\le2N$ and $y(\eta)\le m$;
\item $r\colon\{\text{finite places}\}\times\fitSq\to K^\times$ (the ratios), $z_A,z'_A\colon\fitSq\to K^\times$,
$z_T\colon\fitSq\times\Z\to K^\times$, $z_{Td}\colon\fitSq\times\{0,1\}\times\Z\to K^\times$, integers
$g_T(\eta,s,w)$, and a function $c$ from the finite places to $\{0,1\}$;
\item a function $d_l\ge0$ on the finite places, finite sets $S_d,S_e$ of primes and numbers $N_0(\eta,s,l)$ which
are $\ge0$ for $1\le s\le h(\eta)$ and $l\in S_d$;
\item a finite set $S_0$ with functions $c_0,e,q\colon S_0\to\N$ such that $c_0\le e$ on $S_0$.
\end{itemize}
For a finite place $w$ and $\eta\in\fitSq$ put
\begin{multline}\label{eq:pf:fit:sqbound}
\beta_w(\eta)=|v_w(z_A(\eta))|+|v_w(z'_A(\eta))|\\
+4\sum_{s=1}^{h(\eta)}\Bigl(\bigl|v_w(z_T(\eta,s))+g_T(\eta,s,w)\bigr|
+\bigl|v_w(z_{Td}(\eta,c(w),s))+g_T(\eta,s,w)\bigr|\Bigr)\\
+5\sum_{s=1}^{h(\eta)}\fitwt_{S_d}(N_0(\eta,s,\cdot))(w)+d(\eta)\fitwt_{S_e}(l\mapsto l^{-2})(w)\\
+6d_l(w)+2y(\eta)\sum_{q\in S_0}c_0(q)v_w(q(q)),
\end{multline}
and for real $B,B_N,B_\delta$ define the \emph{grid budget}
\begin{equation}\label{eq:pf:fit:budget}
\fitbudget(B,B_N,B_\delta)=m^2\Bigl(2(2+8N)B+10NB_N+2N+6B_\delta+4m\sum_{q\in S_0}e(q)\log q(q)\Bigr).
\end{equation}

\begin{proposition}[Grid estimate]\label{pf:fit:gridheight}
Assume that for real numbers $B$, $B_\delta$ and $B_N\ge0$:
\begin{enumerate}[label=(\roman*)]
\item $\fitht_0(z_A(\eta))\le B$ and $\fitht_0(z'_A(\eta))\le B$ for all $\eta$; $\fitht_{g_T(\eta,s,\cdot)}(z_T(\eta,s))\le B$
and $\fitht_{g_T(\eta,s,\cdot)}(z_{Td}(\eta,c,s))\le B$ for all $\eta,s,c$;
\item $\sum_{l\in S_d}N_0(\eta,s,l)\log l\le B_N$ for $1\le s\le h(\eta)$;
\item $\sum_{w\in F}\ell_wd_l(w)\le B_\delta$ for every finite set $F$ of finite places.
\end{enumerate}
Then for every finite set $W_0$ of finite places such that $v_w(r(w,\eta))\le\beta_w(\eta)$ for $w\in W_0$ and all
$\eta$,
\[
\sum_{w\in W_0}\ell_w\,v_w\Bigl(\prod_{\eta\in\fitSq}r(w,\eta)\Bigr)\le\fitbudget(B,B_N,B_\delta).
\]
\formallabel{Formalism.AtG.gridHeightBoundIndexed}
\formallabel{Formalism.AtG.grid_squares_index_estimate}
\formallabel{Formalism.ValuationGridEstimates.grid_index_bound}
\formallabel{Formalism.AtG.j3Wt}
\formallabel{Formalism.AtG.grid_ell_bound}
\end{proposition}

\begin{proof}
\emph{Normalisation.} Every finite place $w$ lies over a unique prime number $\ell(w)$, and $e_w=v_w(\ell(w))\ge1$. For a
prime $l\ne\ell(w)$ we have $v_w(l)=0$ (as $l$ and $\ell(w)$ are coprime), hence
$\fitwt_S(c)(w)=c(\ell(w))e_w$ if $\ell(w)\in S$ and $0$ otherwise. Put $\lambda_w=\ell_we_w$, and divide all
valuations at $w$ by $e_w$. Replacing $N_0(\eta,s,l)$ by $N_1(\eta,s,l)=N_0(\eta,s,l)$ for $1\le s\le h(\eta)$ and
$l\in S_d$ and by $0$ otherwise, the hypothesis $v_w(r(w,\eta))\le\beta_w(\eta)$ becomes
\begin{multline*}
\frac{v_w(r(w,\eta))}{e_w}\le\sum_{x\in X(\eta,c(w))}\frac{|u_x(w)|}{e_w}+5\sum_{s=1}^{h(\eta)}N_1(\eta,s,\ell(w))\\
+\frac{d(\eta)}{\ell(w)^2}+6\frac{d_l(w)}{e_w}+2y(\eta)\sum_{q\in S_0}c_0(q)\frac{v_w(q(q))}{e_w},
\end{multline*}
where $X(\eta,c)$ is a set of $2+8h(\eta)\le2+8N$ ``terms'' $x$ (the two numbers $z_A(\eta),z'_A(\eta)$ and four
copies of each $z_T(\eta,s)$, $z_{Td}(\eta,c,s)$, $1\le s\le h(\eta)$) and $u_x(w)$ is $v_w$ of the term plus the
shift $g_T$ where present.

\emph{Fibres.} For a prime $l$ and a finite set $F$ of places over $l$, \cref{pf:fit:prodformula}(b) gives
$\sum_{w\in F}\lambda_w=\sum_{w\in F}\ell_wv_w(l)\le\log l$. Hence for every $\varphi\colon\N\to\R_{\ge0}$,
$\sum_{w\in W_0}\lambda_w\varphi(\ell(w))\le\sum_{l\in\ell(W_0)}\varphi(l)\log l$.

\emph{Summation.} Multiply the normalised inequality by $\lambda_w\ge0$ and sum over $w\in W_0$ and $\eta$. Since
$\lambda_w|u_x(w)|/e_w=\ell_w|u_x(w)|$, splitting $W_0$ according to the value of $c(w)$ and using (i), the terms give
at most $\sum_{c}\sum_{x\in X(\eta,c)}B\le2(2+8N)B$ for each $\eta$. By the fibre bound and (ii), the second terms
give at most $5h(\eta)B_N\le5NB_N$. By the fibre bound and \cref{pf:fit:primes}, the third give at most
$d(\eta)\sum_{l}\log l/l^2\le2N$. By (iii) the fourth give at most $6B_\delta$. By \cref{pf:fit:prodformula}(b),
$\sum_{w\in W_0}\ell_wv_w(q(q))\le\log q(q)$, so the last give at most $2m\sum_qe(q)\log q(q)$ (as $c_0\le e$ and
$y\le m$). Summing over the $m^2$ pairs $\eta$ and using $v_w(\prod_\eta r(w,\eta))=\sum_\eta v_w(r(w,\eta))$ gives
the bound $m^2\bigl(2(2+8N)B+5NB_N+2N+6B_\delta+2m\sum e\log q\bigr)\le\fitbudget(B,B_N,B_\delta)$.
\end{proof}

\begin{theorem}[Slope bound from grid data]\label{pf:fit:passage}
Let $R,X_E,\kappa,M,X',\Psi,T,a$ be as in \cref{pf:fit:largest}, let $d_F\colon T\to\R$, and let $E_M,H,B_F,P$ be real
numbers with $\fitmumax(M)\le E_M$, $h(\Psi)\le H$ (frames $X_{E,\kappa}$, $X'$) and $\sum_{w\in T}\ell_wd_F(w)\le B_F$.
Let the data and bounds $B,B_N,B_\delta$ of \cref{pf:fit:gridheight} be given, with $v_w(r(w,\eta))\le\beta_w(\eta)$ and
$a(w)\le v_w(\prod_\eta r(w,\eta))+d_F(w)$ for $w\in T$. If
$E_M+H+2(\fitbudget(B,B_N,B_\delta)+B_F)\le P$, then $\fitmumax(R;X_E)\le-\fitmlog(\kappa)+P$.
\formallabel{Formalism.AtG.height_h0_index_bound}
\end{theorem}

\begin{proof}
Apply \cref{pf:fit:gridheight} with $W_0=T$ and then \cref{pf:fit:largest} with $r(w)=\prod_\eta r(w,\eta)$ and
$G=\fitbudget(B,B_N,B_\delta)$.
\end{proof}

% ---------------------------------------------------------------------------
\subsection{Functional data and the conditional slope bound}\label{sec:pf:fit:cond}

In this section $K$ is a totally complex number field, $S$ an indefinite datum over a well-formed definite side, and we
write $\iota=\iota_{\mu_0}$, $P=P_{\mu_0}$, $X_Q=X_Q(\mu_0)$ and $Q_p=Q_p(\mu_0)$ when no confusion is possible.

\begin{definition}[Functional data]\label{pf:fit:fundata}
A \emph{functional datum} $O$ for the definite side consists of:
\begin{itemize}
\item an integer $n$, a frame $X'$ on $K^n$, an $\OO_K$-submodule $M\subseteq K^n$ and, for every type $t$, a matrix
$V_t\in M_{n\times n_V(t)}(K)$;
\item a finite set $T_f$ of finite places and integers $d_F(w)$;
\item an integer $m$, \emph{square ratios} $r_t(\eta)\in K$ for every type $t$ and $\eta\in\fitSq=\{1,\dots,m\}^2$, and
the data $d,h,y,z_A,z'_A,z_T,z_{Td},g_T,c,d_l,S_d,S_e,N_0$ of \cref{sec:pf:fit:grid};
\item a set with a finite subset $S_0$, a distinguished element $q_*$ and functions $c_0,q$ to $\N$;
\item integers $f,k,t,g_1,g_2,p_d,p'_d\ge0$, $d_E\in\Z$, a function $\mathrm{cl}\colon\N^2\to\Z$ and a real number
$B_w$.
\end{itemize}
Its \emph{constants} are defined as follows, for real $C,x,b$ and with $L=\log b$:
\begin{gather*}
r=k+m-1,\quad n_r=k+m,\quad s=k+m+5,\quad N=k+m(t+s+1),\\
n_p=(r+1)+m(r+5),\qquad L_c=\max_{0\le\tau<f}\max\bigl(\mathrm{cl}(\tau,2)-\mathrm{cl}(\tau,k),0\bigr),\\
\gamma(C)=\max\bigl(\lceil CN^3+Cn_r^7+Cn_r^3L_c^{1/2}\rceil,\ 2N^3+2s+3,\ \lceil48n_r^2m^{1/2}\rceil\bigr),\\
R_p(C)=(m+1)(\gamma(C)+1)+L_c+n_p,\qquad E_M(C)=4fn_pR_p(C)\log\bigl(n_pR_p(C)\bigr),\\
\Pi(C,b)=Cf^2n_r^6\log\bigl(b\,n_rg_1g_2p_dp'_d|d_E|\bigr),\\
e(q)=\begin{cases}\lfloor\Pi(C,b)/\log q(q_*)\rfloor+1&q=q_*,\\2&q\ne q_*,\end{cases}\\
\Sigma_0=\sum_{q\in S_0}e(q)\log q(q),\qquad c_S=\sum_{q\in S_0}c_0(q)\log q(q),\\
B_{\mathrm{FP}}=CfbNm(mL+m\Sigma_0)+Cfm^2N^2(L+\Sigma_0)b,\\
B_{\mathrm{IPa}}=CfmNbL+Cm^2N^2\Sigma_0,\qquad B_{\mathrm{IP}}=B_{\mathrm{IPa}}+3CfmN^2L,\\
B_\Delta=2CfN\log(2fN)+2CfNB_wL+4CfN^2B_wc_S,\qquad B_{\mathrm{TH}}(x)=Cfb(L+\log x),\\
B_{\mathrm{Ar}}=2CfNB_w\bigl(L+\log(2NB_w)\bigr)+CfN^2\bigl(L+\log(2fNB_w)\bigr)+CN^2c_S,\\
\fitpol(C,b)=10^5C^3fm^2N^2\bigl(rNL+\Pi(C,b)\bigr).
\end{gather*}
\end{definition}

\begin{definition}[Leaf conditions]\label{pf:fit:leafconds}
Let $O$ be a functional datum, $S$ an indefinite datum with map $\mu$, and $C,x,C_H,b$ real numbers, $L=\log b$. The
\emph{leaf conditions} are:
\begin{enumerate}[label=(L\arabic*),leftmargin=*]
\item\label{pf:fit:L:MM} (minimal mass) $M$ is finitely generated and non-zero, the $X'(v)$ are invertible and
$\fitmumax(M;X')\le E_M(C)$;
\item\label{pf:fit:L:C} (unique map) there is $\Psi\in M_{n\times D}(K)$ with
$\Psi\iota_tP_t=\bigl(\prod_{\eta\in\fitSq}r_t(\eta)\bigr)^{-1}V_t$ for $t=\mu_0$ and for every $t=\mu(w)$;
\item\label{pf:fit:L:Ker} (kernel) $\ker V_{\mu_0}=U_{\mu_0}$;
\item\label{pf:fit:L:FP} (fixed point) for every finite place $w$,
$\bigl(V_{\mu(w)}(M_{\mu(w)})\bigr)_{(w)}\subseteq(w^{-d_F(w)}M)_{(w)}$; $d_F(w)=0$ for $w\notin T_f$; and
$2\sum_{w\in T_f}\ell_wd_F(w)\le B_{\mathrm{FP}}$;
\item\label{pf:fit:L:IP} (infinite places) every $\nu\in M_{n\times d_Q(\mu_0)}(K)$ with $\nu P=V_{\mu_0}$ has
height $h(\nu)\le B_{\mathrm{IPa}}$ for the frames $X_Q$ and $X'$; and there is a function $g$ on the infinite places
with $g(v)\le\sum_{j\in J}\log\lambda_{j,\mu_0,p,v}$ for all $v$ and all $p$ with $Q_p\ne0$, and
$\fitmi{g}\ge-3CfmN^2L$;
\item\label{pf:fit:L:Sq} (one square) $r_{\mu_0}(\eta)=1$ and $r_{\mu(w)}(\eta)\ne0$ for all $w,\eta$; the elements of
$S_e$ are primes; $d(\eta)=2h(\eta)\le2N$ and $y(\eta)+1\le m$; $c_0(q)\le e(q)$ for $q\in S_0$; and
$v_w(r_{\mu(w)}(\eta))\le\beta_w(\eta)$ for all $w,\eta$, with $\beta_w$ as in \eqref{eq:pf:fit:sqbound};
\item\label{pf:fit:L:Den} (density) the elements of $S_d$ are primes, and for $1\le s\le h(\eta)$:
$N_0(\eta,s,l)\ge0$ for $l\in S_d$ and $\sum_{l\in S_d}N_0(\eta,s,l)\log l\le B_{\mathrm{TH}}(x)$;
\item\label{pf:fit:L:Int} (integrality) $d_l\ge0$; $\sum_{w\in F}\ell_wd_l(w)\le B_\Delta$ for every finite $F$; and
$z_A(\eta),z'_A(\eta)\ne0$ with $v_w(z_A(\eta)),v_w(z'_A(\eta))\ge-3d_l(w)$ for all $w$;
\item\label{pf:fit:L:Arch} (archimedean) $\fitmi{\log^+|\sigma_vz_A(\eta)|}\le B_{\mathrm{Ar}}$ and
$\fitmi{\log^+|\sigma_vz'_A(\eta)|}\le B_{\mathrm{Ar}}$;
\item\label{pf:fit:L:Thb} (theta bound) $z_T(\eta,s)\ne0$ and $z_{Td}(\eta,c,s)\ne0$ for $1\le s\le h(\eta)$, and for all
$\eta,s,c$ the heights $\fitht_{g_T(\eta,s,\cdot)}(z_T(\eta,s))$ and $\fitht_{g_T(\eta,s,\cdot)}(z_{Td}(\eta,c,s))$ are at most
$B_{\mathrm{TH}}(x)$;
\item\label{pf:fit:L:Const} (ranks) $2\le m$, $m+3\le k$ and $m^4\le f\le k$;
\item\label{pf:fit:L:Lc} (lattice constant) $L_c\le b$;
\item\label{pf:fit:L:Size} (degrees and sizes) $mfN^2\le b$, $80r^4\le N$, $n_p\le2mr$, $R_p(C)\le CmN^3b$ and
$\log(n_pR_p(C))\le4L$;
\item\label{pf:fit:L:Mod} (prime power moduli) $q(q)$ is a prime power for $q\in S_0$, and $\Sigma_0\le2\Pi(C,b)$;
\item\label{pf:fit:L:CS} $c_S\le45fL$;
\item\label{pf:fit:L:Wt} (weights) $1\le B_w\le Cb$ and $\log B_w\le CL$;
\item\label{pf:fit:L:Pol} (polynomial) $\fitpol(C,b)\le x^{C_H}$.
\end{enumerate}
\end{definition}

\begin{lemma}[A Gram inequality]\label{pf:fit:gram}
Let $P$ be a finite set, $a_p\in\C^d$ and $A_p\in\C^{D}$ for $p\in P$, $s_p\in\R$ and $m\in\R$, such that
$A_p^*A_{p'}=s_p\,a_p^*a_{p'}$ for all $p,p'$ and $s_p\ge m$ whenever $a_p\ne0$. Then
$\bigl|\sum_pA_p\bigr|^2\ge m\bigl|\sum_pa_p\bigr|^2$.
\formallabel{Formalism.L2.Indef.Tree.gram_core}
\end{lemma}

\begin{proof}
If $a_p^*a_{p'}\ne0$, then conjugating $A_p^*A_{p'}=s_pa_p^*a_{p'}$ and comparing with $A_{p'}^*A_p=s_{p'}a_{p'}^*a_p$ gives
$s_p=s_{p'}$; moreover $a_p\ne0$, so $s_p\ge m$. Put $c_p=\max(s_p-m,0)^{1/2}$. Then
$s_pa_p^*a_{p'}=m\,a_p^*a_{p'}+c_pc_{p'}a_p^*a_{p'}$ for all $p,p'$ (both sides vanish if $a_p^*a_{p'}=0$). Summing,
$|\sum A_p|^2=\sum_{p,p'}s_pa_p^*a_{p'}=m|\sum a_p|^2+|\sum c_pa_p|^2\ge m|\sum a_p|^2$.
\end{proof}

\begin{lemma}[Height of the map at the infinite places]\label{pf:fit:ip}
Let $S$ satisfy the row conditions of \cref{sec:pf:fit:setting}, let $g$ be a function on the infinite places with
$g(v)\le\sum_j\log\lambda_{j,\mu_0,p,v}$ for all $v$ and all $p$ with $Q_p\ne0$, let $X'$ be a frame on $K^n$ and
$\Psi\in M_{n\times D}(K)$ injective. Then:
\begin{enumerate}[label=(\alph*)]
\item for every $v$ and $q\in K^{d_Q(\mu_0)}$, $e^{g(v)}\|q\|^2_{X_Q,v}\le\|\iota q\|^2_{X_{E,\kappa},v}$;
\item $h(\Psi)\le h(\Psi\eta)-\fitmi{g}$, the heights being taken for the frames $(X_{E,\kappa},X')$ and $(X_Q,X')$.
\end{enumerate}
\formallabel{Formalism.L2.Indef.Tree.ip_point}
\formallabel{Formalism.L2.Indef.Tree.ip_first}
\end{lemma}

\begin{proof}
(a) Since $K^{d_Q(\mu_0)}=\bigoplus_pQ_p$, write $q=\sum_pq_p$ with $q_p\in Q_p$. By the row conditions
$\iota(Q_p)=E_p$ and $\iota$ is injective, so $\iota^{-1}(E_p)=Q_p$, and \eqref{eq:pf:fit:scalar} gives
$\langle\iota q_p,\iota q_{p'}\rangle_{X_E,v}=\kappa_vs_p\langle q_p,q_{p'}\rangle_{X_Q,v}$ with
$s_p=\prod_j\lambda_{j,\mu_0,p,v}$. Let $a_p$ and $A_p$ be the coordinate vectors of $q_p$ for $X_Q$ at $v$ and of
$\iota q_p$ for $X_{E,\kappa}$ at $v$; since the norms of $X_{E,\kappa}$ are $\kappa_v^{-1/2}$ times those of $X_E$,
$A_p^*A_{p'}=s_pa_p^*a_{p'}$. If $a_p\ne0$ then $q_p\ne0$, so $Q_p\ne0$ and $e^{g(v)}\le
\exp\sum_j\log\lambda_{j,\mu_0,p,v}=s_p$ (all $\lambda>0$). \Cref{pf:fit:gram} with $m=e^{g(v)}$ gives (a).

(b) Let $v$ be an infinite place. Every $y\in K^D$ is $\iota q$ for some $q$, and by (a)
$\|\Psi y\|_{X',v}=\|\Psi\iota q\|_{X',v}\le\|\Psi\iota\|_v\|q\|_{X_Q,v}\le\|\Psi\iota\|_ve^{-g(v)/2}\|y\|_{X_{E,\kappa},v}$.
By \lref{Formalism.L2.Geo.ExtensionFiltration.opNormAt_le_of_points} (an inequality
$\|\varphi(y)\|_v\le c\|y\|_v$, $c\ge0$, for all $K$-rational $y$ bounds the operator norm at $v$ by $c$, the rational
vectors being dense), $\|\Psi\|_v\le\|\Psi\iota\|_ve^{-g(v)/2}$. Since $\Psi$ is injective and $D>0$, $\|\Psi\|_v>0$,
hence $\log\|\Psi\|_v^2\le\log\|\Psi\iota\|_v^2-g(v)$; averaging over $v$ gives (b).
\end{proof}

\begin{lemma}[Injectivity]\label{pf:fit:psiinj}
If $\Psi\iota P=V_{\mu_0}$, $\ker V_{\mu_0}=U_{\mu_0}$, $\iota$ is bijective and $P$ is surjective, then $\Psi$ is
injective.
\formallabel{Formalism.L2.Indef.Tree.psi_inj}
\end{lemma}

\begin{proof}
If $\Psi y=\Psi y'$, write $y=\iota Pu$ and $y'=\iota Pu'$. Then $V_{\mu_0}(u-u')=0$, so $u-u'\in U_{\mu_0}=\ker P$ and
$y=y'$.
\end{proof}

\begin{lemma}[Local inclusion]\label{pf:fit:hencelocal}
Let $w$ be a finite place, $t$ a type, $\fitRR$ an $\OO_K$-submodule of $K^D$ with
$\fitRR_{(w)}=(\iota_t(\Lambda_t))_{(w)}$, $\Psi\in M_{n\times D}(K)$, $0\ne c\in K$, $\delta\in\Z$ and $V\in
M_{n\times n_V(t)}(K)$ with $\Psi\iota_tP_t=c^{-1}V$ and $(V(M_t))_{(w)}\subseteq(w^{-\delta}M)_{(w)}$. Then
$(\Psi(\fitRR))_{(w)}\subseteq(w^{-(v_w(c)+\delta)}M)_{(w)}$.
\formallabel{Formalism.L2.Indef.Tree.hence_local}
\formallabel{Formalism.L2.Indef.Tree.unit_at_val}
\end{lemma}

\begin{proof}
First, there is $\tau\in\OO_K\smallsetminus w$ with $\tau c^{-1}\in w^{-v_w(c)}$: the fractional ideal
$J=w^{-v_w(c)}(c)$ has $w$-adic valuation $0$, and by the Chinese remainder theorem it contains an element
$\tau\in\OO_K\smallsetminus w$; then $\tau c^{-1}\in Jc^{-1}=w^{-v_w(c)}$. Now let $y$ with $sy=\Psi r$, $r\in\fitRR$,
$s\notin w$. There are $s_1\notin w$ and $z\in M_t$ with $s_1r=\iota_tP_tz$, hence $s_1sy=c^{-1}Vz$, and $s_2\notin w$
with $s_2Vz\in w^{-\delta}M$. Then $\tau s_2s_1sy=(\tau c^{-1})(s_2Vz)\in w^{-v_w(c)-\delta}M$ and $\tau s_2s_1s\notin w$.
\end{proof}

\begin{lemma}[The numerical inequality]\label{pf:fit:size}
Let $O$ be a functional datum and $C,x,C_H,b$ real numbers with $C\ge1$, $x\ge1$, $L=\log b\ge1$ and
$\log b+\log x\le3\log b$, satisfying \ref{pf:fit:L:Const}, \ref{pf:fit:L:Size}, \ref{pf:fit:L:Mod}, \ref{pf:fit:L:CS},
\ref{pf:fit:L:Wt} and \ref{pf:fit:L:Pol}. Put $\bar B=\max(B_{\mathrm{Ar}}+6B_\Delta,B_{\mathrm{TH}}(x))$ and let
$\fitbudget$ be the grid budget \eqref{eq:pf:fit:budget} for $m,N,S_0,e,q$ of $O$. Then
\[
E_M(C)+B_{\mathrm{IP}}+2\bigl(\fitbudget(\bar B,B_{\mathrm{TH}}(x),B_\Delta)+B_{\mathrm{FP}}/2\bigr)\le x^{C_H}b.
\]
\formallabel{Formalism.L2.Indef.Tree.size_le}
\formallabel{Formalism.L2.Indef.Tree.size_arith}
\end{lemma}

\begin{proof}
This is an elementary verification. From \ref{pf:fit:L:Const} and the definition of $N$ one gets $1\le f\le k\le r+1$,
$r\le N$, $m\le N$ and $mf\ge2$; with \ref{pf:fit:L:Size}, $b\ge mfN^2\ge2fN\ge2N$, so $b\ge1$; also
$\log x\le2L$, $\Sigma_0\ge0$, $c_S\ge0$, $n_p\ge1$ and $R_p\ge1$. Then, step by step:
$E_M\le32C^3fm^2rN^3bL$ (from $n_p\le2mr$, $R_p\le CmN^3b$, $\log(n_pR_p)\le4L$);
$B_\Delta\le186C^2f^2N^2bL$ (from $\log(2fN)\le L$, $B_w\le Cb$, $c_S\le45fL$);
$B_{\mathrm{Ar}}\le57C^3fN^2bL$ (from $L+\log(2NB_w)\le4CL$ and $L+\log(2fNB_w)\le4CL$, consequences of
$\log B_w\le CL$ and $\log(2fN)\le L$); $B_{\mathrm{TH}}\le3CfbL$; hence $\bar B\le1176C^3f^2N^2bL$. Inserting these
bounds, $\Sigma_0\le2\Pi$, $f\le r$ and $m\le C^3fN^2b$ into the definitions of $B_{\mathrm{IP}}$, $B_{\mathrm{FP}}$ and
$\fitbudget$, every term on the left is bounded by a multiple of $C^3fm^2rN^3bL$ or of $C^3fm^2N^2\Pi b$, and the
multiples add up to less than $10^5$. Hence the left side is at most $\fitpol(C,b)\,b\le x^{C_H}b$ by
\ref{pf:fit:L:Pol}.
\end{proof}

\begin{theorem}[Conditional slope bound]\label{pf:fit:flb}
Let $S$ be an indefinite datum over a well-formed definite side, satisfying the row conditions, and assume that
$\fitRR$ is a lattice with $\fitRR_{(w)}=(\iota_{\mu(w)}(\Lambda_{\mu(w)}))_{(w)}$ for every $w$. Let $O$ be a
functional datum and $C,x,C_H,b$ real numbers with $C\ge1$, $x\ge1$, $1\le\log b$ and $\log b+\log x\le3\log b$, such
that the leaf conditions \ref{pf:fit:L:MM}--\ref{pf:fit:L:Pol} hold. Then $\fitmumax(\fitRR)\le x^{C_H}b-\fitmlog(\kappa)$;
in particular $S$ satisfies $\fitFUN(N)$ for all numbers $N$ with $hh=x^{C_H}b$, and $\fitFUN(N_Z(C_0,C_H))$ for every
$C_0$ and every scale $Z$ with entries $x$ and $b$.
\formallabel{Formalism.L2.Indef.Tree.functional_lower_bound_of_conditions}
\formallabel{Formalism.L2.Indef.Tree.functional_lower_bound_of_conditions_numeric}
\end{theorem}

\begin{proof}
\emph{The map.} Let $\Psi$ be given by \ref{pf:fit:L:C}. Since $r_{\mu_0}(\eta)=1$ by \ref{pf:fit:L:Sq},
$\Psi\iota P=V_{\mu_0}$, and by \ref{pf:fit:L:Ker} and \cref{pf:fit:psiinj} (with the scalar conditions and the
surjectivity of $P$) $\Psi$ is injective.

\emph{Denominators.} Let $T=\mathrm{Tw}\cup T_f$ and $a(w)=v_w\bigl(\prod_\eta r_{\mu(w)}(\eta)\bigr)+d_F(w)$. For every
$w$, \cref{pf:fit:hencelocal} with $t=\mu(w)$, $c=\prod_\eta r_{\mu(w)}(\eta)\ne0$ (by \ref{pf:fit:L:Sq}), the
identity of \ref{pf:fit:L:C} and the inclusion of \ref{pf:fit:L:FP} gives $(\Psi(\fitRR))_{(w)}\subseteq
(w^{-a(w)}M)_{(w)}$. If $w\notin T$, then $\mu(w)=\mu_0$ by the row conditions, so the ratios are $1$, $d_F(w)=0$ and
$a(w)=0$. By \cref{pf:fit:locglob}, $\Psi(\fitRR)\subseteq\fitden_{T,a}M$.

\emph{Height.} By \ref{pf:fit:L:IP}, applied to $\nu=\Psi\iota$ (which satisfies $\nu P=V_{\mu_0}$), and
\cref{pf:fit:ip}(b), $h(\Psi)\le B_{\mathrm{IPa}}+3CfmN^2L=B_{\mathrm{IP}}$ for the frames $X_{E,\kappa}$ and $X'$.

\emph{Remaining bounds.} $\sum_{w\in T}\ell_wd_F(w)=\sum_{w\in T_f}\ell_wd_F(w)\le B_{\mathrm{FP}}/2$ by \ref{pf:fit:L:FP}.
By \ref{pf:fit:L:Int}, \ref{pf:fit:L:Arch} and \cref{pf:fit:prodformula}(c), $\fitht_0(z_A(\eta))$ and
$\fitht_0(z'_A(\eta))$ are at most $B_{\mathrm{Ar}}+6B_\Delta\le\bar B$; by \ref{pf:fit:L:Thb} the heights of $z_T$ and
$z_{Td}$ are at most $B_{\mathrm{TH}}(x)\le\bar B$. Moreover $B_{\mathrm{TH}}(x)\ge0$ since $b\ge mfN^2\ge0$ and
$\log b+\log x\ge0$.

\emph{Conclusion.} \Cref{pf:fit:passage}, applied with $R=\fitRR$, $E_M=E_M(C)$ (by \ref{pf:fit:L:MM}),
$H=B_{\mathrm{IP}}$, $B_F=B_{\mathrm{FP}}/2$, $r(w,\eta)=r_{\mu(w)}(\eta)$, $B=\bar B$, $B_N=B_{\mathrm{TH}}(x)$
(\ref{pf:fit:L:Den}), $B_\delta=B_\Delta$ (\ref{pf:fit:L:Int}), the local bounds of \ref{pf:fit:L:Sq} and
$P=x^{C_H}b$ (\cref{pf:fit:size}), gives $\fitmumax(\fitRR;X_E)\le-\fitmlog(\kappa)+x^{C_H}b$. The numbers
$N_Z(C_0,C_H)$ have $hh=x^{C_H}b$ for every $C_0$.
\end{proof}

% ---------------------------------------------------------------------------
\subsection{Fitted functions and the reduction to sixteen leaves}\label{sec:pf:fit:assembly}

We now attach functional data to the points of the families. Fix a point $p$ of a family, with datum $D$, coefficient
field $K=\mathit{LL}$ (totally complex), real datum $R$ and first side $S_{\mathrm f}$ over the definite side
$\mathcal D_p$ of $p$; let $C_F\in\R$.

\begin{definition}[Fitted functions]\label{pf:fit:funin}
A \emph{function datum} $\phi$ at $(p,C_F)$ consists of a grid datum $G=G_\phi$ at level $(1,C_F,b)$, a Weil datum $W$ over
$G$ which is over $R$, an integer $n$, a frame $X'$ on $K^n$, an $\OO_K$-submodule $M\subseteq K^n$, matrices
$V_t\in M_{n\times n_V(t)}(K)$ for the types $t$, square ratios $r_t(\eta)\in K$ ($\eta\in\fitSq$, $m=m(D)$), and
data $d_l,S_d,S_e,N_0,z_A,z'_A,z_T,z_{Td},g_T,c,B_w$ as in \cref{pf:fit:fundata}. For $\zeta$ a function on types
with values in $K^\times$, a finite set $T_f$ and $d_F$, the functional datum $\phi(\zeta,T_f,d_F)$ is $O$ with the entries of
$\phi$, the matrices $\zeta_t^{-1}V_t$, the given $T_f,d_F$, and the following entries fixed by $D$ and $G$: $m,f,k,t$ are the
ranks of $D$; $d(i,j)=2r+(2j+1)t+(2i+1)s$ (indices from $0$) with $r=k+m-1$ and $s=k+m+5$; $h=d/2$; $y(i,j)=m-1-j$; $S_0$
is the set of places of $G$, $q_*=v_0$, $c_0(v_0)=1$, $c_0(q)=2$ otherwise, $q(q)=\Nm q$; $\mathrm{cl}$ is the cloud of
Harish-Chandra parameters of $D$; $g_1,g_2$ are the type primes, $p_d,p'_d$ the auxiliary primes of the host and $d_E$
the discriminant of $E$. We write $O_\phi=\phi(1,\emptyset,0)$.
\end{definition}

A function datum is tied to a choice of the remaining data (\cref{sec:pf:fit:setting}). For a choice $a$ for $(G,W,n)$
the \emph{tie} $\sfname{TieAt}(\phi,a)$ asserts: $\mathrm{OK}(a)$; the \emph{test condition}: $X'$ is a metric of the inverse
transpose of the Gram matrix of the forms of the choice and $M$ is the dual of the lattice of these forms; the
\emph{value condition}: for every type $t$ in use, every $x\in K^{n_V(t)}$ and every $b\le n$,
\begin{equation}\label{eq:pf:fit:tieV}
\sigma_{v_0}\bigl((V_tx)_b\bigr)=v_K\bigl(\mathrm{ch}(a),t,e_{V_0}(x)\bigr)_b,
\end{equation}
where $v_K(\mathrm{ch},t,\cdot)$ is the vector of class sums of the choice and $e_{V_0}(x)$ the form attached to $x$; and
the \emph{square condition}: the square data of $\phi$ are a square datum of the choice
($\sfname{IsSqOf}$, \cref{sec:pf:fit:inputs}). Let $\mathcal T(p,C_F,\phi)$ be the statement that $\sfname{TieAt}(\phi,a)$ holds for
some choice $a$, and $\mathcal T^{\fitFit}(p,C_F,\phi)$ the statement that $G$ fits $p$ and $\mathcal T(p,C_F,\phi)$.

\begin{definition}[Leaf statements]\label{pf:fit:leafst}
For a property $Q(p,C_F,\phi,C,C_H)$ of points, function data and constants, and a property $\Xi(p,C_F,\phi)$, let
$\mathcal L(\Xi,Q)$ be the statement: there is $C_{F0}$ such that for every $C_F\ge C_{F0}$
\begin{multline*}
\fitClosed\Bigl(\sfname{FamHyp},\ (\mathcal F,C_H)\mapsto\Bigl(\sfname{DIVinN}(\mathcal F)\Rightarrow\exists C_0\
\forall C\ge C_0:\\ \fitEv\bigl(\mathcal F,\ p\mapsto\forall\phi:\ \Xi(p,C_F,\phi)\Rightarrow Q(p,C_F,\phi,C,C_H)\bigr)\Bigr)\Bigr).
\end{multline*}
With $\mathrm{True}$ the trivial property, the sixteen \emph{leaves} are the statements $\mathcal L(\mathcal T^{\fitFit},Q)$ (or
$\mathcal L(\mathrm{True},Q)$ where indicated) for the following $Q$, where $O=O_\phi$, $\mu$ is the map of $S_{\mathrm f}$,
and $x,b$ are the entries of the scale of $p$:
\begin{enumerate}[label=(F\arabic*),leftmargin=*]
\item\label{pf:fit:F:MM} minimal mass: \ref{pf:fit:L:MM} for $O$ and $C$;
\item\label{pf:fit:F:C} unique map: there is $\zeta$ with $\zeta_{\mu_0}=1$, each $\zeta_t$ a root of unity, such that
\ref{pf:fit:L:C} holds for $S_{\mathrm f}$ and $\phi(\zeta,\emptyset,0)$;
\item\label{pf:fit:F:Ker} kernel: \ref{pf:fit:L:Ker} for $O$;
\item\label{pf:fit:F:FP} fixed point: there are $T_f,d_F$ such that \ref{pf:fit:L:FP} holds for $\phi(1,T_f,d_F)$, $\mu$, $C$ and
$b$;
\item\label{pf:fit:F:IPa} height of the map: the first half of \ref{pf:fit:L:IP} for $\phi$ (with the matrix $V_{\mu_0}$ of
$\phi$);
\item\label{pf:fit:F:IPb} infinite places, with $\Xi=\mathrm{True}$: the second half of \ref{pf:fit:L:IP} for $S_{\mathrm f}$
and the constants $f,m,N$ of $O$;
\item\label{pf:fit:F:Sq} unique square root: \ref{pf:fit:L:Sq};
\item\label{pf:fit:F:Den} density: \ref{pf:fit:L:Den};
\item\label{pf:fit:F:Int} integrality: \ref{pf:fit:L:Int};
\item\label{pf:fit:F:Arch} archimedean condition: \ref{pf:fit:L:Arch};
\item\label{pf:fit:F:Thb} theta bound: \ref{pf:fit:L:Thb};
\item\label{pf:fit:F:Wt} weight bound: \ref{pf:fit:L:Wt};
\item\label{pf:fit:F:DesH} descent height, with $\Xi=\mathrm{True}$: \ref{pf:fit:L:Const} and \ref{pf:fit:L:Lc};
\item\label{pf:fit:F:Rank} rank gap, with $\Xi=\mathrm{True}$: \ref{pf:fit:L:Size};
\item\label{pf:fit:F:Ram} ramified sums, with $\Xi=\mathrm{True}$: \ref{pf:fit:L:Mod} and \ref{pf:fit:L:CS};
\item\label{pf:fit:F:Pol} descent polynomial, with $\Xi=\mathrm{True}$: \ref{pf:fit:L:Pol} with $C_H$;
\end{enumerate}
all conditions for $O$ and the constants $(C,x,C_H,b)$. Moreover:
\begin{itemize}
\item $\sfname{ExistsFunctionIn}(\mathrm{Fld},\mathcal T)$: there is $C_{F0}$ such that for every $C_F\ge C_{F0}$,
closedly on $\sfname{FamHyp}$, under $\sfname{DIVinN}$, eventually at the points $p$: for every grid datum $G$ at level
$(1,C_F,b)$ and every Weil datum $W$ over $G$ and over $R$ such that $G$ fits $p$ and $\mathrm{Fld}(p,C_F,G,W)$, there is a
function datum $\phi$ with $G_\phi=G$ and $\mathcal T(p,C_F,\phi)$;
\item $\sfname{ScaleCondition}$: closedly on $\sfname{FamHyp}$, eventually at the points, $1\le x$, $1\le\log b$ and
$\log b+\log x\le3\log b$;
\item $\sfname{BundleCondition}$: at every point $p$ with $\fitPC(p)$, the definite side $\mathcal D_p$ is well formed and
$S_{\mathrm f}$ satisfies the row conditions.
\end{itemize}
\end{definition}

\begin{lemma}\label{pf:fit:leafand}
If $\mathcal L(\Xi,Q)$ and $\mathcal L(\Xi,Q')$, then $\mathcal L(\Xi,Q\wedge Q')$. If $\mathcal L(\mathrm{True},Q)$,
then $\mathcal L(\Xi,Q)$ for every $\Xi$.
\formallabel{GlobalTheory.BigTop.Fun.LeafConditions.and}
\formallabel{GlobalTheory.BigTop.Fun.LeafConditions.of_noTie}
\end{lemma}

\begin{proof}
Take the larger of the two thresholds $C_{F0}$, the conjunction of the two closed statements (closed statements are
stable under conjunction, \lref{Formalism.MainNotions.Closed.and}), the larger of the two constants $C_0$ and the
larger of the two thresholds $x_0$ of $\fitEv$. The second statement is monotonicity in $\Xi$.
\end{proof}

\begin{lemma}\label{pf:fit:rootunity}
Let $O$ satisfy \ref{pf:fit:L:FP} for $\mu,C,b$, and let $\zeta_t\in K$ be roots of unity. Then the functional datum
obtained by replacing every $V_t$ by $\zeta_t^{-1}V_t$ satisfies \ref{pf:fit:L:FP} for $\mu,C,b$.
\formallabel{GlobalTheory.finitePlaceBound_inv_rootOfUnity_smul}
\formallabel{GlobalTheory.matLat_smul_le_of_isIntegral}
\end{lemma}

\begin{proof}
$\zeta_t^{-1}$ is a root of unity, hence integral over $\Z$, so $\zeta_t^{-1}V_t(M_t)\subseteq V_t(M_t)$, the latter being an
$\OO_K$-module; localisation preserves inclusions.
\end{proof}

\begin{theorem}[Reduction to the leaves]\label{pf:fit:leaves}
Let $\mathrm{Fld}$ be a property of grid and Weil data at points. Assume $\sfname{BundleCondition}$,
$\sfname{ScaleCondition}$, $\sfname{ExistsFunctionIn}(\mathrm{Fld},\mathcal T)$ and the sixteen leaves
\ref{pf:fit:F:MM}--\ref{pf:fit:F:Pol}. Then $\sfname{ConstantFunctionFits}(\mathrm{Fld})$ holds.
\formallabel{GlobalTheory.BigTop.Fun.fun0Fit_of_leaves_alt}
\end{theorem}

\begin{proof}
By \cref{pf:fit:leafand} the sixteen leaves combine into one statement $\mathcal L(\mathcal T^{\fitFit},Q_{\mathrm{all}})$ with
threshold $a$, where $Q_{\mathrm{all}}$ is the conjunction of \ref{pf:fit:F:MM}--\ref{pf:fit:F:Pol}; let $b'$ be the
threshold of $\sfname{ExistsFunctionIn}$, and $C_{F0}=\max(a,b')$. Fix $C_F\ge C_{F0}$. Closedly on $\sfname{FamHyp}$ we
have the conjunction of the three statements; fix a large family $\mathcal F$, a value $C_H$, and assume
$\sfname{DIVinN}(\mathcal F)$. Let $C_0$ be the constant of the combined leaf statement and put $C_1=\max(C_0,1)$.
Eventually at the points $p$ (with $\fitPC(p)$) we then have: all leaves hold for $C_1$ and every function datum $\phi$ with
$\mathcal T^{\fitFit}(p,C_F,\phi)$; the conclusion of $\sfname{ExistsFunctionIn}$; and the scale condition.

Let $p$ be such a point with $\sfname{HasGrid}(\mathrm{Fld},p,C_F)$: there are $G,W$ over $R$ with $G$ fitting $p$ and
$\mathrm{Fld}(p,C_F,G,W)$. By $\sfname{ExistsFunctionIn}$ there is $\phi$ with $G_\phi=G$ and $\mathcal T(p,C_F,\phi)$; as $G$ fits $p$,
$\mathcal T^{\fitFit}(p,C_F,\phi)$ holds and the leaves hold for $\phi$ and $C_1$. Let $\zeta$ be given by \ref{pf:fit:F:C} and
$T_f,d_F$ by \ref{pf:fit:F:FP}, and consider $O=\phi(\zeta,T_f,d_F)$. By $\sfname{BundleCondition}$ the definite side is well
formed and $S_{\mathrm f}$ satisfies the row conditions, and by \cref{pf:fit:glue} (applied with $L_0=\iota_{\mu_0}
(\Lambda_{\mu_0})$; the local lattices differ from it only at the finitely many twisted places) $\fitRR$ is a lattice with
the local components of \cref{pf:fit:flb}. We check the leaf conditions for $O$ with the constant $C_1$. All of them
except \ref{pf:fit:L:C}, \ref{pf:fit:L:Ker}, \ref{pf:fit:L:FP} and the first half of \ref{pf:fit:L:IP} involve only
entries of $O$ which coincide with those of $O_\phi$, and hold by the leaves. Condition \ref{pf:fit:L:C} only involves the
matrices $\zeta_t^{-1}V_t$ and the ratios, and holds by \ref{pf:fit:F:C}. Since $\zeta_{\mu_0}=1$, the matrix of $O$ at
$\mu_0$ is $V_{\mu_0}$, so \ref{pf:fit:L:Ker} and the first half of \ref{pf:fit:L:IP} hold by \ref{pf:fit:F:Ker} and
\ref{pf:fit:F:IPa}. Finally \ref{pf:fit:L:FP} for $\phi(\zeta,T_f,d_F)$ follows from \ref{pf:fit:F:FP} for $\phi(1,T_f,d_F)$ by
\cref{pf:fit:rootunity}. The scale condition
gives $x\ge1$, $\log b\ge1$ and $\log b+\log x\le3\log b$, and $C_1\ge1$. By \cref{pf:fit:flb}, $S_{\mathrm f}$
satisfies $\fitFUN(N_p(C,C_H))$ for every $C$.
\end{proof}

% ---------------------------------------------------------------------------
\subsection{The inputs}\label{sec:pf:fit:inputs}

The leaves are deduced from statements about grid data and about choices, which hold uniformly for all data
of large scale. We first fix the quantifiers.

\begin{definition}[Uniform statements]\label{pf:fit:uniform}
Let $Q(G,\mu_0,C,x,b)$ be a property of grid data with constants. $\fitGBig(Q)$ is the statement: there is $C_{F0}$ such
that for all $\eps$ and $C_F\ge C_{F0}$ there are $C_0,x_0$ such that $Q(G,\mu_0,C,x,b)$ holds for all $C\ge C_0$, all
global data $D$ and size data with exponent $\eps$, scale $x\ge x_0$ and $b=e^{x^\eps}$ satisfying the \emph{standing
conditions} $\sfname{Rec}$ (the size rows $\sfname{SizeRows}$, the host hypothesis, $F$ the maximal real subfield of
$E$, and $|\mathrm{logDer}(1)|\le x^\eps$ for the logarithmic derivative at $1$ attached to the size data), all level
characters with base type $\mu_0$ and all grid data $G$ at level $(1,C_F,b)$. For a property $Q(\mathrm{ch})$ of choices,
$\sfname{Big}(Q)$ is the statement: there is $C_{F0}$ such that for all $\eps$ and $C_F\ge C_{F0}$ there is $x_0$ such
that $Q(\mathrm{ch})$ holds for all data $D$ and size data with exponent $\eps$ and scale $x\ge x_0$, all families $A$,
real data $R$ and faces $F$ in \emph{Hida setting} ($\sfname{HidaSetting}$: size rows, Eisenstein tie, host conditions,
and the conditions on the characters of the face), and all admissible choices $\mathrm{ch}$ at $R$ at scale $(C_F,b)$.
$\fitBigC(Q)$, for a property $Q(F,\mathrm{ch},C,b)$, is the analogous statement with the order of quantifiers
$\exists C_{F0}\,\forall C_F\ge C_{F0}\,\exists C_0\,\forall\eps,C\ge C_0\,\exists x_0$, and with the field condition
$\sfname{FieldOK}(\mathrm{ch})$ among the hypotheses.
\end{definition}

The following statements about points are used besides; ``closedly, eventually'' means
$\fitClosed(\sfname{FamHyp},(\mathcal F,C_H)\mapsto\fitEv(\mathcal F,\cdot))$.
\begin{itemize}
\item $\sfname{SizeAt}$: closedly, eventually, the size data $L_{x,i}$ of the member satisfy $\sfname{SizeRows}$.
\item $\sfname{SettingAtPoint}$: closedly, eventually, if the Eisenstein tie $\sfname{ESerTie}$ of $D_{x,i}$ holds, then
$L_{x,i}$, the real datum and the face of the point are in Hida setting.
\item $\sfname{LogDer}$: closedly, eventually, $|\mathrm{logDer}(1)|\le x^\eps$ for $L_{x,i}$.
\end{itemize}

\begin{lemma}[Sizes, setting and the logarithmic derivative]\label{pf:fit:setting}
$\sfname{SizeAt}$, $\sfname{SettingAtPoint}$ and $\sfname{LogDer}$ hold.
\formallabel{GlobalTheory.T.size_at'}
\formallabel{GlobalTheory.BigTop.Fun.SizeProof.sizeAt_alt}
\formallabel{GlobalTheory.T.set_at_of_size_at'}
\formallabel{GlobalTheory.BigTop.Fun.setAt_of_sizeAt_alt}
\formallabel{GlobalTheory.BigTop.Fun.SetFromSize.setting_of_size_alt}
\formallabel{GlobalTheory.BigTop.Fun.SetFromSize.finite_tup}
\formallabel{GlobalTheory.T.abs_logDerivativeAtOne_le_x_rpow_eps}
\formallabel{GlobalTheory.logDeriv_LFunction_at_one_le_along_family_alt}
\end{lemma}

\begin{proof}
$\sfname{SizeAt}$: for a family of $\sfname{FamHyp}$ with exponent $\eps$ and $x\ge\max(44^4,8^{4/\eps})$, the size rows
of $L_{x,i}$ follow from the parameter choice of the family, its fixed constants and size conditions, and the rank and
scale windows \eqref{eq:families:rankval} (\lref{GlobalTheory.BigTop.Fun.SizeProof.sizeRows_fam}: a direct
verification of the numerical inequalities defining the size rows, from bounds for the Harish-Chandra parameters, the
gap and the length of the cloud). $\sfname{SettingAtPoint}$: given the size rows and the Eisenstein tie, the remaining
conditions of the Hida setting are part of the data of the point: its host condition, the conditions on the map $\mu$, the
condition on its group of characters, the guard condition $\sfname{QGuard}$ of its real datum, the standard characters, and
the core clauses of the Hecke choice, which follow from the weak characterisation of the characters of the point
(\lref{GlobalTheory.OneMapBridge.heckeChoiceClauses_of_weakCharacterisation_at_face}: the weak characterisation at a
face, together with the existence condition of the real datum, gives these clauses and the conditions that $\mu$ equals
the base type off the twisted places and that every type is the base type or a value of $\mu$). The set of types is
finite: by the last two conditions every type is the base type or the image under $\mu$ of one of the finitely many
twisted places. $\sfname{LogDer}$: the bound $|\mathrm{logDer}(1)|\le x^\eps$ is one of the conditions of the parameter
choice of every member of a family in $\sfname{FamP}(\eps)$, so it holds for all $x>x_1$; the hypothesis $F$ of the
formal statement is not used.
\end{proof}

\paragraph{The square statements.}
The following twelve properties of grid data (used through $\fitGBig$) concern the ``square'' of a grid: for index
pairs $\eta=(i,j)$ and types $t$, the constants $\gamma_\eta(t)$ of the matched square theta integrals, the Hecke characters
$\nu_H(G,t,i,j)$ and $\mathrm{fam}(G,\mu_0,i,j,c)$ attached to the grid, their completed $L$-functions $\Lambda$, the
\emph{toric values} (values $\Lambda(s)$ at integers $1\le s\le h(\eta)$, $h(\eta)$ the half rank of $\eta$, normalised by
CM periods of a lattice $\Lambda_d$ of discriminant $-d$), their \emph{denominator data} (functions $n$ on primes), and
\emph{bottom pairs} $(P,P')$ and \emph{bottom triples} of theta data at the bottom of the tower, with norms $\|P\|$. We refer
to \cite{LeanSource2026} for the definitions of these objects.

The \emph{square datum} $\sfname{IsSqOf}$ of the square condition of a tie is defined as follows. For a choice with Weil datum
$W$, grid choice and real match, and for the base type $\mu_0$ and the map $\mu$ of the point, a family of square data
$(r,z_A,z'_A,z_T,z_{Td},g_T,c,N_0,S_d,S_e,d_l,B_w)$ is a \emph{square datum} if: for every $\eta$ there are non-zero
constants $\gamma_\eta(t)$ of the matched square integrals at $\mu_0$ and at the types $\mu(w)$ with $\sigma(r_t(\eta))\gamma_\eta(\mu_0)=\gamma_\eta(t)$;
$z_A(\eta)$ and $z'_A(\eta)$ map to the norms of a bottom pair at $\mu_0$; $z_T=1$; for $1\le s\le h(\eta)$ the numbers
$z_{Td}(\eta,c,s)$ are toric values at $(\eta,s)$ with valuation shift $g_T(\eta,s)$ and $N_0(\eta,s,\cdot)$ is a denominator
datum, while outside this range $z_{Td}=1$, $g_T=0$ and $N_0=0$; $c(w)=1$ if and only if $p_d\in w$; $S_d$ is the set of
primes $l$ with $N_0(\eta,s,l)\ne0$ for some $(\eta,s)$; $S_e$ is the set of primes below the places $w$ with $\mu(w)\ne\mu_0$;
$d_l$ is an integrality exponent of the grid; and $B_w=\max\mathrm{Wt}(D)$.

The twelve statements are, with $B_{\mathrm{TH}}=Cf(D)b(\log b+\log x)$:
\begin{enumerate}[label=(S\arabic*),leftmargin=*]
\item\label{pf:fit:S:DenSum} $\sfname{DenSum}$: every denominator datum $n$ of the toric value at $(\eta,s)$, $1\le
s\le h(\eta)$, is non-negative on primes and satisfies $\sum_{l\in S}n(l)\log l\le B_{\mathrm{TH}}$ for every finite set
$S$ of primes;
\item\label{pf:fit:S:Thb} $\sfname{ThbSt}$: if $z_d(c)$ ($c\in\{0,1\}$) are toric values at $(\eta,s)$ with valuation
shift $g$, then $\fitht_g(1)\le B_{\mathrm{TH}}$, $z_d(c)\ne0$ and $\fitht_g(z_d(c))\le B_{\mathrm{TH}}$;
\item\label{pf:fit:S:One} $\sfname{OneSqSt}$: for every Weil datum $W$ over $G$ with standard $\psi$, grid choice, real
match, square values at $\mu_0$ in $K$, and every square datum $(\mu,r,z_A,z'_A,z_T,z_{Td},g_T,\allowbreak c,N_0,S_d,S_e,d_l,B_w)$
satisfying the toric estimate $\sfname{ThA}$ below: $r_{\mu(w)}(\eta)\ne0$ and
$v_w(r_{\mu(w)}(\eta))\le\beta_w(\eta)$ for all $w,\eta$, with $\beta_w$ as in \eqref{eq:pf:fit:sqbound} for $d=2h$,
$y(i,j)=m-1-j$, the places of $G$ with their conductor exponents $c_0$, and $q=\Nm$;
\item\label{pf:fit:S:BotSmall} $\sfname{BotSmall}$: the level of $G$ is good, $G$ has no twisted places at $\mu_0$, and all
bottom pairs are \emph{small} with respect to $N(D)$ and the maximal weight $\max\mathrm{Wt}(D)$;
\item\label{pf:fit:S:BotNe} $\sfname{BotNe}$: both members of every bottom pair at $\mu_0$ have a non-zero coordinate;
\item\label{pf:fit:S:Field} $\sfname{FieldDat}$: $E$ is a CM field, abelian over $\Q$, the CM type of the host has the
field data of discriminant $-d$, the conductor of $G$ satisfies the base condition, and an integral period lattice of
discriminant $-d$ exists;
\item\label{pf:fit:S:TorEx} $\sfname{TorEx}$: at every $(\eta,s)$ and $c$ there is a non-zero toric value $x=\Lambda(s)$ for
the $L$-function of $\mathrm{fam}(G,\mu_0,i,j,c)$, whose period-normalised value $\beta(x)$ is algebraic;
\item\label{pf:fit:S:DenEx} $\sfname{DenEx}$: at every $(\eta,s)$ there is a denominator datum with finite support;
\item\label{pf:fit:S:BotEx} $\sfname{BotEx}$: bottom pairs at $\mu_0$ exist for every $\eta$;
\item\label{pf:fit:S:LNe} $\sfname{LNe}$: the $L$-functions of $\nu_H(G,t,i,j)$ do not vanish at the integers
$1-e_0(\eta)\le s\le h(\eta)$, and those of $\mathrm{fam}(G,\mu_0,i,j,c)$ do not vanish at the integers $s\ge1$;
\item\label{pf:fit:S:TriEx} $\sfname{TriEx}$: for every Weil datum with standard $\psi$ and every $\eta$ there is a bottom
triple with its constants;
\item\label{pf:fit:S:Tower} $\sfname{Tower}$: the ratio $\gamma_\eta(t)/\gamma_\eta(\mu_0)$ of square constants, multiplied by a product of
ratios $\Lambda(e+k+1)/\Lambda_0(e+k+1)$ of $L$-values for $\nu_H(G,t,i,j)$ and $\nu_H(G,\mu_0,i,j)$, equals the image of an
element of $K$ which is a unit away from the level, times a product of ``step ratios'', the ratio of the constants of the
bottom triples and a second product of $L$-value ratios (a Witt tower formula).
\end{enumerate}
The \emph{toric estimate} $\sfname{ThA}(G,\mu_0,\mu)$ asserts: for every finite place $w$ with $\mu(w)\ne\mu_0$, the level
element of $G$ at $\mu(w)$ is a $w$-adic unit, and for every prime $l\in w$, all $\eta$, $1\le s\le h(\eta)$, and $z\in K$
with $z\Lambda_l(s)=\Lambda(s)$, where $\Lambda$ is the $L$-function of $\nu_H(G,\mu(w),i,j)$ and $\Lambda_l$ that of
$\nu_H(G,\mu_0,i,j)$ twisted at $l$, one has $v_w(z)\le v_w(l)/l^2$.

\begin{theorem}[The square statements]\label{pf:fit:sqtop}
Howe duality at the real place, $F$, the functional equation of Hecke $L$-functions in absolute value form,
$\sfname{ToricDenominator}$, $\sfname{Transport}$, $\sfname{ZetaDenominators}$, $\sfname{Refl}$, $\sfname{ThetaIntegralsAgree}$
and $\sfname{LocalMatching}$ imply $\fitGBig$ of each of \ref{pf:fit:S:DenSum}--\ref{pf:fit:S:Tower}.
\formallabel{GlobalTheory.T.square_leading_term_bound}
\formallabel{GlobalTheory.SqTop.square_top_term_joined_estimates}
\formallabel{GlobalTheory.T.completedHeckeL_bounded_in_vertical_strips}
\formallabel{GlobalTheory.T.triEx_of_howe_duality_real}
\formallabel{GlobalTheory.T.tower_of_theta_matching_lNe}
\end{theorem}

\begin{proof}
The proof assembles the following results; we write each statement for its $\fitGBig$ version.
\begin{enumerate}[label=(\arabic*),leftmargin=*]
\item Completed Hecke $L$-functions are bounded in vertical strips
(\lref{GlobalTheory.HeckeStrip.completedLBoundedInStrips}), and the CM theta norm hypotheses $\sfname{Inp}(D)$
hold for every datum (\lref{GlobalTheory.T.cmThetaNormBoundHypotheses_satisfied}).
\item $\sfname{FieldDat}$ holds unconditionally (\lref{GlobalTheory.SqTop.fieldDat}).
\item $\sfname{LNe}$ follows from the functional equation, (1) and $\sfname{FieldDat}$
(\lref{GlobalTheory.SqTop.lNe_of_functionalEquation}), and so does the bound
$\log|\Lambda(s)|\le B_{\mathrm{TH}}$ for the $L$-functions of the characters with permuted weights, $\sfname{LBound}$
(\lref{GlobalTheory.T.heckeLValue_upperBound_permutedWeights_of_FE}).
\item $\sfname{TorEx}$ follows from $\sfname{Transport}$, $\sfname{FieldDat}$ and $\sfname{LNe}$
(\lref{GlobalTheory.ToricExist.torEx}); with $\sfname{Transport}$ and $\sfname{LBound}$ one obtains the bound $\sfname{BetaArch}$ for the
constant coefficient of the minimal polynomial of the normalised toric values,
$\log|a_0|\le\deg\cdot12B_{\mathrm{TH}}$ (\cref{pf:fit:betaarch}).
\item $\sfname{DenEx}$ and $\sfname{DenSum}$ follow from $\sfname{ZetaDenominators}$ and $\sfname{FieldDat}$
(\lref{GlobalTheory.DenExist.denEx}, which takes the denominator data given by the zeta denominators, and
\lref{GlobalTheory.DenSumPf.denSum}).
\item $\sfname{ThbSt}$ follows from $\sfname{ToricDenominator}$, $\sfname{FieldDat}$, $\sfname{TorEx}$, $\sfname{DenEx}$,
$\sfname{DenSum}$ and $\sfname{BetaArch}$ (\lref{GlobalTheory.T.toricValueHeightBound_of_denominatorsMeans}).
\item $\sfname{BotSmall}$ holds unconditionally (\lref{GlobalTheory.SqTop.BotSmallProof.botSmall_holds}); with
$F$, $\sfname{ThetaIntegralsAgree}$, $\sfname{LocalMatching}$, (1) and $\sfname{LNe}$ it gives the bottom constant statement
$\sfname{BotConst}$ (\cref{pf:fit:botconst}: the ratio of the constants of the bottom triples
at $t$ and $\mu_0$ is the image of a non-zero $\gamma\in K$ with $v_w(\gamma)\le v_w(z)+v_w(z')+6d_l(w)$ away from the
level, $z,z'$ the norms of the bottom pairs).
\item Bottom pairs exist at $\mu_0$ and at every type ($\sfname{BotEx}$, $\sfname{BotExT}$) by (1)
(\lref{GlobalTheory.T.bottomPairs_exist_baseTuple} and \lref{GlobalTheory.T.bottomPairsExistEveryTuple_of_cmThetaHypotheses}).
\item The norms of bottom pairs are, up to a non-zero constant independent of the type, products of $L$-values
$\prod_{1-e_0\le s\le h_0}\Lambda(s)$ ($\sfname{NormBig}$, from $F$, (1) and $\sfname{LNe}$, \cref{pf:fit:normbig}, through the
Rallis inner product formula \cite{Rallis1984}).
\item A lower bound for the valuations of ratios of $L$-values at the level primes, $\sfname{RhoLow}$, follows from
$\sfname{ToricDenominator}$, $\sfname{FieldDat}$ and $\sfname{TorEx}$ (\lref{GlobalTheory.T.lValueRatio_lower_bound_gauge}).
\item $\sfname{TriEx}$ follows from Howe duality, $F$, $\sfname{ThetaIntegralsAgree}$, $\sfname{LocalMatching}$ and
$\sfname{LNe}$ (\lref{GlobalTheory.SqTop.triEx_holds_alt}); $\sfname{BotNe}$ follows from these, (1), $\sfname{BotSmall}$ and
$\sfname{NormBig}$ (\lref{GlobalTheory.T.bottomTheta_nonvanishing}); and $\sfname{Tower}$ from $F$, the two statements
of \cref{thm:theorem-K} and $\sfname{LNe}$ (\lref{GlobalTheory.SqTower.tower_alt_bis}).
\item Finally $\sfname{OneSqSt}$ follows from (1), $\sfname{LNe}$, $\sfname{TorEx}$, $\sfname{BotSmall}$, $\sfname{BotExT}$,
$\sfname{NormBig}$, $\sfname{BotNe}$, $\sfname{TriEx}$, $\sfname{Tower}$, $\sfname{BotConst}$, $\sfname{Refl}$ and
$\sfname{RhoLow}$ (\lref{GlobalTheory.T.square_ratio_valuation_without_norm_formula}): by $\sfname{Tower}$ and
$\sfname{BotConst}$ the square ratio $r_t(\eta)$ is a product of $L$-value ratios, step ratios, a unit away from the level and
a ratio of bottom constants, whose valuations are bounded by the toric values, the denominator data and the norms of
bottom pairs, $\sfname{Refl}$ being used to pass from $\Lambda(s)$ to $\Lambda(1-s)$.\qedhere
\end{enumerate}
\end{proof}

\paragraph{Size bounds for class sums.}
For a choice $\mathrm{ch}$ with grid $G$ put $a_1=Cf(D)N(D)b\log b+Cm(D)N(D)^2\sigma(G)$ and
$b_1=Cf(D)m(D)N(D)^2(\log b+\sigma(G))b$, where $N(D)=k+m(t+s+1)$ and $\sigma(G)=\sum_{v\in S_0}\mathrm{depth}(v)\log\Nm v$
is the logarithmic depth sum of the grid; it is the number $\Sigma_0$ of \cref{pf:fit:fundata}, for the entries of
\cref{pf:fit:funin}, with $C$ replaced by $C_F$; since the exponent of the base place is non-decreasing in the constant,
$\sigma(G)\le\Sigma_0$ for $C\ge C_F$. A \emph{test lattice} of the
choice is a pair $(X',M)$ as in the test condition of the tie. The three size statements are $\fitBigC$ of:
\begin{itemize}
\item $\sfname{FinAt}$: there is $d_l\ge0$ with finite support and $\sum_{w\in F}\ell_wd_l(w)\le m\,b_1/2$ for every finite
$F$, such that for every test lattice $(X',M)$, every type $t$ and place $w$ with $w$ \emph{away} from the level at $t$, every
$x$ in the lattice of type $t$ and every $y\in K^n$ with $\sigma_{v_0}(y)=v_K(\mathrm{ch},t,e(x))$, there is $s\in K^\times$
with $v_w(s)\le d_l(w)$ and $sy\in M$;
\item $\sfname{AwayAt}$: every finite place $w$ is away from the level at the type $\mu(w)$;
\item $\sfname{InfAt}$: for every test lattice $(X',M)$, every frame $X$ of the hermitian form $H_V(\mu_0)$ and every $V$
with $\sigma_{v_0}(Vx)=v_K(\mathrm{ch},\mu_0,e(x))$, $h(V)\le m\,a_1$ for the frames $X$ and $X'$.
\end{itemize}

\begin{theorem}[Size bounds for class sums]\label{pf:fit:sumsize}
$F$ implies $\fitBigC(\sfname{FinAt})$, $\fitBigC(\sfname{AwayAt})$ and $\fitBigC(\sfname{InfAt})$.
\formallabel{GlobalTheory.T.sum_size_bound_all_places}
\formallabel{GlobalTheory.SumSize.Join.sum_size_finite_away_archimedean_combined}
\end{theorem}

\begin{proof}
The claims are defined in \cref{sec:pf:fit:sizes}. The hypotheses $\sfname{Inp}(D)$ hold for all $D$
(\lref{GlobalTheory.T.cmThetaNormBoundHypotheses_satisfied}); $\sfname{SumOKK}$ follows from $F$
(\lref{GlobalTheory.T.sumOKK_of_theta_local_inputs}) and $\sfname{DepthK}$ from $F$ and $\sfname{SumOKK}$ (\cref{pf:fit:depth});
a dictionary of local bounds follows from $F$ and $\sfname{Inp}$ (\lref{GlobalTheory.T.sum_size_dictionary}).
\emph{Finite places.} $\sfname{FinScal}$ (\lref{GlobalTheory.T.finite_and_infinite_scalar_bounds}), $\sfname{FinCoef}$
(\lref{GlobalTheory.Adopt.finCoefK_of_theta_inputs}), $\sfname{FinDat}$ (\cref{pf:fit:findat}), $\sfname{FinGam}$
(\lref{GlobalTheory.T.gamma_factor_loss_at_finite_places}), $\sfname{FinTop}$ and $\sfname{FinRow}$ (from $F$, $\sfname{Inp}$,
the dictionary and $\sfname{DepthK}$: \lref{GlobalTheory.T.sum_size_finite_top_term_integral} and \lref{GlobalTheory.T.sum_size_bound_finite_places}),
$\sfname{FinPair}$ (\cref{pf:fit:fieldpair}), $\sfname{FinPsi}$ (\lref{GlobalTheory.Adopt.test_vector_value_factorization})
and $\sfname{FinForm}$ (\cref{pf:fit:formval}) give $\fitBigC(\sfname{FinAt})$ by \cref{pf:fit:finasm}. $\fitBigC(\sfname{AwayAt})$
follows from $F$ (\lref{GlobalTheory.Adopt.twisted_local_factor_away_from_level}).
\emph{Infinite places.} $\sfname{InfScal}$ (\lref{GlobalTheory.T.finite_and_infinite_scalar_bounds}), $\sfname{InfTens}$
(\cref{pf:fit:tensor}), $\sfname{InfCount}$ (\cref{pf:fit:sumcount}), $\sfname{InfClass}$ (\cref{pf:fit:classcount}),
$\sfname{InfTop}$ (from $F$, $\sfname{Inp}$, the dictionary and $\sfname{DepthK}$:
\lref{GlobalTheory.T.kernel_conjugation_mean_archimedean_top_term}), $\sfname{InfRow}$ (\cref{pf:fit:opposite}),
$\sfname{InfPair}$ (\cref{pf:fit:fieldpair}), $\sfname{InfPsi}$ (\lref{GlobalTheory.Adopt.test_vector_value_factorization}) and
$\sfname{InfForm}$ (\cref{pf:fit:formval}) give $\fitBigC(\sfname{InfAt})$ by \cref{pf:fit:infasm}.
\end{proof}

\paragraph{Row properties.}
The five row properties are the $\sfname{Big}$ statements of \cref{sec:families:theta}: $\sfname{SumOKK}$ (the finite sums
of the choice are meaningful), $\sfname{ConstK}$ (the matched square integrals are non-zero and $c(i,j,t)\ne0$),
$\sfname{LinK}$ (all column--row pairings converge), $\sfname{FormulaK}$ (under $\sfname{FieldOK}$ there are $\kappa\in\C$
and roots of unity $\zeta_t\in K$ such that $\sigma(\zeta_t)c_{\mathrm{Tot}}(t)\,\psi_K(b,\Theta_t(j_Xx))=\kappa\,
v_K(t,x)_b$ for all good rational points $x$ and all $b$, where $c_{\mathrm{Tot}}(t)=\prod_{i,j}c(i,j,t)$, $\psi_K(b,\cdot)$ are
the column functionals of the choice and $\Theta_t$ the theta maps) and $\sfname{SumRatK}$ (under $\sfname{FieldOK}$ the class
sums lie in $\sigma(K)$); further $\sfname{InjK}$: if $\psi_K(b,u)=0$ for all $b$, then $u=0$.

\begin{lemma}[Row properties]\label{pf:fit:rows}
$F$, $\sfname{ThetaIntegralsAgree}$ and $\sfname{LocalMatching}$ imply $\sfname{SumOKK}$, $\sfname{ConstK}$, $\sfname{LinK}$,
$\sfname{FormulaK}$ and $\sfname{SumRatK}$; $F$ and $\sfname{FunctionalsSeparate}$ imply $\sfname{InjK}$.
\formallabel{GlobalTheory.T.local_sum_estimates_joined}
\formallabel{GlobalTheory.RowClaim.JoinF.joined_local_sum_estimates}
\formallabel{GlobalTheory.T.thetaLift_row_properties}
\formallabel{GlobalTheory.RowTest.twisted_convergence_with_real_places}
\end{lemma}

\begin{proof}
The first statement is \lref{GlobalTheory.RowClaim.Join.thetaLift_row_properties_conj} (the five row properties follow
from $F$, the two statements of \cref{thm:theorem-K} and three statements on square theta integrals), applied with the three
statements supplied by \lref{GlobalTheory.T.twisted_theta_integral_converges},
\lref{GlobalTheory.T.twisted_local_factors_match} and \lref{GlobalTheory.T.twisted_square_theta_integrals_eq}. For
the second, the twisted theta lifts are cuspidal by $F$ (\lref{GlobalTheory.T.twisted_thetaLift_cuspidal}); hence the
pairings with real places converge for all choices, which is the convergence statement required by
\lref{GlobalTheory.RowTest.injK_of_separates} (separation and convergence imply $\sfname{InjK}$).
\end{proof}

\paragraph{Toric estimate and square ratios.}
\begin{lemma}\label{pf:fit:toric}
\begin{enumerate}[label=(\alph*)]
\item $\sfname{TwistedTransport}$, $\sfname{SettingAtPoint}$, $\sfname{LogDer}$, $\fitGBig(\sfname{FieldDat})$ and
$\fitGBig(\sfname{TorEx})$ imply $\mathcal L(\mathcal T^{\fitFit},\ \sfname{ThA}(G_\phi,\mu_0,\mu))$.
\item $F$, $\sfname{ThetaIntegralsAgree}$, $\sfname{LocalMatching}$, $\sfname{SettingAtPoint}$, $\sfname{LogDer}$,
$\sfname{ConstK}$ and $\fitGBig$ of $\sfname{LNe}$, $\sfname{TriEx}$ and $\sfname{Tower}$ imply that, for large $C_F$,
closedly, eventually at the points, for all grids $G$, Weil data $W$ and choices $a$ with $\mathrm{OK}(a)$, the square ratios
$\mathrm{sqRatio}(\mathrm{ch}(a),i,j,t)$ lie in $\sigma(K)$.
\end{enumerate}
\formallabel{GlobalTheory.T.toric_ratio_family_estimate'}
\formallabel{GlobalTheory.ThmAUpper.thmA_upper_alt}
\formallabel{GlobalTheory.T.sq_ratio_mem_range_emb_ll}
\formallabel{GlobalTheory.BigTop.Fun.square_ratio_at_point_of_leading_term_alt}
\formallabel{GlobalTheory.BigTop.Fun.sqRatio_mem_of}
\end{lemma}

\begin{proof}
(a) $\sfname{TwistedTransport}$ gives the transport of toric ratios under $\Aut(\C)$
(\lref{GlobalTheory.ToricTransport.toricRatioConjugates_of_transport}). By \cref{pf:fit:transfer} below,
$\sfname{FieldDat}$ and $\sfname{TorEx}$ hold at the grids of tied function data. At a place $w$ with $\mu(w)\ne\mu_0$ the
level element is a $w$-adic unit (\lref{GlobalTheory.T.qn_not_mem_asIdeal_of_mu_ne}), and the bound
$v_w(z)\le v_w(l)/l^2$ is \lref{GlobalTheory.T.cm_abelian_field_data_of_toric_ratio_alt}, which expresses $z$ through
toric sums of the CM field, uses the transport of toric ratios and bounds the valuations of twisted theta values
(\lref{GlobalTheory.T.twistedThetaValue_valuation_le_div_sq}).
(b) By \cref{pf:fit:transfer}(b) the row property $\sfname{ConstK}$ holds at the choice, and the statements $\sfname{LNe}$,
$\sfname{TriEx}$, $\sfname{Tower}$ and the factorisation of the bottom triple integrals
(\lref{GlobalTheory.T.bottom_triple_integral_eq_const_mul_local}, from $F$, \cref{thm:theorem-K} and $\sfname{LNe}$) hold
at its grid. The square ratio is $c(i,j,t)/c(i,j,\mu_0)$, a ratio of constants of matched square integrals. By
$\sfname{Tower}$ it equals $\sigma(\zeta)$ times a product of step ratios, the ratio $\lambda(t)/\lambda(\mu_0)$ of the
constants of a bottom triple, and a quotient of products of $L$-value ratios $\Lambda(s)/\Lambda_0(s)$ at integers $1\le s\le
h(\eta)$, which are non-zero by $\sfname{LNe}$. The step ratios are by definition images of elements of $K$, the $L$-value
ratios lie in $\sigma(K)$ by the field condition, and so does $\lambda(t)/\lambda(\mu_0)$ by the rationality of the constants
of definite bottom triples, which follows from the factorisation of the bottom triple integrals, the positivity of the
hermitian forms involved and the rationality of the corresponding sums, which is part of the field condition.
\end{proof}

% ---------------------------------------------------------------------------
\subsection{Proof of the leaves}\label{sec:pf:fit:leafproofs}

\paragraph{Transfer.}
At a point $p$, a function datum $\phi$ with $\mathcal T^{\fitFit}(p,C_F,\phi)$ comes with a choice $a$ satisfying $\mathrm{OK}(a)$, in
particular the Eisenstein tie of the datum. The uniform statements therefore apply at $\phi$.

\begin{lemma}[Transfer]\label{pf:fit:transfer}
Assume $\sfname{SettingAtPoint}$ and $\sfname{LogDer}$.
\begin{enumerate}[label=(\alph*)]
\item If $\fitGBig(Q)$, then $\mathcal L\bigl(\mathcal T^{\fitFit},\ Q(G_\phi,\mu_0,C,x,b)\bigr)$.
\item If $\sfname{Big}(Q)$, then for large $C_F$, closedly, eventually at the points, $Q(\mathrm{ch}(a))$ holds for all grid
data $G$ at level $(1,C_F,b)$, Weil data $W$, $n$ and choices $a$ for $(G,W,n)$ with $\mathrm{OK}(a)$.
\end{enumerate}
\formallabel{GlobalTheory.T.forallLeaves_fitted_of_gBig}
\formallabel{GlobalTheory.BigTop.Fun.gBig_atTied_alt}
\end{lemma}

\begin{proof}
(a) Let $C_{F0}$ be the threshold of $\fitGBig(Q)$ and $C_F\ge C_{F0}$. Closedly on $\sfname{FamHyp}$ we have the conjunction
of $\sfname{SettingAtPoint}$ and $\sfname{LogDer}$. Given a large family with exponent $\eps$, let $C_0,x_0$ be the
constants of $\fitGBig(Q)$ for $(\eps,C_F)$, and let $C\ge C_0$. At a point over $(x,i)$, with $x$ larger than $x_0$ and the
thresholds of the two statements, and a function datum $\phi$ with a tie $a$: $\mathrm{OK}(a)$ contains the Eisenstein tie, so
the Hida setting holds; its size rows, host conditions and the identification of $F$ with the maximal real subfield, with
the bound of $\sfname{LogDer}$, are the standing conditions; the level character of the real datum has base type $\mu_0$;
and $G_\phi$ is a grid datum at level $(1,C_F,b)$. Hence $Q(G_\phi,\mu_0,C,x,b)$. (b) is proved in the same way, the Hida setting
being the hypothesis of $\sfname{Big}$.
\end{proof}

\paragraph{Leaves from the square statements.}
\begin{lemma}\label{pf:fit:leafsq}
Assume $\sfname{SettingAtPoint}$ and $\sfname{LogDer}$. Then:
\begin{enumerate}[label=(\alph*)]
\item $\fitGBig(\sfname{DenSum})$ implies the density leaf \ref{pf:fit:F:Den};
\item $\fitGBig(\sfname{ThbSt})$ implies the theta bound leaf \ref{pf:fit:F:Thb};
\item $\fitGBig(\sfname{BotSmall})$ and $\fitGBig(\sfname{BotNe})$ imply the integrality leaf \ref{pf:fit:F:Int};
\item $\fitGBig(\sfname{OneSqSt})$ and the toric estimate of \cref{pf:fit:toric}(a) imply the unique square root leaf
\ref{pf:fit:F:Sq};
\item the archimedean leaf \ref{pf:fit:F:Arch} holds.
\end{enumerate}
\formallabel{GlobalTheory.T.leafDensity_of_logDeriv_bound_and_denSum}
\formallabel{GlobalTheory.BigTop.Fun.den_tie_alt}
\formallabel{GlobalTheory.T.leaf_theta_bound_of_set_at}
\formallabel{GlobalTheory.BigTop.Fun.thb_tie_alt}
\formallabel{GlobalTheory.T.leafIntegrality_of_logDeriv_bound}
\formallabel{GlobalTheory.BigTop.Fun.int_tie_alt}
\formallabel{GlobalTheory.T.oneSq_of_log_derivative_bound_and_oneSqSt}
\formallabel{GlobalTheory.BigTop.Fun.oneSq_of_log_derivative_bound_thA}
\formallabel{GlobalTheory.BigTop.Fun.OneSqTie.lf_std_alt}
\formallabel{GlobalTheory.T.arch_tie_of_setAt_logDeriv}
\formallabel{GlobalTheory.BigTop.Fun.arch_tie_of_logDeriv_bound}
\end{lemma}

\begin{proof}
In each case we apply \cref{pf:fit:transfer}(a) and read the conclusion at the grid of $\phi$ through the square datum of the
tie; the entries $d,h,y,S_0,c_0,q,f,N$ of $O_\phi$ are those of \cref{pf:fit:funin}, and the half rank of $\eta$ is
$h(\eta)=d(\eta)/2$.

(a) For $1\le s\le h(\eta)$, $N_0(\eta,s,\cdot)$ is a denominator datum, and $S_d$ consists of primes; $\sfname{DenSum}$ gives
\ref{pf:fit:L:Den} with $B_{\mathrm{TH}}(x)=Cfb(\log b+\log x)$.

(b) For $1\le s\le h(\eta)$, $\sfname{ThbSt}$ applied to the toric values $z_{Td}(\eta,\cdot,s)$ with shift $g_T(\eta,s)$ gives
$z_{Td}(\eta,c,s)\ne0$, $\fitht_{g_T}(z_{Td}(\eta,c,s))\le B_{\mathrm{TH}}$ and, as $z_T=1\ne0$,
$\fitht_{g_T}(z_T(\eta,s))\le B_{\mathrm{TH}}$; in particular $B_{\mathrm{TH}}\ge0$ (take $s=1\le h(\eta)$). Outside this range
$z_T=z_{Td}=1$ and $g_T=0$, so the heights are $0\le B_{\mathrm{TH}}$.

(c) The statement used is $\fitGBig$ of $\sfname{BotSmall}\wedge\sfname{BotNe}$ together with $1\le C_F\le C$ (by taking
$C_0\ge C_F$). The bottom pairs of the square datum are small and have non-zero coordinates, and the integrality exponent
$d_l$ of the grid provides, for such pairs, the integrality bound of \ref{pf:fit:L:Int} with $B_\Delta$ replaced by
$C_FX$, where $X=2fN\log(2fN)+2fNB_w\log b+4fN^2B_w\log(\mathrm{lev}(G))$, $\log\mathrm{lev}(G)=c_S$ and $B_w=\max\mathrm{Wt}(D)$.
Since $B_\Delta=CX$, $X\ge0$ and $C\ge C_F$, \ref{pf:fit:L:Int} follows.

(d) The Weil datum of a fitted point has the standard additive character (a condition of the Hida setting). Applying
$\sfname{OneSqSt}$ to the square datum of the tie, whose map $\mu$ satisfies $\sfname{ThA}$ by \cref{pf:fit:toric}(a), gives
$r_{\mu(w)}(\eta)\ne0$ and the bound $v_w(r_{\mu(w)}(\eta))\le\beta_w(\eta)$. The square datum gives $r_{\mu_0}=1$ and that
$S_e$ consists of primes. Finally $d(i,j)=2r+(2j+1)t+(2i+1)s$ is even ($t$ and $s$ being odd), $d(i,j)\le
2(r+(j+1)t+(i+1)s)\le2N$, $y(i,j)+1=m-j\le m$, and $c_0(q)\le e(q)$ since $c_0(q_*)=1\le e(q_*)$ and $c_0(q)=2=e(q)$ otherwise.

(e) By \cref{pf:fit:archgrid}, $\fitGBig$ holds for the statement: if
$z_A,z'_A$ map to the norms of bottom pairs at $\mu_0$, then $\fitmi{\log^+|\sigma_vz_A(\eta)|}$ and
$\fitmi{\log^+|\sigma_vz'_A(\eta)|}$ are at most
$2CfNB_w(\log b+\log(2NB_w))+CfN^2(\log b+\log(2fNB_w))+CN^2\log\mathrm{lev}(G)$ with $B_w=\max\mathrm{Wt}(D)$. This
expression is $B_{\mathrm{Ar}}$, and the square datum identifies $z_A,z'_A$ with norms of bottom pairs.
\end{proof}

\paragraph{Leaves from the size bounds.}
\begin{lemma}\label{pf:fit:leafsize}
Assume $\sfname{SettingAtPoint}$.
\begin{enumerate}[label=(\alph*)]
\item $\fitBigC(\sfname{FinAt})$ and $\fitBigC(\sfname{AwayAt})$ imply the fixed point leaf \ref{pf:fit:F:FP}.
\item $\sfname{BundleCondition}$ and $\fitBigC(\sfname{InfAt})$ imply the leaf \ref{pf:fit:F:IPa}.
\item $\sfname{SizeAt}$ implies the weight leaf \ref{pf:fit:F:Wt}.
\end{enumerate}
\formallabel{GlobalTheory.T.leafFixedPoint_of_finK_awayK}
\formallabel{GlobalTheory.Adopt.leafFixedPoint_of_finK_awayK}
\formallabel{LeafFixedPointPairing.GlobalTheory.CompanionTheory.LeafFixedPointPlaceNorm.fp_tie_alt}
\formallabel{LeafFixedPointPairing.GlobalTheory.CompanionTheory.LeafFixedPointPlaceNorm.leaf_at}
\formallabel{LeafFixedPointPairing.GlobalTheory.CompanionTheory.LeafFixedPointPlaceNorm.local_clause}
\formallabel{LeafFixedPointPairing.GlobalTheory.CompanionTheory.LeafFixedPointPlaceNorm.sum_clause}
\formallabel{LeafFixedPointPairing.GlobalTheory.CompanionTheory.LeafFixedPointPlaceNorm.sig_le}
\formallabel{GlobalTheory.FinitePlacesProbe.sigma_mono_in_constant}
\formallabel{GlobalTheory.T.leafHeightMapBound_of_infK}
\formallabel{GlobalTheory.QuotientNormCoisometry.leafHeightMapBound_tie_of_infK}
\formallabel{GlobalTheory.QuotientNormCoisometry.htMap_quot}
\formallabel{GlobalTheory.QuotientNormCoisometry.norm_quot}
\formallabel{GlobalTheory.QuotientNormCoisometry.bound_le}
\formallabel{GlobalTheory.T.weight_of_size_bound'}
\formallabel{GlobalTheory.BigTop.Fun.wt_tie_alt}
\formallabel{GlobalTheory.BigTop.Fun.maxWt_le}
\end{lemma}

\begin{proof}
As in \cref{pf:fit:transfer}, the size statements apply at the choice $a$ of the tie, for $C$ larger than $C_F$, $0$ and their
constants $C_0$; note $\sigma(G)\le\Sigma_0$ for $C\ge C_F$ because the depth of the base place is non-decreasing in the
constant.

(a) Let $d_l$ be given by $\sfname{FinAt}$, $T_f$ its (finite) support and $d_F(w)=\lfloor d_l(w)\rfloor$; then $d_F=0$
off $T_f$. Local inclusion: let $z\in(V_{\mu(w)}(M_{\mu(w)}))_{(w)}$, say $s_0z=V_{\mu(w)}x$ with $x\in M_{\mu(w)}$ and
$s_0\notin w$. By the value condition \eqref{eq:pf:fit:tieV}, $y=V_{\mu(w)}x$ satisfies the hypothesis of $\sfname{FinAt}$ at
$t=\mu(w)$, and $w$ is away from the level at $\mu(w)$ by $\sfname{AwayAt}$; so there is $s\in K^\times$ with
$v_w(s)\le d_l(w)$, hence $v_w(s)\le d_F(w)$, and $sy\in M$. Then $y=s^{-1}(sy)\in(w^{-d_F(w)}M)_{(w)}$ and so is $z$.
Sum: $2\sum_{w\in T_f}\ell_wd_F(w)\le2\sum_{w\in T_f}\ell_wd_l(w)\le m\,b_1\le Cfm^2N^2(\log b+\Sigma_0)b\le B_{\mathrm{FP}}$.

(b) Let $\nu P=V_{\mu_0}$. Choose a frame $X$ of the positive definite form $H_V(\mu_0)$, i.e.\
$(X(v)^{-1})^*X(v)^{-1}=\sigma_v(H_V(\mu_0))$. Since $X_Q$ is the quotient metric, $U=X_Q(v)^{-1}\sigma_v(P)X(v)$ satisfies
$UU^*=X_Q(v)^{-1}\sigma_v(PH_V^{-1}P^*)X_Q(v)^{-*}=1$, so $U$ is a co-isometry and $\|BU\|=\|B\|$ for every matrix $B$
($\|BU\|^2=\|BUU^*B^*\|$). As $X'(v)^{-1}\sigma_v(\nu P)X(v)=(X'(v)^{-1}\sigma_v(\nu)X_Q(v))U$, the operator norms of $\nu$ for
$(X_Q,X')$ and of $\nu P=V_{\mu_0}$ for $(X,X')$ coincide at every $v$, and $h(\nu)=h(V_{\mu_0})$. By the value condition,
$\sfname{InfAt}$ applies to $V_{\mu_0}$: $h(V_{\mu_0})\le m\,a_1=CfmNb\log b+Cm^2N^2\sigma(G)\le B_{\mathrm{IPa}}$.

(c) The square datum gives $B_w=\max\mathrm{Wt}(D)$, the maximum of $1$ and the absolute values of the weights $w_A,w'_A$ of
the datum. By the size rows these are bounded by $2e+1+\mathrm{gap}$ and $2\,\mathrm{cen}+1+\mathrm{gap}$ for entries $e$,
$\mathrm{cen}$ of the Harish-Chandra data, which gives $1\le B_w\le10b$ and $b\ge675$. For $C\ge10$, $B_w\le Cb$ and
$\log B_w\le\log10+\log b\le2\log b\le C\log b$.
\end{proof}

\paragraph{Leaves from the row properties.}
We use three facts about the theta maps of a point \cite{LeanSource2026}: (i) the complex realisations $i_C(t)$ of the types
satisfy $i_C(\mu_0)(\sigma_{v_0}(\iota_tz))=i_C(t)(\sigma_{v_0}(z))$, and $i_C(\mu_0)$ is injective
(\lref{GlobalTheory.T.iC_iota_mulVec_eq_iC}); (ii) the theta vector of the form attached to $y\in K^{n_V(t)}$ is
$\Theta_t(j_X(e_{V_0}(y)))=i_C(t)(\sigma_{v_0}(P_ty))$ (\lref{GlobalTheory.T.theta_j_ev0_eq_ic_pq}); (iii) these forms are
good rational points (\lref{GlobalTheory.T.eV0_mem_goodF'}). Moreover, by the square condition of the tie and the
admissibility of the choice, for every $\eta$ the constant $c(\eta,\mu_0)$ of the choice is a non-zero constant of the
matched square integrals at $\mu_0$ and $\sigma(r_t(\eta))c(\eta,\mu_0)=c(\eta,t)$; hence
$c_{\mathrm{Tot}}(\mu_0)\ne0$ and $\sigma(\prod_\eta r_t(\eta))=c_{\mathrm{Tot}}(t)/c_{\mathrm{Tot}}(\mu_0)$ for every type $t$,
with $c_{\mathrm{Tot}}$ as in $\sfname{FormulaK}$. (Every type is \emph{in use}, that is, equal to $\mu_0$ or to a value of
$\mu$, by the conditions on $\mu$ at the point.)

\begin{lemma}\label{pf:fit:leafrow}
Assume $\sfname{BundleCondition}$, $\sfname{SettingAtPoint}$, $\sfname{InjK}$, $\sfname{ConstK}$ and $\sfname{FormulaK}$.
Then the kernel leaf \ref{pf:fit:F:Ker} holds; if moreover $\sfname{LinK}$ holds, the unique map leaf \ref{pf:fit:F:C} holds.
\formallabel{GlobalTheory.T.leafKernel_tie_of_claim_inputs}
\formallabel{GlobalTheory.BigTop.Fun.leafKernel_of_linK_injK_constK_formula2K}
\formallabel{GlobalTheory.BigTop.Fun.ker_mulVecLin_eq_U_of_formulaAt}
\formallabel{GlobalTheory.T.oneMap_of_bundle_setAt_lin_inj_const_formula}
\formallabel{GlobalTheory.BigTop.Fun.oneMap_of_bundle_setAt_kernel_conditions}
\formallabel{GlobalTheory.BigTop.Fun.exists_unit_root_ze_leafC_of_square_nonzero}
\formallabel{GlobalTheory.complex_matrix_descends_of_rational_on_tuple}
\formallabel{GlobalTheory.BigTop.Fun.TieModel.ratio_eq}
\end{lemma}

\begin{proof}
By \cref{pf:fit:transfer}(b) the row properties hold at the choice $a$ of the tie. $\sfname{FormulaK}$ (the field condition
being part of $\mathrm{OK}(a)$) gives $\kappa\in\C$ and roots of unity $\zeta_t\in K$ with
$\sigma(\zeta_t)c_{\mathrm{Tot}}(t)\psi_K(b,\Theta_t(j_Xx))=\kappa v_K(t,x)_b$ for good rational $x$; by $\sfname{ConstK}$,
$c_{\mathrm{Tot}}(t)\ne0$. Combining with \eqref{eq:pf:fit:tieV}, (ii) and (iii), for all $t$ in use, $y$ and $b$:
\begin{equation}\label{eq:pf:fit:vpsi}
\sigma(\zeta_t)c_{\mathrm{Tot}}(t)\,\psi_K\bigl(b,i_C(t)(\sigma_{v_0}(P_ty))\bigr)=\kappa\,\sigma\bigl((V_ty)_b\bigr).
\end{equation}
\emph{$\kappa\ne0$.} By $\sfname{BundleCondition}$ the side is well formed, so $d_Q(\mu_0)>0$ and $P_{\mu_0}$ is surjective;
take $y$ with $P_{\mu_0}y$ a standard basis vector. Then $i_C(\mu_0)(\sigma(P_{\mu_0}y))\ne0$ by (i), so some
$\psi_K(b,\cdot)$ does not vanish on it by $\sfname{InjK}$, and \eqref{eq:pf:fit:vpsi} forces $\kappa\ne0$.

\emph{Kernel.} For $t=\mu_0$, \eqref{eq:pf:fit:vpsi}, $\kappa\ne0$, the injectivity of $\sigma$ and of $i_C(\mu_0)$ and
$\sfname{InjK}$ show: $V_{\mu_0}y=0$ if and only if $\psi_K(b,i_C(\mu_0)(\sigma(P_{\mu_0}y)))=0$ for all $b$, if and only if
$P_{\mu_0}y=0$. Thus $\ker V_{\mu_0}=U_{\mu_0}$.

\emph{Unique map.} By $\sfname{LinK}$ the pairings defining $\psi_K(b,\cdot)$ converge, so $\psi_K(b,\cdot)$ is additive and
compatible with the scalar action twisted by the base embedding; hence
$\psi_K(b,i_C(\mu_0)(\sigma(z)))=\sum_k\sigma(z_k)\psi_K(b,i_C(\mu_0)(e_k))$ for $z\in K^D$. Let $\Psi_{\C}$ be the complex
$n\times D$ matrix with entries $c_0\psi_K(b,i_C(\mu_0)(e_k))$, $c_0=\kappa^{-1}\sigma(\zeta_{\mu_0})c_{\mathrm{Tot}}(\mu_0)$, put
$\zeta'_t=\zeta_t/\zeta_{\mu_0}$ and $\rho_t=\zeta'_t\prod_\eta r_t(\eta)$. By (i), \eqref{eq:pf:fit:vpsi} and
$\sigma(\prod_\eta r_t(\eta))=c_{\mathrm{Tot}}(t)/c_{\mathrm{Tot}}(\mu_0)$,
\[
\Psi_{\C}\,\sigma(\iota_tP_t)=\sigma(\rho_t^{-1}V_t)\qquad\text{for all types }t.
\]
The map $\iota_{\mu_0}$ is bijective (a row condition of $S_{\mathrm f}$, by $\sfname{BundleCondition}$) and $P_{\mu_0}$ is
surjective, so $A=\iota_{\mu_0}P_{\mu_0}$ has a right inverse $A'$ over $K$. Put $\Psi=\rho_{\mu_0}^{-1}V_{\mu_0}A'\in
M_{n\times D}(K)$; then $\sigma(\Psi)=\Psi_{\C}\,\sigma(AA')=\Psi_{\C}$, and by the injectivity of $\sigma$ the display gives
$\Psi\iota_tP_t=\rho_t^{-1}V_t=(\prod_\eta r_t(\eta))^{-1}(\zeta'^{-1}_tV_t)$ for all $t$. The $\zeta'_t$ are roots of unity
with $\zeta'_{\mu_0}=1$, which is \ref{pf:fit:F:C}.
\end{proof}

\begin{lemma}\label{pf:fit:leafmm}
$\sfname{SettingAtPoint}$ implies the minimal mass leaf \ref{pf:fit:F:MM}.
\formallabel{GlobalTheory.T.leafMinimalMass_of_setAt}
\formallabel{GlobalTheory.BigTop.Fun.LeafMMTie.mm_tie_alt}
\formallabel{GlobalTheory.BigTop.Fun.LeafMMTie.leafMM_of_test}
\formallabel{GlobalTheory.FunObj.IsTestLat.gridTestLattice_largestSlope_bound}
\end{lemma}

\begin{proof}
By the test condition of the tie, $M$ is the dual of the sum of the lattices of the forms of the choice, indexed by pairs of
test indices and weight indices, and $X'$ is a metric of the inverse transpose of their Gram matrix. By
\cref{pf:fit:testlat}(c), in the dual form (b), it suffices that the spread of the weights of the grid is at most $R_p(C)\ge1$:
then $M$ is non-zero and finitely generated and $\fitmumax(M;X')\le B(f,n_p,R_p(C))=4fn_pR_p(C)\log(n_pR_p(C))=E_M(C)$, with
$f$ the number of embedding types and $n_p$ the rank. For every embedding type, the difference of two highest weights of the grid is at
most $(v_2-v_k)+1+(m+1)\mathrm{gap}+n_p$, where $v_2,v_k$ are entries of the Harish-Chandra parameter and $\mathrm{gap}$ is
the gap of the datum (\lref{GlobalTheory.BigTop.Fun.LeafMMTie.La_sub_le}). Since $v_2-v_k\le L_c$ (the length of the
cloud), $\mathrm{gap}\le\gamma(C)$ for $C\ge1$ and $1+(m+1)\gamma(C)\le(m+1)(\gamma(C)+1)$, the spread is at most $R_p(C)$;
also $R_p(C)\ge(m+1)(\gamma(C)+1)\ge1$. Hence \ref{pf:fit:L:MM}.
\end{proof}

\begin{lemma}\label{pf:fit:leafipb}
$\sfname{LambdaLowerBound}_2$ implies the leaf \ref{pf:fit:F:IPb}.
\formallabel{GlobalTheory.T.leafInfinitePlaceBound_of_lamLow}
\formallabel{GlobalTheory.BigTop.Fun.ipb_of_lamLow_alt}
\formallabel{GlobalTheory.InfinitePlacesProbe.clauseB_of_lifts}
\formallabel{GlobalTheory.BigTop.Fun.IPbProof.La_le_three}
\formallabel{Formalism.L2.DIV.Arith4S.Lp_le_La}
\formallabel{Formalism.L2.DIV.Arith4S.log_g_le}
\formallabel{Formalism.MainNotions.DIV.Proofs.HHParts.nV_le}
\end{lemma}

\begin{proof}
The two sides of a point share $\mu$, $\kappa$ and the scalars $\lambda$. $\sfname{LambdaLowerBound}_2$ gives, eventually,
$l_6(j)\le\log\lambda_{j,\mu_0,p,v}$ for all $j,p,v$ and every $C$, where for the numbers of a point
$l_6(j)=\sum_\rho\log\mathrm{hook}(j+1,\rho)-C_{\mathrm{sz}}fn_V^2L_p$ (the term $C\,l_{6u}$ vanishes). The hook sum is
non-negative and $n_V\le N$. Next, $L_p=\log(b\,\mathrm{rd}_Eg_1g_2n_V)$ and $L_a=\log x+\log b+\log\mathrm{rd}_E$,
where $1\le\mathrm{rd}_E$. The parameter choice gives $Y$ with $\log Y\le C_{YL}\log b$ and $g_i\le2Y$, so
$\log g_i\le\log2+C_{YL}\log b$; and $n_V=k+m-1+t$ is at most the size polynomial $1+f+k+m+t+N+\log b$, which is at most
$x^{c}$, $c=\max(C_t,5)+3$, for $x\ge9$ (\lref{Formalism.MainNotions.DIV.Proofs.szPol_le_rpow}: under the size
conditions this polynomial bound holds). As $L_a\ge\log9+\log b\ge2\log2$, we get
$L_p\le L_a+L_a+(2\log2+2C_{YL}\log b)+c\log x\le(3+2C_{YL}+c)L_a$, and
$L_a\le3\log b$ because the size rows give $x\le n\le b$ and $\mathrm{rd}_E\le b$ (both $n$ and $\mathrm{rd}_E$ are at most
a conductor quantity of the host which the size rows bound by $b$); thus
$L_p\le3(6+2C_{YL}+\max(C_t,5))\log b$. Hence, for
$C\ge\max(C_{\mathrm{sz}},0)(6+2C_{YL}+\max(C_t,5))$, $l_6(j)\ge-3CfN^2\log b$. Summing over the at most $m$ indices
$j\in J$, the constant function $g=-3CfmN^2\log b$ satisfies $g\le\sum_j\log\lambda_{j,\mu_0,p,v}$ and $\fitmi{g}=g$.
\end{proof}

\paragraph{Numerical leaves.}
\begin{lemma}\label{pf:fit:leafnum}
$\sfname{ScaleCondition}$ holds, and so do the leaves \ref{pf:fit:F:DesH}, \ref{pf:fit:F:Rank}, \ref{pf:fit:F:Ram} and
\ref{pf:fit:F:Pol}.
\formallabel{GlobalTheory.T.scaling_function_bound'}
\formallabel{GlobalTheory.BigTop.Fun.existsIntegralMultiple_holds_alt}
\formallabel{GlobalTheory.BigTop.Fun.Close.desHLeaf_alt}
\formallabel{GlobalTheory.T.rank_gap_bound}
\formallabel{GlobalTheory.BigTop.Fun.rankGapBound_alt}
\formallabel{GlobalTheory.T.ramified_sum_bound'}
\formallabel{GlobalTheory.BigTop.Fun.ramified_sums_seventh_bound_alt}
\formallabel{GlobalTheory.BigTop.Fun.RamSum.leaves_of_sum}
\formallabel{GlobalTheory.BigTop.Fun.RamSum.ramifiedPlaces_sum_log_absNorm_le}
\formallabel{GlobalTheory.BigTop.Fun.RamSum.le_PPF}
\formallabel{GlobalTheory.T.polynomial_descent_estimate'}
\formallabel{GlobalTheory.BigTop.Fun.desPol_alt}
\formallabel{GlobalTheory.BigTop.Fun.PolBound.pol_core}
\formallabel{GlobalTheory.BigTop.Fun.PolBound.factors_le}
\end{lemma}

\begin{proof}
These are elementary consequences of the definitions $x>x_1>1$, $b=e^{x^\eps}$, of the rank windows
\eqref{eq:families:rankval} and of the size rows; we indicate the inputs.

\emph{Scale.} $\log b=x^\eps\ge1$, and $\log x\le2x^\eps$ for $x$ larger than a constant depending on $\eps$.

\emph{Descent height.} $2\le m$, $m+3\le k$ and $f\le k$ are conditions on the ranks; $m^4\le f=2n$ follows from
$m\le2(x\delta_1)^{1/2}$, $n\ge x^{2+8\eps}$ and $x^{8\eps}\ge8$. $L_c\le b$: the entries of the Harish-Chandra parameters lie in
an interval of length at most the size constant $S_c$ of the parameter choice, and $S_c\le b$.

\emph{Rank gap.} With $t\ge x^{377}$ (from the parameter choice) and $t\le x^{\lceil|C_t|\rceil+1}$, the windows give the
polynomial bounds $mfN^2\le b$, $80r^4\le N$, $n_p\le2mr$, $R_p(C)\le CmN^3b$ and $\log(n_pR_p(C))\le4\log b$ for $x$ larger
than a constant depending on $\eps$, $C$ and $C_t$ (\lref{GlobalTheory.SizeLeafCheck.tree_leaf_parameter_arithmetic_of_bounds} and \lref{GlobalTheory.BigTop.Fun.design_lines_of_ranks}).

\emph{Ramified sums.} The base place of the grid lies over the first modulus $q_0$ and the other places of $S_0$ over primes
dividing $g_1g_2p_dp'_d$; as $\sum_{v\mid n}\log\Nm v\le[F:\Q]\log n=f\log n$ and $q_0,p_d,p'_d<X'\le b$, $g_1,g_2\le b^9$, one
gets $\sum_{q\in S_0}\log\Nm q\le21f\log b$. The moduli $q(q)=\Nm q$ are prime powers. For $C\ge42$ one has
$42f\log b\le\Pi(C,b)$, and since $e(q_*)\log q(q_*)\le\Pi+\log q(q_*)$, $e(q)=2$ otherwise and $c_0\le2$, this gives
$\Sigma_0\le\Pi+42f\log b\le2\Pi$ and $c_S\le42f\log b\le45f\log b$.

\emph{Polynomial.} With $T=\lceil|C_t|\rceil+1$, the size rows give $f,m,k,t,\log b\le x^{T+9}$ and
$\log(bn_rg_1g_2p_dp'_d|d_E|)\le20x^{T+9}$; an elementary computation gives $\fitpol(C,b)\le1.3\cdot10^{10}C^4x^{16(T+9)}$,
which is at most $x^{16(T+9)+1}\le x^{C_H}$ for $x\ge1.3\cdot10^{10}C^4$, because the largeness of the family in the
sense of \cref{def:families:closed} includes $C_H\ge16(T+9)+1$.
\end{proof}

\paragraph{Existence of fitted functions and the bundle condition.}
\begin{lemma}\label{pf:fit:leafex}
\begin{enumerate}[label=(\alph*)]
\item $\sfname{SettingAtPoint}$, $\sfname{LogDer}$, $\sfname{SumOKK}$, $\sfname{SumRatK}$, $\sfname{ConstK}$, the statement of
\cref{pf:fit:toric}(b), the integrality statement $\sfname{IntSt}$ for CM theta norms, and $\fitGBig$ of $\sfname{FieldDat}$,
$\sfname{TorEx}$, $\sfname{DenEx}$, $\sfname{BotEx}$ and $\sfname{BotSmall}$ imply $\sfname{ExistsFunctionIn}(\sfname{Fld},\mathcal T)$.
\item $\sfname{BundleCondition}$ holds.
\end{enumerate}
\formallabel{GlobalTheory.T.existsFunctionIn_of_sqRatio_range}
\formallabel{GlobalTheory.BigTop.Fun.ex_of_claims_alt}
\formallabel{GlobalTheory.BigTop.Fun.bundle_of_rows_alt}
\formallabel{Formalism.IndefR.Links.primRows3_data}
\end{lemma}

\begin{proof}
(a) Let $G,W$ be over $R$, $G$ fitting $p$, and $(n,a)$ with $\mathrm{OK}(a)$ (the field $\sfname{Fld}$). At the point the row
properties hold for $\mathrm{ch}(a)$, the square ratios lie in $\sigma(K)$, and, with the standing conditions as in
\cref{pf:fit:transfer}, the statements $\sfname{TorEx}$, $\sfname{DenEx}$, $\sfname{BotEx}$, $\sfname{BotSmall}$ and
$\sfname{FieldDat}$ hold at $G$; $\sfname{IntSt}$ gives an integrality exponent $d_l$ for the good level of $G$ (using
$\log\mathrm{rd}_E\le\log b$ from the size rows). For every $(\eta,s)$ with $1\le s\le h(\eta)$, $\sfname{TorEx}$ gives a
non-zero toric value $x_0$, and for $c\in\{0,1\}$ the toric values of the twisted characters are $\sigma(z)x_0$ with $z\in K$
(by the condition on square values contained in $\mathrm{OK}(a)$); the valuation shift $g_T$ is defined from the minimal
polynomial of the period-normalised value of $x_0$, which is algebraic by $\sfname{TorEx}$; $\sfname{DenEx}$ gives denominator
data $N_0$ with finite support. With these data the function datum $\phi$
with $G_\phi=G$ and $\sfname{TieAt}(\phi,a)$ is constructed by \lref{GlobalTheory.BigTop.Fun.TieModel.exists_tie} (from the choice:
$n$ the number of test indices, $X'$ and $M$ the test metric and lattice, $V_t$ the rational matrices of the class sums given by
$\sfname{SumOKK}$ and $\sfname{SumRatK}$, $r_t$ the square ratios, $z_A,z'_A$ the norms of bottom pairs, $z_T=1$, and the data
above).
(b) At a point with $\fitPC(p)$, the primitive links of the point give the conditions on the complex data at the base level
(\lref{Formalism.IndefR.Links.complex_data_conditions_base_level_of_primitive_links}), which contain the well-formedness of the
definite side, and the primitive rows give a datum with rows whose underlying indefinite datum is $S_{\mathrm f}$
(\lref{Formalism.IndefR.Links.primRows3}; that its underlying datum is $S_{\mathrm f}$ holds by construction).
\end{proof}

% ---------------------------------------------------------------------------
\subsection{Proof of the main theorem}\label{sec:pf:fit:main}

The functional equation enters in absolute value form: $\sfname{HeckeFunctionalEquationAbs}$ asserts that for a totally
complex field $E$ with a CM type, a unitary Hecke character $\chi$ with non-trivial infinity type and conductor exponents
$a$, the completed $L$-functions $\Lambda$ of $\chi$ and $\Lambda'$ of $\chi^{-1}$ satisfy
$\Lambda(s)=\epsilon\,(|d_E|\Nm\mathfrak f)^{1/2-s}\Lambda'(1-s)$ with $|\epsilon|=1$, $\mathfrak f$ the conductor. It follows from
$\sfname{HeckeTateFunctionalEquation}$ by \lref{GlobalTheory.Published.HeckeTateFunctionalEquation.abs} (the archimedean
and finite root numbers have absolute value $1$, up to the factor $(|d_E|\Nm\mathfrak f)^{1/2-s}$).

\begin{theorem}\label{pf:fit:close}
$F$, $\sfname{HeckeFunctionalEquationAbs}$, the five components of $\sfname{JoinOut}$, $\sfname{ThetaIntegralsAgree}$,
$\sfname{LocalMatching}$, $\sfname{FunctionalsSeparate}$ and $\sfname{LambdaLowerBound}_2$ imply
$\sfname{ConstantFunctionFits}(\sfname{Fld})$.
\formallabel{GlobalTheory.BigTop.Fun.Close.fittedFunction_at_zero_close_alt}
\end{theorem}

\begin{proof}
We collect the inputs.
\begin{enumerate}[label=(\arabic*),leftmargin=*]
\item By \cref{pf:fit:sqtop} (Howe duality at the real place being a component of $F$), $\fitGBig$ holds for the twelve
square statements \ref{pf:fit:S:DenSum}--\ref{pf:fit:S:Tower}.
\item By \cref{pf:fit:sumsize}, $\fitBigC$ holds for $\sfname{FinAt}$, $\sfname{AwayAt}$ and $\sfname{InfAt}$.
\item By \cref{pf:fit:rows}, the row properties $\sfname{SumOKK}$, $\sfname{ConstK}$, $\sfname{LinK}$, $\sfname{FormulaK}$,
$\sfname{SumRatK}$ and $\sfname{InjK}$ hold.
\item By \cref{pf:fit:setting}, $\sfname{SizeAt}$, $\sfname{SettingAtPoint}$ and $\sfname{LogDer}$ hold.
\item The CM theta norm hypotheses hold (\lref{GlobalTheory.T.cmThetaNormBoundHypotheses_satisfied}), hence the
integrality statement $\sfname{IntSt}$ (\lref{GlobalTheory.T.classSumDenominatorBound_of_cmThetaHypotheses}) and the
archimedean statement $\sfname{ArchSt}$ (\lref{GlobalTheory.ArchTop.archSt_of_inp}).
\item By \cref{pf:fit:toric}, the toric estimate holds at tied function data, and the square ratios lie in $\sigma(K)$.
\end{enumerate}
We apply \cref{pf:fit:leaves} with $\mathrm{Fld}=\sfname{Fld}$ and the tie $\mathcal T$ of \cref{sec:pf:fit:assembly}:
$\sfname{BundleCondition}$ and $\sfname{ExistsFunctionIn}(\sfname{Fld},\mathcal T)$ hold by \cref{pf:fit:leafex} (with (1), (3), (4),
(5), (6)); $\sfname{ScaleCondition}$ and the leaves \ref{pf:fit:F:DesH}, \ref{pf:fit:F:Rank}, \ref{pf:fit:F:Ram},
\ref{pf:fit:F:Pol} by \cref{pf:fit:leafnum}; \ref{pf:fit:F:MM} by \cref{pf:fit:leafmm}; \ref{pf:fit:F:C} and
\ref{pf:fit:F:Ker} by \cref{pf:fit:leafrow}; \ref{pf:fit:F:FP}, \ref{pf:fit:F:IPa} and \ref{pf:fit:F:Wt} by
\cref{pf:fit:leafsize}; \ref{pf:fit:F:IPb} by \cref{pf:fit:leafipb}; and \ref{pf:fit:F:Sq}, \ref{pf:fit:F:Den},
\ref{pf:fit:F:Int}, \ref{pf:fit:F:Arch}, \ref{pf:fit:F:Thb} by \cref{pf:fit:leafsq}.
\end{proof}

\begin{proof}[Proof of \cref{pf:fit:main}]
The statement $\sfname{JoinOut}$ is the conjunction of the five statements $\sfname{ToricDenominator}$,
$\sfname{Transport}$, $\sfname{TwistedTransport}$, $\sfname{ZetaDenominators}$ and $\sfname{Refl}$, and the functional
equation $\sfname{HeckeTateFunctionalEquation}$ implies its absolute value form. Apply \cref{pf:fit:close}.
\end{proof}

% ---------------------------------------------------------------------------
\subsection{Norms of bottom pairs}\label{sec:pf:fit:normbig}

The remaining subsections prove the statements on square terms, size bounds and test lattices which were used above.
We begin with $\sfname{NormBig}$.

\paragraph{Pairs and their norms.}
Let $D$ be a global datum, with CM field $E$ and totally real subfield $F$. For $N\ge1$ let $G_N$ be the group of finite
adelic points of the definite unitary group of rank $N$ of the host, $K_N\subseteq G_N$ its integral subgroup and, for an
open subgroup $J\subseteq K_N$, $\mathrm{Cls}(N,J)$ the finite set of double cosets of $G_N$ modulo the rational points and
$J$; $\Gamma_h$ is the (finite) stabiliser of the class of $h$, and
$\mathrm{mass}(N)=\sum_{x\in\mathrm{Cls}(N,K_N)}|\Gamma_x|^{-1}$. A \emph{pair} $P$ with parameters $(\delta,t_0)$ consists
of ranks $a\le d$ with $d$ even; a CM curve with parameter $\delta$, whose lattice has area $A_\delta>0$, with a two-torsion
point $t_0$; a Weil representation $\Omega$ of $G_a\times G_d$ with its vacuum datum; weights $n(\tau)\in\N$ for the embedding
types $\tau$; a vector $\Phi_P$ fixed by $K_a$; and a subgroup $J\subseteq K_d$ of finite index on which $\Phi_P$ is an
eigenvector with eigenvalues of absolute value $1$. For $h\in G_d$ and $T$ in the finite non-empty set $\mathrm{Tup}(a,d,n)$
of weight tuples, the \emph{theta coordinate} of $P$ is
$\mathrm{coord}_P(h,T)=\sum_{x\in\mathrm{Cls}(a,K_a)}|\Gamma_x|^{-1}\,(\partial_T\,\Omega(x,h)\Phi_P)(0)$, where $\partial_T$
is the differential operator of the minor polynomial $p_T$ of $T$, and the \emph{norm} of $P$ is
\begin{equation}\label{eq:pf:fit:pairnorm}
\|P\|=\sum_{h\in\mathrm{Cls}(d,J)}|\Gamma_h|^{-1}\sum_{T}|\mathrm{coord}_P(h,T)|^2 .
\end{equation}
Let $G$ be a grid, $t$ a type and $(i,j)$ an index pair, and let $b_+>b_-\ge2$ be the bottom ranks of $(i,j)$; $b_++b_-$ is
even (\lref{GlobalTheory.FunObj.SqNum.Model.bot_ranks}). Two pairs $(P,P')$ form a \emph{bottom pair} at $t$ if $P$ has
ranks $(b_+,b_++b_-)$ and $P'$ has ranks $(b_-,b_++b_-)$; $P$ carries the characters $\beta_H(G,i,j)^{-1}$ and
$\nu_H(G,t,i,j)^{-1}$, and $P'$ the characters $\alpha_H(t,i,j)$ and $\nu_H(G,t,i,j)^{-1}$; $P$ and $P'$ are standard data for
the modulus $\mathfrak m_t$ of the grid at $t$ (the conductor ideal of the grid times the primes of the twisted places of
$t$) and for its conjugate, respectively; and their levels are block levels for $\mathfrak m_t$. Put $e_0=(b_+-b_-)/2$ and
$h_0=(b_++b_-)/2$. The statement $\sfname{NormBig}$ is $\fitGBig$ of: for all $\delta,t_0$ and $(i,j)$ there is $A\ne0$ such
that $\|P\|=A\prod_{s=1-e_0}^{h_0}\Lambda(s)$ for every type $t$, every bottom pair $(P,P')$ at $t$ and every completed
$L$-function $\Lambda$ of $\nu_H(G,t,i,j)$.

\paragraph{Constants.}
For ranks $a<d$ and weights $n$ let $N(a,n)=\prod_\tau(\pi/A_\delta)^{a\,n(\tau)}>0$; let $\nu(a,d)$ be the norm of the
vacuum vector, a power of the squared norm of the theta-null value of the CM lattice at $t_0$, which is non-zero
(\lref{GlobalTheory.ConjJet.thetaZeroNormSq_ne_zero_cmLattice}); let $\Phi(a,d,n)=\sum_T\langle p_T,p_T\rangle$, with the
Fock inner product $\langle p,q\rangle=\sum_\alpha\alpha!\,c_\alpha(p)\overline{c_\alpha(q)}$ of polynomials; for a set $S$ of
finite places of $F$ let $d_S(a,d;s)=\prod_{v\notin S}d_v(a,d;s)$ be the Euler product of the local normalising factors of
the doubling integral, which does not vanish for $\Re s>0$ (\lref{GlobalTheory.LFun.dSfun_ne_zero_of_re_pos}); and put
\[
c(a,d,n)=N(a,n)\,\mathrm{mass}(d)\,\mathrm{mass}(a)\,\Phi(a,d,n)\,\nu(a,d)\big/d_\emptyset\bigl(a,d;\tfrac{d-a}2\bigr).
\]

\begin{lemma}[Petersson norms of bottom pairs]\label{pf:fit:pet}
Assume $F$ and $\sfname{Inp}(D)$, and let the size data satisfy the standing conditions. Let $(P,P')$ be a bottom pair at the
type $t$ for $(i,j)$, let $(a,d)$ and $n$ be the ranks and weights of $P$, and let $\Lambda$ be a completed $L$-function of
$\nu_H(G,t,i,j)$. Then
\[
\|P\|=c(a,d,n)\prod_{s=1-e_0}^{h_0}\Lambda(s).
\]
\formallabel{GlobalTheory.NormRatio.petersson_norm_bottom_theta_lift_eq}
\end{lemma}

\begin{proof}
Here $a=b_+$ and $d=b_++b_-$, so $1\le a<d$, $d=2h_0$ and $a=h_0+e_0$. Let $S_t$ be the set of places of the grid at $t$.
Since the auxiliary primes of the host are prime to the moduli $\mathrm{QS}\cdot q$ (a standing condition), $\mathfrak m_t$ is a
\emph{good modulus} for $S_t$: it is non-zero, prime to its complex conjugate and to $2q$, and its prime divisors are the places
of $S_t$ (\lref{GlobalTheory.NormRatio.goodMod_modT}). The level $J$ of $P$ is a block level for $\mathfrak m_t$, hence a
congruence level (it contains the principal congruence subgroup of level $\Nm\mathfrak m_t$;
\lref{GlobalTheory.NormExact.isCongLevel_of_isBlockLevel}).

\emph{The dictionary.} By \lref{GlobalTheory.NormExact.exact_dictionary} (for a pair with congruence level and under
$\sfname{Inp}(D)$), there are an \emph{exact datum} $X$ for $P$ --- a realisation of $P$ by a Weil representation of adelic
unitary groups $\mathbf G_A,\mathbf G_B$ of ranks $a,d$, with a Hecke character $\chi_B$ and theta lifts $\theta_T$ of the
Fock vectors of the tuples $T$ --- and a maximal level $K_0$ of $\mathbf G_B$ with Haar measure $\mu_0$; and for every such
choice, $\theta_T(\iota_Bh)=c_1\,\mathrm{coord}_P(h,T)$ with $|c_1|^2N(a,n)=1$, and
\begin{equation}\label{eq:pf:fit:dict}
\|P\|=N(a,n)\,[K_d:J]\sum_T\langle\theta_T,\theta_T\rangle_{\mu_0},
\end{equation}
where $\langle\cdot,\cdot\rangle_{\mu_0}$ is the Petersson inner product on $\mathbf G_B$.

\emph{The character.} The finite part of $\chi_B$ is $c^{-e}$, where $c$ is the character of the Weil datum of $P$ and $e$
its exponent, and $c^{e}=\nu_H(G,t,i,j)^{-1}$ on finite ideals by the character condition of a bottom pair; hence $\chi_B$
corresponds to $\nu_H(G,t,i,j)$
(\lref{GlobalTheory.T.heckeCharEq_of_finitePart_power_inverse}). As $F$ is the maximal real subfield of $E$, $\chi_B$ is
unramified outside $S_t$ and, for every $s\in\C$, $\Lambda(s)$ is the value at $s$ of the holomorphic continuation of the
partial $L$-function $L^{S_t}(s,\chi_B^{-1})$ (\lref{GlobalTheory.T.squareNu_unramified_off_pair_partial_L_primitive}).

\emph{The Rallis inner product formula.} By \lref{GlobalTheory.NormRatio.theta_lift_petersson_norm_eq_rallis_integral}
(for an exact datum with maximal level, a good modulus $\mathfrak m$ with a standard datum, $1\le a<d$ and $\chi_B$ unramified
outside $S$, the doubling integral identity for anisotropic unitary groups \cite{Rallis1984}, with the volumes of the
definite quotients expressed as masses), for every $T$
\[
\langle\theta_T,\theta_T\rangle_{\mu_0}=\mathrm{mass}(d)\,\mathrm{mass}(a)\,
\frac{\prod_{k=0}^{a-1}\Lambda\bigl(\tfrac{d-a}2+\tfrac a2-k\bigr)}{d_{S_t}\bigl(a,d;\tfrac{d-a}2\bigr)}\,
\langle p_T,p_T\rangle\,\ell(\mathfrak m_t,a,d-a)\,\nu(a,d),
\]
where $\ell(\mathfrak m,a,a')$ is the local norm of the standard datum at the places dividing $\mathfrak m$. Since
$\tfrac{d-a}2+\tfrac a2-k=h_0-k$ and $a=h_0+e_0$, the numerator is $\prod_{s=1-e_0}^{h_0}\Lambda(s)$ (reindex $s=h_0-k$).

\emph{Local factors.} The relative index $[K_d:J]$ of a block level equals the local index
$\iota(\mathfrak m_t,a,d-a)=\prod_{w\mid\mathfrak m_t}\iota_w$, a product of explicit local indices of block subgroups
(\lref{GlobalTheory.StdLoc.relIndex_blockLevel}), and the local index, the local norm and the normalising factors satisfy
$\iota(\mathfrak m_t,a,a')\,\ell(\mathfrak m_t,a,a')\,d_\emptyset(a,a+a';\tfrac{a'}2)=d_{S_t}(a,a+a';\tfrac{a'}2)$ for $a'\ge1$
(\lref{GlobalTheory.T.localIndex_localNorm_normalizingFactor_eq_partial}: the Euler factors at the places of $S_t$ are the
products of the local index and the local norm). Substituting into \eqref{eq:pf:fit:dict} and summing over $T$,
\[
\|P\|=N(a,n)\,\iota\,\mathrm{mass}(d)\,\mathrm{mass}(a)\,\frac{\prod_s\Lambda(s)}{d_{S_t}}\,\Phi(a,d,n)\,\ell\,\nu(a,d)
=c(a,d,n)\prod_{s=1-e_0}^{h_0}\Lambda(s),
\]
because $\iota\ell/d_{S_t}=1/d_\emptyset$ at $s=\tfrac{d-a}2$, where both normalising factors are non-zero since $\Re\tfrac{d-a}2>0$.
\end{proof}

\begin{lemma}\label{pf:fit:petratio}
Assume $F$ and $\sfname{Inp}(D)$ for all $D$.
\begin{enumerate}[label=(\alph*)]
\item $\fitGBig$ holds for the statement: for bottom pairs $(P,P')$ at $t$ and $(P_1,P_1')$ at $t'$ (for the same $\delta,t_0$
and $(i,j)$) and completed $L$-functions $\Lambda$ of $\nu_H(G,t,i,j)$ and $\Lambda'$ of $\nu_H(G,t',i,j)$,
$\|P\|\prod_{s=1-e_0}^{h_0}\Lambda'(s)=\|P_1\|\prod_{s=1-e_0}^{h_0}\Lambda(s)$.
\item If moreover $\sfname{LNe}$ holds, then $\fitGBig$ holds for the statement: $\|P\|\ne0$ for every bottom pair $(P,P')$.
\end{enumerate}
\formallabel{GlobalTheory.T.theta_lift_petersson_norm_ratio_formula}
\formallabel{GlobalTheory.NormRatio.big_theta_petersson_norm_ratio_eq}
\formallabel{GlobalTheory.T.big_theta_lift_ne_zero_of_petersson_norm}
\formallabel{GlobalTheory.NormBigNe.big_theta_petersson_norm_ne_zero}
\formallabel{GlobalTheory.NormRatio.fockSum_ne_zero}
\formallabel{GlobalTheory.NormRatio.fockIp_self_re_pos}
\formallabel{GlobalTheory.NormRatio.std_mass_ne_zero}
\formallabel{GlobalTheory.NormExact.normConst_pos}
\formallabel{GlobalTheory.VacuumNorm.vacNorm_ne_zero}
\formallabel{GlobalTheory.NormBigProof.squareBottomHalfRankSum_le_halfRank}
\end{lemma}

\begin{proof}
(a) holds for all data satisfying the standing conditions. Bottom pairs at two types have the same ranks and the same weights
(\lref{GlobalTheory.BottomReal.bottom_real_eq}: the ranks are fixed by $(i,j)$, and the weights agree because the
characters $\nu_H(G,t,i,j)$ and $\nu_H(G,t',i,j)$ agree on the finite ideles of the unit group); so the constants $c(a,d,n)$
of $P$ and $P_1$ coincide, and the identity follows from two applications of \cref{pf:fit:pet}.

(b) Let the thresholds be those of $\sfname{LNe}$. For a bottom pair at $t$, $\sfname{LNe}$ gives a completed $L$-function
$\Lambda$ of $\nu_H(G,t,i,j)$ with $\Lambda(s)\ne0$ for integers $1-e_0\le s\le\mathrm{hr}(i,j)$, where
$\mathrm{hr}(i,j)\ge h_0$ is the half rank (the bottom ranks are obtained from the ranks of $(i+1,j+1)$ and $(i,j)$ by subtracting
non-negative even numbers). By \cref{pf:fit:pet} it suffices that $c(a,d,n)\ne0$. Here $N(a,n)>0$ and $\nu(a,d)\ne0$;
$\mathrm{mass}(N)>0$, being a sum of the positive numbers $|\Gamma_x|^{-1}$ over the finite non-empty set $\mathrm{Cls}(N,K_N)$
(\lref{GlobalTheory.NormSum.cls_finite_card_le} and \lref{GlobalTheory.StabOrder.gam_ne_zero_dvd_invariant}: the class set is
finite and the stabilisers are finite); $\Re\Phi(a,d,n)>0$ because $\langle p,p\rangle=\sum_\alpha\alpha!|c_\alpha(p)|^2>0$ for
the non-zero minor polynomials and $\mathrm{Tup}(a,d,n)\ne\emptyset$ as $a\le d$; and $d_\emptyset(a,d;\tfrac{d-a}2)\ne0$.
\end{proof}

\begin{theorem}\label{pf:fit:normbig}
$F$, $\sfname{Inp}(D)$ for all $D$, and $\sfname{LNe}$ imply $\sfname{NormBig}$.
\formallabel{GlobalTheory.T.normBig_of_lNe}
\formallabel{GlobalTheory.NormBigProof.normBig_alt}
\end{theorem}

\begin{proof}
Take the maximum of the thresholds of $\sfname{LNe}$ and of the two statements of \cref{pf:fit:petratio}. Fix
$\delta,t_0,(i,j)$. If there is no bottom pair at any type, $A=1$ works. Otherwise fix a bottom pair $(P_0,P_0')$ at some type
$t_*$ and, by $\sfname{LNe}$, a completed $L$-function $\Lambda_0$ of $\nu_H(G,t_*,i,j)$ with $\Lambda_0(s)\ne0$ for
$1-e_0\le s\le h_0$. Put $A=\|P_0\|/\prod_s\Lambda_0(s)$; then $A\ne0$ by \cref{pf:fit:petratio}(b), and for every bottom pair
$(P,P')$ at a type $t$ and completed $L$-function $\Lambda$ of $\nu_H(G,t,i,j)$, \cref{pf:fit:petratio}(a) gives
$\|P\|\prod_s\Lambda_0(s)=\|P_0\|\prod_s\Lambda(s)$, i.e.\ $\|P\|=A\prod_s\Lambda(s)$.
\end{proof}

% ---------------------------------------------------------------------------
\subsection{Bottom constants}\label{sec:pf:fit:botconst}

We prove $\sfname{BotConst}$. For a grid $G$ let $\mathrm{lev}_Q(G)=\prod_{v\in S_0}\Nm v^{c(v)}$ be its level number and, for
a type $t$, $Q_n(G,t)=\prod_w\Nm w$ the product of the norms of the twisted places of $t$; at the base type there are no twisted
places, so $Q_n(G,\mu_0)=1$ and $\mathfrak m_{\mu_0}$ is the conductor ideal of the grid
(\lref{GlobalTheory.SqTop.Model.base_tuple}). A pair is \emph{small} (for $N$ and $B$) if $d\le2N$ and $n(\tau)\le2B$ for
all $\tau$; under $\sfname{BotSmall}$ all bottom pairs are small for $N=\mathrm{bigN}(D)$ and $B$ the maximal weight
$\max\mathrm{Wt}(D)$.

\paragraph{Integrality exponents.} For a prime $p$ and $r\ge0$ call $x\in\C$ \emph{$p$-bounded by $r$} if $x$ is algebraic and
$|\lambda(x)|_p\le r$ for every embedding $\lambda\colon\Qbar\to\C_p$; $p$-bounded numbers are stable under sums (with the
maximum of the bounds) and products (with the product of the bounds), integers are $p$-bounded by $1$, and complex conjugation
permutes the embeddings. \emph{Exponents} $(S,e)$ with bound $\beta$ are a finite set $S$ of primes and a function $e\ge0$
supported on $S$ with $\sum_{l\in S}e(l)\log l\le\beta$. They are \emph{coordinate-integral} for $(N,B,Q)$ if for every prime
$p$, every $Q'$ with $v_p(Q')\le v_p(Q)$ and every small pair $P$ at level $Q'$: every class of $\mathrm{Cls}(d,J)$ has a
representative integral at $p$; $|\Gamma_h|\ne0$ and $v_p(|\Gamma_h|)\le e(p)$ for all $h$; and $\mathrm{coord}_P(h,T)$ is
$p$-bounded by $p^{e(p)}$ for every $h$ integral at $p$ and every $T$. An \emph{integrality exponent} of the grid is a function
$d_l=\mathrm{wt}_{S,e}$, $\mathrm{wt}_{S,e}(w)=\sum_{l\in S}e(l)v_w(l)$, for exponents which are coordinate-integral for
$(\mathrm{bigN}(D),\max\mathrm{Wt}(D),\mathrm{lev}_Q(G))$, with an explicit bound $\beta$, and satisfy a similar integrality
condition for norms.

\begin{lemma}[Valuations of proportionality constants]\label{pf:fit:lemmaC}
Assume $\sfname{Inp}(D)$, let $(S,e)$ be exponents which are coordinate-integral for $(N,B,Q)$, and let $w$ be a finite place
of $K$ over the prime $p$. Let $P_1,P_1'$ be pairs at level $Q_1$ and $P_2,P_2'$ pairs at level $Q_2$, all small, with
$v_p(Q_1),v_p(Q_2)\le v_p(Q)$, $d(P_1')=d(P_1)$ and $d(P_2')=d(P_2)$. Let $b_1,b_2,v\in\C$ and integer matrices $M_1,M_2$
satisfy, for all $h$, $T$, $T'$,
\begin{gather*}
\overline{\mathrm{coord}_{P_1}(h,T)}=b_1\sum_{T'}M_1(T,T')\,\mathrm{coord}_{P_1'}(h,T'),\\
\mathrm{coord}_{P_2'}(h,T')\,b_2=\sum_TM_2(T',T)\,\overline{\mathrm{coord}_{P_2}(h,T)},
\end{gather*}
with $b_2\ne0$, $v$ algebraic and $|\lambda(v)|_p=1$ for all $\lambda$. If $\gamma,z,z'\in K$ satisfy $b_1=\sigma(\gamma)b_2v$,
$\sigma(z)=\|P_1\|$, $\sigma(z')=\|P_2'\|$ and $z,z'\ne0$, then $v_w(\gamma)\le v_w(z)+v_w(z')+6\,\mathrm{wt}_{S,e}(w)$.
\formallabel{GlobalTheory.T.valuation_proportionalityRatio_le_bottomPairs}
\formallabel{GlobalTheory.NormSum.LemmaC.vw_le_of_bottom_pairs}
\end{lemma}

\begin{proof}
Let $r=p^{e(p)}$. If $b_1=0$, all coordinates of $P_1$ vanish and $\|P_1\|=0$ by \eqref{eq:pf:fit:pairnorm}, contradicting
$z\ne0$; hence $b_1\ne0$, and $\gamma\ne0$, $v\ne0$. The norm \eqref{eq:pf:fit:pairnorm} does not depend on the representatives
of the classes (\lref{GlobalTheory.T.definiteThetaPair_terms_norms_independent_representatives}), so we may compute
$\|P_1\|$ with representatives $h_q$ integral at $p$. By the first relation,
$|\mathrm{coord}_{P_1}(h,T)|^2b_1^{-1}=\mathrm{coord}_{P_1}(h,T)\sum_{T'}M_1(T,T')\mathrm{coord}_{P_1'}(h,T')$, hence
\[
\|P_1\|\,b_1^{-1}=\sum_q|\Gamma_{h_q}|^{-1}\sum_T\mathrm{coord}_{P_1}(h_q,T)\sum_{T'}M_1(T,T')\,\mathrm{coord}_{P_1'}(h_q,T'),
\]
and every summand is $p$-bounded by $r\cdot r\cdot1\cdot r$ (coordinate integrality; $h_q$ is integral at $p$ for $P_1'$ as
well). So $\|P_1\|b_1^{-1}$ is $p$-bounded by $r^3$, and in the same way, using the second relation,
$\|P_2'\|\,b_2=\sum_q|\Gamma_{h_q}|^{-1}\sum_{T'}\overline{\mathrm{coord}_{P_2'}(h_q,T')}\sum_TM_2(T',T)
\overline{\mathrm{coord}_{P_2}(h_q,T)}$ is $p$-bounded by $r^3$. Put $x=zz'\gamma^{-1}\in K^\times$. Then
$\sigma(x)=\|P_1\|\|P_2'\|\sigma(\gamma)^{-1}=(\|P_1\|b_1^{-1})(\|P_2'\|b_2)\,v$ is $p$-bounded by $r^6=p^{6e(p)}$. An element
$x\in K^\times$ all of whose $p$-adic images have absolute value at most $p^c$ satisfies $v_w(x)\ge-c\,v_w(p)$ at the places
$w\mid p$ (\lref{GlobalTheory.PadicPlaceBound.neg_mul_vw_le_vw}); so $v_w(z)+v_w(z')-v_w(\gamma)=v_w(x)\ge-6e(p)v_w(p)$.
Finally $\mathrm{wt}_{S,e}(w)=e(p)v_w(p)$, because $v_w(l)=0$ for primes $l\ne p$ and $e(p)=0$ if $p\notin S$.
\end{proof}

The statement $\sfname{BotConst}$ is $\fitGBig$ of the following. Let $W$ be a Weil datum of the grid with the standard
additive character, $S$ a grid choice whose numbers lie in $\sigma(K)$ at $\mu_0$, $(i,j)$ an index pair, $t$ a type, $X$ a
bottom triple for $(i,j)$ with constants $c_0(\cdot)$ (the constants of the factorisation of its triple integrals at the types),
$(P_0,P_0')$ a bottom pair at $\mu_0$ at level $\mathrm{lev}_Q(G)$ and $(P,P')$ a bottom pair at $t$ at level
$\mathrm{lev}_Q(G)Q_n(G,t)$ (with the same $\delta,t_0$), $z,z'\in K^\times$ with $\sigma(z)=\|P\|$ and $\sigma(z')=\|P_0'\|$,
and $d_l$ an integrality exponent of $G$. Then there is $\gamma\in K^\times$ with $c_0(\mu_0)\ne0$,
$\sigma(\gamma)c_0(\mu_0)=c_0(t)$ and $v_w(\gamma)\le v_w(z)+v_w(z')+6d_l(w)$ for every finite place $w$ of $K$ with
$Q_n(G,t)\notin w$.

\begin{theorem}\label{pf:fit:botconst}
$F$, $\sfname{ThetaIntegralsAgree}$, $\sfname{LocalMatching}$, $\sfname{Inp}(D)$ for all $D$, $\sfname{BotSmall}$ and
$\sfname{LNe}$ imply $\sfname{BotConst}$.
\formallabel{GlobalTheory.T.botConst_of_theta_matching}
\formallabel{GlobalTheory.SqTop.botConst_of_inputs_alt}
\formallabel{GlobalTheory.SqTop.BotConstPf.lamC_div_norm}
\end{theorem}

\begin{proof}
We use $\fitGBig$ of the conjunction of three statements (taking the maximum of the thresholds):
\begin{enumerate}[label=(\roman*)]
\item (ratio of bottom constants) $c_0(t)\ne0$, $c_0(\mu_0)\ne0$, and there is $\gamma\in K$ with $\sigma(\gamma)c_0(\mu_0)=c_0(t)$
(\lref{GlobalTheory.T.bottom_constant_ratio_eq_local_factor_ratio}, from $F$, $\sfname{ThetaIntegralsAgree}$,
$\sfname{LocalMatching}$, $\sfname{Inp}$, $\sfname{BotSmall}$ and $\sfname{LNe}$, where the ratio $c_0(t)/c_0(\mu_0)$ is
identified with a ratio of local factors);
\item (first occurrence) there is $\kappa\ne0$ such that for every type $s$ and every bottom pair $(P,P')$ at $s$ there are an
algebraic $u\ne0$ with $|\lambda(u)|_p=1$ for all primes $p\nmid Q_n(G,s)$ and all $\lambda$, and integer matrices $M,M'$ with
\begin{gather*}
\overline{\mathrm{coord}_P(h,T)}=\kappa uc_0(s)\sum_{T'}M(T,T')\mathrm{coord}_{P'}(h,T'),\\
\mathrm{coord}_{P'}(h,T')\,\kappa uc_0(s)=\sum_TM'(T',T)\overline{\mathrm{coord}_P(h,T)}
\end{gather*}
(\lref{GlobalTheory.T.rallis_inner_product_formula_first_occurrence}, from the same hypotheses except $\sfname{LNe}$);
\item $\sfname{BotSmall}$: all bottom pairs are small, and there are no twisted places at $\mu_0$.
\end{enumerate}
Let $\gamma$ be as in (i); $\gamma\ne0$ because $c_0(t)\ne0$. Let $w$ be a finite place of $K$ with $Q_n(G,t)\notin w$ and $p$
the prime under $w$; then $p\nmid Q_n(G,t)$, and $p\nmid Q_n(G,\mu_0)=1$. Apply (ii) to $(P,P')$ at $t$ (giving $u_1,M$) and to
$(P_0,P_0')$ at $\mu_0$ (giving $u_2,M_2$). Then $v=u_1/u_2$ is algebraic with $|\lambda(v)|_p=1$ for all $\lambda$, and
$v_p(\mathrm{lev}_Q(G)Q_n(G,t))=v_p(\mathrm{lev}_Q(G))$. Write $d_l=\mathrm{wt}_{S,e}$ with $(S,e)$ coordinate-integral for
$(\mathrm{bigN}(D),\max\mathrm{Wt}(D),\mathrm{lev}_Q(G))$, and apply \cref{pf:fit:lemmaC} with $Q=\mathrm{lev}_Q(G)$,
$(P_1,P_1')=(P,P')$, $(P_2,P_2')=(P_0,P_0')$, $b_1=\kappa u_1c_0(t)$, $b_2=\kappa u_2c_0(\mu_0)\ne0$ and this $v$: indeed
$b_1=\sigma(\gamma)b_2v$ because $\sigma(\gamma)c_0(\mu_0)=c_0(t)$. This gives $v_w(\gamma)\le v_w(z)+v_w(z')+6d_l(w)$.
\end{proof}

% ---------------------------------------------------------------------------
\subsection{The archimedean bound for normalised toric values}\label{sec:pf:fit:betaarch}

We prove $\sfname{BetaArch}$. Let $d$ be the discriminant parameter of the host ($E$ contains an imaginary quadratic field of
discriminant $-d$). For an index pair $(i,j)$ and $1\le s\le\mathrm{hr}(i,j)$ let $\kappa_s(\tau)\ge1$ be the weights of the
square at $s$, $M_s=2s|I|+2\sum_\tau\kappa_s(\tau)$ (with $I$ the set of embedding types), and $C_0=C_0(2s,\kappa_s)>0$ the CM
value constant. For a toric value $y$ of the character $\mathrm{fam}(G,\mu_0,i,j,c)$ at $s$ and an integral period lattice
$\Lambda$ for $d$ (a lattice of the imaginary quadratic field, \cite{LeanSource2026}) the \emph{normalised toric value} is
\[
\beta=\beta(s,\Lambda,y)=\frac{(C_0\,y)^{12}}{\Delta(\Lambda)^{M_s}},
\]
$\Delta$ being the discriminant of a lattice. $\sfname{BetaArch}$ is $\fitGBig$ of: for all such $(i,j,s,c,y,\Lambda)$, with
$a_0(\beta)$ the constant coefficient of the minimal polynomial of $\beta$ over $\Q$,
\[
\log|a_0(\beta)|\le\deg(\beta)\cdot12B_{\mathrm{TH}}(C,x,b),\qquad B_{\mathrm{TH}}(C,x,b)=Cfb(\log b+\log x).
\]
$\sfname{LBound}$ is $\fitGBig$ of: $\log|\Lambda'(s)|\le B_{\mathrm{TH}}(C,x,b)$ for the completed $L$-function $\Lambda'$ of
every value character with the weights $\kappa_s\circ\pi$ ($\pi$ a permutation of the embedding types) and the conductor of the
grid; $\sfname{TorEx}$ is $\fitGBig$ of: there is a toric value $x_s\ne0$ of $\mathrm{fam}(G,\mu_0,i,j,c)$ at $s$, equal to
$\Lambda(s)$ for every completed $L$-function $\Lambda$ of this character, such that $\beta(s,\Lambda,x_s)$ is algebraic for every
integral lattice $\Lambda$. $\sfname{Transport}$ is the hypothesis of \cref{sec:pf:fit:setting} (the main theorem of complex
multiplication \cite{Shimura1971} in the form used in the development): for an imaginary quadratic field $K$ of discriminant
$-d<-4$, an embedding $\varphi_K$ and a lattice $\Lambda_K$ of the unit ideal, there are ideles $t(\sigma)$, ideals
$\mathfrak a(\sigma)$ and lattices $\Lambda_\sigma$ of $\mathfrak a(\sigma)$, for $\sigma\in\Aut(\C)$ fixing $\varphi_K(K)$,
with $j(\Lambda_\sigma)=\sigma(j(\Lambda_K))$, such that the classes of the $\mathfrak a(\sigma)$ are \emph{equidistributed}
($j(\Lambda_K)$ is algebraic; $\mathfrak a(\sigma),\mathfrak a(\sigma')$ have the same class if and only if $\sigma
j(\Lambda_K)=\sigma'j(\Lambda_K)$; and for every finite Galois extension $L$ containing $\varphi_K(K)$ and $j(\Lambda_K)$, each
class is attained by exactly $[L:\Q]/2h_K$ embeddings of $L$), and such that toric values of CM fields $E\supseteq K$ are
transported by $\sigma$ (algebraicity, and the conjugation formula used below).

\begin{theorem}\label{pf:fit:betaarch}
$\sfname{Transport}$, $\sfname{FieldDat}$, $\sfname{LBound}$ and $\sfname{TorEx}$ imply $\sfname{BetaArch}$.
\formallabel{GlobalTheory.T.betaConjugate_mean_logAbs_bound_of_transport}
\formallabel{GlobalTheory.SqTop.betaArch_of_transport}
\formallabel{GlobalTheory.BetaArchPf.toric_eq_L}
\end{theorem}

\begin{proof}
Take the maximum of the thresholds of the three hypotheses and of the constants $C_x,x_4$ given (for the exponent $\eps$) by
\lref{GlobalTheory.T.log_cmConstant_add_logDerivative_le_squareTopBound}: under the standing conditions, for $x\ge x_4$,
\begin{equation}\label{eq:pf:fit:cmconst}
\log C_0+M_s\,\frac{\log d+|\mathrm{logDer}(1)|}2\le B_{\mathrm{TH}}(C_x,x,b),
\end{equation}
where $\mathrm{logDer}(1)$ is the logarithmic derivative at $1$ of the $L$-function of the quadratic character of discriminant
$-d$. Let $C\ge C_2+C_x$, with $C_2$ the constant of $\sfname{LBound}$, and fix $(i,j,s,c,y,\Lambda)$.

\emph{Set-up.} By $\sfname{FieldDat}$, $E$ is a CM field, Galois over $\Q$ and unramified at all finite places over its
maximal real subfield, $d>4$ is a prime, $-d$ is a square in $E$, the CM type $\varphi$ has at least two embedding types, and the
conductor of the grid satisfies the base condition. By $\sfname{TorEx}$ the toric value is unique, so $y=x_s\ne0$, and $\beta$ is algebraic; $\beta\ne0$ since
$C_0>0$ and $\Delta(\Lambda)\ne0$. There are an imaginary quadratic subfield $K\subseteq E$ of discriminant $-d$ and an
embedding $\varphi_K$, non-real, with $\varphi_\tau|_K=\varphi_K$ for all $\tau$, such that $\Lambda$ is a lattice of the unit
ideal of $K$ (\lref{GlobalTheory.T.exists_imagQuadratic_subfield_of_toricCMFieldData}). Apply $\sfname{Transport}$ to
$(K,d,\varphi_K,\Lambda)$ and obtain $\mathfrak a(\sigma)$, $\Lambda_\sigma$; its transport formula gives, for every $\sigma$
fixing $\varphi_K(K)$, every permutation $\pi$ and every character $\chi'$ conjugate to $\mathrm{fam}(G,\mu_0,i,j,c)$ under
$\sigma$ with weights $\kappa_s\circ\pi$ and every toric sum $T'$ of $\chi'$ at $s$,
\[
|\sigma(\beta)|=(C_0\,|T'|)^{12}\bigl(\Nm(\mathfrak a(\sigma))^6\,|\Delta(\Lambda_\sigma)|\bigr)^{-M_s}.
\]

\emph{Bound for each conjugate.} Let $\sigma$ fix $\varphi_K(K)$. There are a permutation $\pi$ and a character $\chi'$, conjugate
to $\mathrm{fam}(G,\mu_0,i,j,c)$ under $\sigma$, which is a value character with weights $\kappa_s\circ\pi$ for the same conductor
(\lref{GlobalTheory.T.exists_conjugate_cmValueCharacter_permuted_weight}); $\chi'$ has a completed $L$-function $\Lambda'$
(\lref{GlobalTheory.ToricDistribution.exists_isHeckeL_of_hasInf}), and $\Lambda'(s)$ is a toric sum of $\chi'$ at $s$:
by \lref{GlobalTheory.ToricSumValue.toricSum_eq} a toric sum at $s$ is $\Lambda'(s)$ times a product of local factors over
the exceptional places, and here there are none. By $\sfname{LBound}$ (with the constant $C-C_x\ge C_2$), $\log|\Lambda'(s)|\le B_{\mathrm{TH}}(C-C_x)$;
in particular $|\Lambda'(s)|>0$ since $\sigma(\beta)\ne0$. Hence
\[
\log|\sigma(\beta)|\le\Phi(\sigma):=12\log C_0+12B_{\mathrm{TH}}(C-C_x)-M_s\log\bigl(\Nm(\mathfrak a(\sigma))^6|\Delta(\Lambda_\sigma)|\bigr).
\]
The function $\Phi$ only depends on the class of $\mathfrak a(\sigma)$: if $\Lambda_a,\Lambda_b$ are lattices of ideals
$\mathfrak a,\mathfrak b$ in the same class, then $\Nm(\mathfrak a)^6|\Delta(\Lambda_a)|=\Nm(\mathfrak b)^6|\Delta(\Lambda_b)|$
(\lref{GlobalTheory.T.fractionalIdealNorm_pow_six_mul_discriminant_class_invariant}; $\Delta$ is homogeneous of weight $12$).

\emph{Mean over the conjugates.} By the equidistribution of the classes, there are representatives $\sigma_c$, one for each
ideal class $c$ of $K$, fixing $\varphi_K(K)$ with $[\mathfrak a(\sigma_c)]=c$, such that
\[
\log|a_0(\beta)|\le\deg(\beta)\cdot\frac1{h_K}\sum_c\Phi(\sigma_c)
\]
(\lref{GlobalTheory.T.mean_log_norm_conjugates_over_ideal_classes}: $|a_0(\beta)|$ is the product of the absolute values of
the conjugates of $\beta$, and the conjugates are distributed evenly over the classes). By the Chowla--Selberg formula
\cite{ChowlaSelberg1967} in the form \lref{GlobalTheory.T.classMean_log_norm_disc_chowlaSelberg},
\[
\frac1{h_K}\sum_c\log\bigl(\Nm(\mathfrak a_c)^6|\Delta(\Lambda_c)|\bigr)=12\log(2\pi)-6\log d+6\gamma_{\mathrm E}-6\,\mathrm{logDer}(1),
\]
$\gamma_{\mathrm E}$ being Euler's constant. Since $\log(2\pi)>0$, $\gamma_{\mathrm E}\ge0$ and $M_s\ge0$,
\begin{align*}
\frac1{h_K}\sum_c\Phi(\sigma_c)&\le12\log C_0+12B_{\mathrm{TH}}(C-C_x)+6M_s\bigl(\log d+|\mathrm{logDer}(1)|\bigr)\\
&\le12B_{\mathrm{TH}}(C_x)+12B_{\mathrm{TH}}(C-C_x)=12B_{\mathrm{TH}}(C)
\end{align*}
by \eqref{eq:pf:fit:cmconst} and the linearity of $B_{\mathrm{TH}}$ in $C$. This proves the bound.
\end{proof}

% ---------------------------------------------------------------------------
\subsection{The archimedean size of norms of pairs}\label{sec:pf:fit:archnorm}

We prove the archimedean bound for norms of bottom pairs used in \cref{pf:fit:leafsq}(e). For a number field $K$ and a
function $\phi$ on its infinite places, $\fitmi{\phi}=\frac2{[K:\Q]}\sum_v\phi(v)$. For a pair $P$ with parameters
$(\delta,t_0)$, an embedding $\rho$ of a field containing the invariants $g_2,g_3$ of the CM lattice $\Lambda_\delta$ of the
curve gives the conjugate lattice $\Lambda_\delta^\rho$; we write $A_\rho$ for its area and $\nu_\rho$ for the squared norm of its
theta-null value at a two-torsion point not killed by the curve.

\begin{lemma}[Archimedean size of norms]\label{pf:fit:archnorm}
There is an absolute constant $C_A>0$ with the following property. Let $P$ be a pair at level $Q\ge1$, small for $(N,B)$ with
$N\ge1$ and $B\ge1$, let $L\ge1$ with $\log\mathrm{rd}_E\le L$ and $|\mathrm{logDer}(1)|\le L$, and let $x$ be an element of a
totally complex number field $K$ with $\sigma(x)=\|P\|$ for an embedding $\sigma$. Then
\[
\fitmi{\log^+|\sigma_vx|}\le C_A\Bigl(2fNB\bigl(L+\log(2NB)\bigr)+fN^2\bigl(L+\log(2fNB)\bigr)+N^2\log Q\Bigr),
\]
where $f$ is the number of embedding types.
\formallabel{GlobalTheory.T.norm_meanInf_posLog_le_definiteThetaPair}
\formallabel{GlobalTheory.ArchFree.norm.norm_meanInf_posLog_le}
\formallabel{GlobalTheory.ArchFree.st.archSt}
\end{lemma}

\begin{proof}
Let $C_0$ be the constant of \lref{GlobalTheory.T.conjugate_lattice_means_posLog_area_theta_bounded} (for every finite
extension $M$ containing $g_2,g_3$ and the Weierstrass value at $t_0$, and two-torsion points $t_\rho$ of the conjugate
lattices, the averages over the embeddings $\rho\colon M\to\C$ of $\log^+A_\rho^{-1}$ and of $\log^+\nu_\rho$ are at most
$C_0L$), and put $C_A=1100+4C_0$. The coordinates $\mathrm{coord}_P(h,T)$ are algebraic, and for every embedding
$\rho\colon\Qbar\to\C$ and every two-torsion point $t'$ of $\Lambda_\delta^\rho$ not killed by the curve,
\[
|\rho(\mathrm{coord}_P(h,T))|^2\le c_a^2\,Q^{ad}\,(a!^{\,s})^2\,3\,(2fad\,\mathrm{rd}_E+1)^{fad}\,\nu_\rho^{fad}
\Bigl(\frac{2e\pi B}{A_\rho}\Bigr)^{as},
\]
where $c_a=\#\mathrm{Cls}(a,K_a)$ and $s=\sum_\tau n(\tau)$
(\lref{GlobalTheory.T.definiteThetaPair_coordinate_conjugates_bounded}). The invariants $g_2,g_3$ and the Weierstrass value
$\wp(t_0)$ are algebraic, and every conjugate lattice has a two-torsion point not killed by the curve
(\lref{GlobalTheory.T.admJ_and_torsion_data_algebraic} and \lref{GlobalTheory.T.exists_twoTorsionCMPoint_all_conjugates}). Let
$M\subseteq\C$ be a finite Galois extension of $\Q$ containing $\sigma(K)$, all coordinates $\mathrm{coord}_P(h_q,T)$ at fixed
representatives $h_q$ of the classes $q\in\mathrm{Cls}(d,J)$, $g_2$, $g_3$ and $\wp(t_0)$
(\lref{GlobalTheory.ArchNorm.exists_galois}), and choose a two-torsion point $t_\rho$ for each $\rho\colon M\to\C$.

By \eqref{eq:pf:fit:pairnorm}, $\sigma(x)=\sum_{(q,T)}|\Gamma_{h_q}|^{-1}\,z_{q,T}\overline{z_{q,T}}$ with $z_{q,T}=
\mathrm{coord}_P(h_q,T)\in M$ and rational weights $|\Gamma_{h_q}|^{-1}\in(0,1]$. Hence
(\lref{GlobalTheory.ArchMean.meanInf_posLog_le}: for such a representation, the mean of $\log^+|\sigma_vx|$ over the infinite
places of $K$ is at most the average over $\rho\colon M\to\C$ of $\log^+\sum|\Gamma_{h_q}|^{-1}|\rho(z_{q,T})|^2$)
\[
\fitmi{\log^+|\sigma_vx|}\le\frac1{[M:\Q]}\sum_{\rho}\log^+\sum_{(q,T)}|\Gamma_{h_q}|^{-1}|\rho(z_{q,T})|^2 .
\]
Fix $\rho$ and extend it to $\Qbar$. Each term is at most the bound above. The number of terms is
$\#\mathrm{Cls}(d,J)\cdot\#\mathrm{Tup}(a,d,n)$, where $\#\mathrm{Cls}(d,J)\le c_d\,[K_d:J]\le c_dQ^{d^2+4}$, the class numbers
satisfy $c_N\le(12f^2N^4\mathrm{rd}_E^2)^{6fN^2}$ (\lref{GlobalTheory.NormSum.cls_finite_card_le}), and
$\#\mathrm{Tup}(a,d,n)\le2^{ds}$ (\lref{GlobalTheory.MinorBound.minor_bound}). With $d\le2N$, $s\le2fB$ (smallness),
$\log\mathrm{rd}_E\le L$ and $\mathrm{rd}_E\ge1$, an elementary estimate of the product
(\lref{GlobalTheory.ArchNorm.sizes}) gives
\[
\log^+\sum_{(q,T)}|\Gamma_{h_q}|^{-1}|\rho(z_{q,T})|^2\le E_0+4fN^2\log^+\nu_\rho+4fNB\log^+A_\rho^{-1},
\]
with $E_0=1100\bigl(2fNB(L+\log(2NB))+fN^2(L+\log(2fNB))+N^2\log Q\bigr)$. Averaging over $\rho$ and using the bounds $C_0L$ for
the averages of $\log^+\nu_\rho$ and $\log^+A_\rho^{-1}$, the right-hand side becomes at most
$E_0+4C_0fN^2L+4C_0fNBL\le(1100+4C_0)\bigl(2fNB(L+\log(2NB))+fN^2(L+\log(2fNB))+N^2\log Q\bigr)$, all logarithms being
non-negative.
\end{proof}

\begin{theorem}\label{pf:fit:archgrid}
$\fitGBig$ holds for the statement: if $z_A,z'_A$ are functions on the index pairs with values in $K$ such that, for some
$\delta,t_0$ and every $\eta$, there is a bottom pair $(P,P')$ at $\mu_0$ at level $\mathrm{lev}_Q(G)$ with
$\sigma(z_A(\eta))=\|P\|$ and $\sigma(z'_A(\eta))=\|P'\|$, then for every $\eta$, $\fitmi{\log^+|\sigma_vz_A(\eta)|}$ and
$\fitmi{\log^+|\sigma_vz'_A(\eta)|}$ are at most
\[
2CfNB_w\bigl(\log b+\log(2NB_w)\bigr)+CfN^2\bigl(\log b+\log(2fNB_w)\bigr)+CN^2\log\mathrm{lev}_Q(G),
\]
with $N=\mathrm{bigN}(D)$ and $B_w=\max\mathrm{Wt}(D)$.
\formallabel{GlobalTheory.T.archimedeanBound_bottomPairs_norms_grid}
\formallabel{GlobalTheory.ArchGrid.arch_grid}
\end{theorem}

\begin{proof}
Take $C_{F0}=0$, $C_0=C_A$ and $x_0=0$. Under the standing conditions, $L=\log b=x^\eps\ge1$, the size rows give
$\log\mathrm{rd}_E\le\log b$, and $|\mathrm{logDer}(1)|\le\log b$; moreover $N\ge1$ (as $m+3\le k$), $B_w\ge1$ and $f\ge1$.
Bottom pairs at the base type are small for $(N,B_w)$ (\lref{GlobalTheory.T.squareBottomPairs_small_reference_tuple}).
\Cref{pf:fit:archnorm}, applied to $P$ and to $P'$ with $Q=\mathrm{lev}_Q(G)\ge1$, gives the bound with $C_A$ in place of $C$; as
the bracket is non-negative, it holds for every $C\ge C_A$.
\end{proof}

% ---------------------------------------------------------------------------
\subsection{Size bounds for class sums: assembly}\label{sec:pf:fit:sizes}

We now prove the components of \cref{pf:fit:sumsize}. Fix a point in Hida setting and an admissible choice $\mathrm{ch}$ at
scale $(C_F,b)$ satisfying the field condition, with grid $G$, sum choice $S$, class representatives and enumeration of the
test indices. By construction \cite{LeanSource2026}, the $a$-th class sum $v_K(\mathrm{ch},t,F)_a$ of a form $F$ of type $t$ is
$s_{AA}$ times a finite sum, over the cosets $\beta$ in the support of the $a$-th datum $\Phi_a$, the classes $c$ of $\beta$ and
the pure data $d$ of a coset expansion of $\beta$, of products of: the coefficients $\prod_v\mathrm{coef}_v(\beta)$ and the data
$\prod_v\mathrm{dat}_v(\beta)$ of the terms of $\beta$; the inverse stabiliser order $\mathrm{invGam}(c)$; and the last pairing
$\langle u,u'\rangle$ of the image $u$ of the test vector $\psi_C$ of $\Psi_a$ under the $m$ \emph{top maps} $k(1,d_1(j))$ (from
the node $(\mathrm{top},j)$ to $(\mathrm{top},j+1)$ at the points of $c$) with the image $u'$ of the form value $f_t(F)$ under the
$m$ \emph{row maps} $k(1,d_2(i))$. Here $\sigma=\sigma_{v_0}$.

\paragraph{Losses.} A \emph{loss function} with bound $B$ is a function $d\ge0$ on the finite places of $K$ with finite support
and $\sum_{w\in T}\ell_wd(w)\le B$ for every finite $T$. A complex number $z$ has \emph{loss at most $a$ at $w$} if $z=\sigma(y)$
with $y\in K$ and $y=0$ or $v_w(y)\ge-a$; a vector $v$ is \emph{in the lattice $\Lambda$ up to $a$ at $w$} if $v=\sigma(u)$ for a
$K$-rational $u$ and $su\in\Lambda$ for some $s\in K^\times$ with $v_w(s)\le a$. Losses add under products. A place $w$ is
\emph{away} from the level at $t$ if every prime of the level at $t$ lies over $S_0$ or is prime to the residue characteristic of
$w$. Put $b_1(C)=Cfm\,\mathrm{bigN}^2(\log b+\sigma(G))\,b$ and $a_1(C)=Cf\,\mathrm{bigN}\,b\log b+Cm\,\mathrm{bigN}^2\sigma(G)$
(\cref{sec:pf:fit:inputs}); both are linear in $C$.

\paragraph{The factor claims.} Each of the following is a $\fitBigC$ statement about $\mathrm{ch}$, with constants $(C,b)$.
\begin{itemize}
\item $\sfname{DepthK}$: every coset $\beta$ of the support has a coset expansion of depth $e(\beta)$, and $v_l(e(\beta))\le
\mathrm{dep}_l(G)$ for every prime $l$ under a place $w$ away from the level.
\item Finite claims (loss functions with bound $m\,b_1$ for the first four, $b_1$ for the others): $\sfname{FinScal}$ (the scalar
$s_{AA}$), $\sfname{FinCoef}$ (the coefficients), $\sfname{FinDat}$ (the data products, at the column and row points),
$\sfname{FinGam}$ (the numbers $\mathrm{invGam}(c)$) have loss at most $d(w)$ at $w$; $\sfname{FinTop}$, $\sfname{FinRow}$: the
top and row maps send vectors in the lattice of the node up to $a$ at $w$ to vectors in the lattice of the next node up to
$a+d(w)$; $\sfname{FinPair}$: the last pairing of vectors in the lattices up to $a,a'$ has loss at most $a+a'+d(w)$;
$\sfname{FinForm}$: the form values $f_t(F)$ of lattice forms are in the lattice up to $d(w)$; $\sfname{FinPsi}$: the test
vectors are in the lattice of the node $(\mathrm{top},0)$ up to $0$.
\item Infinite claims, at the base type $\mu_0$ and for every embedding $\sigma'$ of $K$: $\sfname{InfScal}$:
$|\sigma'(s_{AA})|\le e^{ma_1}$; $\sfname{InfTens}$: the products of coefficients and data have absolute value at most
$e^{Cm^2\mathrm{bigN}^2\sigma(G)}$; $\sfname{InfCount}$: the number of data, the size of every support and the number of terms of
every coset expansion are at most $e^{ma_1}$; $\sfname{InfClass}$: the number of classes of the grid and of every coset are at
most $e^{ma_1}$; $\sfname{InfTop}$, $\sfname{InfRow}$: there is a function $g$ on the infinite places with $\fitmi{g}\le a_1$
such that the top (resp.\ row) maps multiply the norms of vectors by at most $e^{g(p)}$ at $p$; $\sfname{InfPair}$: the last
pairing satisfies $|\langle u,u'\rangle|^2\le e^{a_1}\|u\|^2\|u'\|^2$; $\sfname{InfPsi}$: the norm of a test vector is bounded by
the norm of its coefficient vector for the Gram matrix of the test forms; $\sfname{InfForm}$:
$\|f_{\mu_0}(F_x)\|^2\le e^{ma_1}\,{}^t\bar x H_V(\mu_0)x$.
\end{itemize}

\begin{theorem}[Finite assembly]\label{pf:fit:finasm}
$\sfname{SumOKK}$, $\sfname{DepthK}$ and the nine finite claims imply $\fitBigC(\sfname{FinAt})$.
\formallabel{GlobalTheory.Adopt.finiteSumDenominatorBound_of_factor_claims}
\formallabel{DenominatorLossAdditivity.GlobalTheory.CompanionTheory.FinDen.fin_of_factors}
\formallabel{DenominatorLossAdditivity.GlobalTheory.CompanionTheory.FinDen.fin_at}
\formallabel{DenominatorLossAdditivity.GlobalTheory.CompanionTheory.FinDen.loss_comb}
\formallabel{DenominatorLossAdditivity.GlobalTheory.CompanionTheory.FinDen.loss_mul}
\end{theorem}

\begin{proof}
Take the maximum $C_{F0}$ of the eleven thresholds, for $C_F\ge C_{F0}$ the constant $16\max_ic_i$ (with $c_i$ the constants of
the claims), and for $C\ge16\max_ic_i$ apply each claim with the constant $C/16$; take the maximum of the thresholds $x_0$.
Let $d_S,d_C,d_D,d_G$ (bound $m\,b_1(C/16)$) and $d_T,d_R,d_F,d_P$ (bound $b_1(C/16)$) be the loss functions of the claims, and
\[
d_l=d_S+d_C+d_D+d_G+m\,d_T+m\,d_R+d_F+d_P .
\]
It is non-negative with finite support, and, as $m\ge1$, $\sum_{w\in T}\ell_wd_l(w)\le4m\,b_1(C/16)+2m\,b_1(C/16)+2b_1(C/16)\le
8m\,b_1(C/16)=m\,b_1(C)/2$. Let $(X',M)$ be a test lattice, $t$ a type, $w$ away from the level at $t$, $x$ a lattice form of type
$t$ and $y\in K^n$ with $\sigma(y_a)=v_K(\mathrm{ch},t,F_x)_a$. By
\lref{DenominatorLossAdditivity.GlobalTheory.CompanionTheory.FinDen.exists_intMod_levelPrimePow_eq_add_mul} ($M$ being the
dual of the lattice of the test forms, it suffices that for every datum $\Phi$, every class $c$ and every vector $u$ of the weight
lattice of the representative of $c$, the class sum core with the test vector $\psi_C(\delta(c,u))$ has loss at most $d_l(w)$ at
$w$), it is enough to bound these losses. The core is a finite sum of products of the factors listed above; by
\lref{DenominatorLossAdditivity.GlobalTheory.CompanionTheory.FinDen.core_loss} (losses add under products and are preserved by
sums; the coset expansions exist by $\sfname{DepthK}$) and the claims, it has loss at most
$d_C+d_D+d_G+m\,d_T+m\,d_R+d_P+d_F$ at $w$, the test vector contributing $0$ by $\sfname{FinPsi}$, the $m$ top maps and the $m$ row
maps contributing $m\,d_T$ and $m\,d_R$. Multiplying by $s_{AA}$ adds $d_S(w)$.
\end{proof}

\begin{theorem}[Archimedean assembly]\label{pf:fit:infasm}
$\sfname{SumOKK}$, $\sfname{DepthK}$ and the nine infinite claims imply $\fitBigC(\sfname{InfAt})$.
\formallabel{GlobalTheory.T.archimedeanSumSizeBound_largeData_from_factors}
\formallabel{GlobalTheory.Adopt.archimedeanSumSizeBound_largeData_from_factors}
\formallabel{DeltaFormSums.GlobalTheory.CompanionTheory.InfAssembly.infK_of_factors}
\formallabel{DeltaFormSums.GlobalTheory.CompanionTheory.InfAssembly.infAt_of}
\end{theorem}

\begin{proof}
As in \cref{pf:fit:finasm} we apply the claims with the constant $C/16$ (now $C\ge0$, $b\ge1$), and write $A=a_1(C/16)$,
$S_g=(C/16)m^2\mathrm{bigN}^2\sigma(G)\le mA$, and $g,g'$ for the functions of $\sfname{InfTop}$ and $\sfname{InfRow}$. Let
$(X',M)$ be a test lattice, $X$ a metric of $H_V(\mu_0)$ and $V$ a matrix with $\sigma(Vx)_a=v_K(\mathrm{ch},\mu_0,F_x)_a$. By
\lref{DeltaFormSums.GlobalTheory.CompanionTheory.InfAssembly.opNorm_bound} (the components of $Vx$ are paired with test
vectors through the dual of the Gram matrix of the test forms, and each class sum is bounded factor by factor by the claims), at
every infinite place $p$
\[
\|V\|_p^2\le\exp\bigl(11mA+2S_g+A+m\,g(p)+m\,g'(p)\bigr).
\]
If $V=0$ there is nothing to prove; otherwise $\|V\|_p>0$, and averaging the logarithms over the infinite places,
$h(V)\le11mA+2S_g+A+m\fitmi{g}+m\fitmi{g'}\le11mA+2mA+mA+mA+mA=16mA=m\,a_1(C)$, using $\fitmi g,\fitmi{g'}\le A$ and $A\le mA$.
\end{proof}

% ---------------------------------------------------------------------------
\subsection{Size bounds for class sums: the factors}\label{sec:pf:fit:factors}

We prove the factor claims used in \cref{pf:fit:sumsize} which carry the main estimates. Throughout, $F=\sfname{ThetaLocalInputs}$;
$N=\mathrm{bigN}(D)$; for a local place $w$ of the grid at a type, $e(w)$ is its coset exponent and $q_w$ its residue
cardinality; and $\sigma(G)=\sum_{v\in S_0}\mathrm{depth}(v)\log\Nm v$.

\begin{theorem}[Depth of cosets]\label{pf:fit:depth}
$F$ and $\sfname{SumOKK}$ imply $\sfname{DepthK}$.
\formallabel{GlobalTheory.T.depthK_of_sumOKK}
\formallabel{GlobalTheory.DepthUnf.depthK_of_hypotheses_alt}
\end{theorem}

\begin{proof}
For a coset $\beta$ of the support of a datum $\Phi$ at the type $t$, a \emph{coset expansion of depth $k$} is an expansion of
the product of the terms of $\beta$ in indicator functions of cosets of depth $k$, and for an idempotent sum $e$ of local
places the product is \emph{depth periodic with period $k$} if it is invariant at depth $k$ \cite{LeanSource2026}; $e(\beta)$
denotes the least depth of a coset expansion. By \lref{GlobalTheory.T.flatTensor_depth_conditions_away_from_twist}, for large
data there is $K>0$ with $v_l(K)\le\mathrm{dep}_l(G)$ for every prime $l$ under a place away from the level at $t$, such that the
product is depth periodic with period $K$ for every idempotent sum. The sum choice gives an expansion $\tau$ of $\beta$ for some
idempotent sum $e$ (\lref{GlobalTheory.FunObj.GridChoice.sumOK_ex}); since the product is depth periodic with period $K$,
there is a coset expansion of depth $K$ (\lref{GlobalTheory.T.exists_cosetIndicatorExpansion_of_depthPeriodic}). Hence
$e(\beta)\le K$ is defined and attained by an expansion $\tau_m$. Conversely, a coset expansion of depth $e(\beta)$ makes the
product depth periodic with period $e(\beta)$ for some idempotent sum $e'$ with an expansion
(\lref{GlobalTheory.T.depthPeriodic_of_cosetIndicatorExpansion}); periods are stable under $\gcd$
(\lref{GlobalTheory.T.flatProductDepthPeriodic_gcd}), so the product is depth periodic with period $\gcd(e(\beta),K)$ for
$e'$, and there is a coset expansion of depth $\gcd(e(\beta),K)$. By minimality $e(\beta)\le\gcd(e(\beta),K)$, so
$e(\beta)\mid K$ and $v_l(e(\beta))\le v_l(K)\le\mathrm{dep}_l(G)$.
\end{proof}

\begin{theorem}[Finite data]\label{pf:fit:findat}
$F$ implies $\sfname{FinDatK}$.
\formallabel{GlobalTheory.T.finDatK_of_theta_inputs}
\formallabel{GlobalTheory.FunDat.finDatK_alt}
\formallabel{GlobalTheory.FunDat.finDatK_of}
\formallabel{GlobalTheory.FunDat.away_cases}
\end{theorem}

\begin{proof}
Two estimates for the local tensors are used, for large data and at local places $w$ of the grid at $t$ which lie in the set
$W_0$ of the grid or whose residue characteristic is prime to the place $\mathfrak P$ of $K$ under consideration (this holds at
the places away from the level: a local place outside $W_0$ is a level place; if it lay over $S_0$ it would be the base place,
which is not a level place, or one of the chosen places $w(v)\in W_0$; so its residue characteristic is prime to $\mathfrak P$ by
the definition of \emph{away}):
\begin{enumerate}[label=(\alph*)]
\item the \emph{flip values} $T_{\mathrm{flip}}(t,\Phi,w,z)$ have loss at most $c_1m^2N^2b\,e(w)\,v_{\mathfrak P}(q_w)$ at
$\mathfrak P$ (\lref{GlobalTheory.T.flips_loss_bounded_at_exceptional_places});
\item if all flip values at $w$ have loss at most $c\ge0$ at $\mathfrak P$, then the data $\mathrm{dat}_w(\beta)(X,Y)$ of the
terms of every coset $\beta$ have loss at most $c+c_2m^2N^2e(w)\,v_{\mathfrak P}(q_w)$
(\lref{GlobalTheory.T.movedTensor_unfolding_integral_valuation_loss_bound}).
\end{enumerate}
Let $K_0=\max(c_1,0)+\max(c_2,0)$, $C\ge\max(2,2K_0)$ and $\kappa=K_0m^2N^2b$. Then each datum at $w$ has loss at most
$\kappa\,e(w)v_{\mathfrak P}(q_w)$ (as $b\ge1$), and the product over the local places has loss at most
$d_D(\mathfrak P)=\kappa\sum_we(w)v_{\mathfrak P}(q_w)$. The function $d_D$ is a loss function with bound $2\kappa\sigma(G)$
(\lref{GlobalTheory.FunDat.loss_dlD}), and $2\kappa\sigma(G)\le C\sigma(G)m^2N^2b\le m\,b_1(C)$ because $f\ge1$.
\end{proof}

\begin{theorem}[The last pairing]\label{pf:fit:fieldpair}
$F$ implies $\sfname{FinPair}$ and $\sfname{InfPair}$.
\formallabel{GlobalTheory.T.field_pair_comparison}
\formallabel{GlobalTheory.FPair.lift_splits_locally}
\formallabel{GlobalTheory.FPair.finPair}
\formallabel{GlobalTheory.FPair.infPair}
\end{theorem}

\begin{proof}
The last pairing is the branching pairing between the representations of highest weights $\lambda_{\mathrm{top}}$ (of the node
$(\mathrm{top},m)$) and $\lambda_{\mathrm{grid}}$ (of the node $(m,m)$); it is defined over $K$, and its complex values are the
images of its $K$-rational values (\lref{GlobalTheory.T.branchingPairing_extendScalars}). For every embedding type the two
weights interlace strictly (\lref{GlobalTheory.WtInterlace.weights_top_grid}).

\emph{Finite places.} For large data there are $D_n\in\N_{>0}$ with $\log D_n\le K_0fN^2b\log b$ and $p^{E_p}\mid D_n$ for every
prime $p$, where $E_p=\sum_\tau E_p(\lambda_{\mathrm{top},\tau})$ is the total exponent of the top weights at $p$
(\lref{GlobalTheory.T.exists_totalExponent_topWeight_bound}). Put $d_P(\mathfrak P)=v_{\mathfrak P}(D_n)\ge0$; it has finite
support and $\sum_{\mathfrak P\in T}\ell_{\mathfrak P}d_P(\mathfrak P)\le\log D_n$ for every finite $T$
(\lref{GlobalTheory.T.placeValuation_localDegree_arithmetic}), which is at most $b_1(C)$ for $C\ge\max(K_0,0)$. If $u=\sigma(u_0)$
and $v=\sigma(v_0)$ with $su_0$ and $s'v_0$ in the lattices of the two nodes, $v_{\mathfrak P}(s)\le a$, $v_{\mathfrak P}(s')\le a'$,
then the branching pairing of lattice vectors is integral up to the exponents:
$v_{\mathfrak P}(\langle su_0,s'v_0\rangle)\ge-v_{\mathfrak P}(p^{E_p})$, $p$ the prime under $\mathfrak P$
(\lref{GlobalTheory.T.branchingPairing_integral_upToExponents}). Hence either $\langle u_0,v_0\rangle=0$ or
$v_{\mathfrak P}(\langle u_0,v_0\rangle)\ge-a-a'-v_{\mathfrak P}(D_n)$, since $p^{E_p}\mid D_n$.

\emph{Infinite places.} Let $P_{\mathrm{br}}=\prod_\tau\mathrm{br}(\lambda_{\mathrm{top},\tau},\lambda_{\mathrm{grid},\tau})^2$ be
the product of the squared branching constants. For large data $\log P_{\mathrm{br}}\le K_0'fNb\log b\le a_1(C)$
(\lref{GlobalTheory.T.archimedean_branching_constant_bound}), and for every embedding $\sigma'$,
$|\sigma'\langle u_0,v_0\rangle|^2\le P_{\mathrm{br}}\|\sigma'u_0\|^2\|\sigma'v_0\|^2$ with the Fock norms of the coefficient vectors
(\lref{GlobalTheory.T.branchingPairing_archimedean_bound_fockNorm}). If $\|\sigma'u_0\|^2\le B$ and $\|\sigma'v_0\|^2\le B'$
this gives $|\sigma'\langle u_0,v_0\rangle|^2\le e^{a_1}BB'$.
\end{proof}

\begin{theorem}[Form values]\label{pf:fit:formval}
$F$ implies $\sfname{FinForm}$ (with the loss function $0$) and $\sfname{InfForm}$.
\formallabel{GlobalTheory.T.form_value_finite_infinite_factorization}
\formallabel{GlobalTheory.FormVal.form_value_eq_finite_mul_infinite}
\formallabel{GlobalTheory.FormVal.form_fin_inf_of}
\end{theorem}

\begin{proof}
For large data: the form values $f_t(F_x)$ of lattice forms $x$ lie in the lattice of the node $(0,m)$ up to $0$ at every place
(\lref{GlobalTheory.T.integralGoodForm_values_mem_scaledDualWeylLattice}); this is $\sfname{FinForm}$ with $d_F=0$, a loss
function for every bound $b_1\ge0$. At the base type, ${}^t\bar xH_V(\mu_0)x\ge0$ and
$\|f_{\mu_0}(F_x)\|^2\le\#P\cdot D_\Gamma\cdot{}^t\bar xH_V(\mu_0)x$ at every embedding, where $\#P$ is the order of the finite
group of the level of the family and $D_\Gamma$ a positive integer attached to the definite side \cite{LeanSource2026} (\lref{GlobalTheory.T.baseTupleGoodForm_values_bounded_by_classSum}); and $b_1\ge0$,
$\#P\cdot D_\Gamma\le e^{ma_1}$ (\lref{GlobalTheory.T.kernelLossBound_nonneg_and_level_card_le_exp}). Together these give
$\sfname{InfForm}$.
\end{proof}

\begin{theorem}[Tensor coefficients and data]\label{pf:fit:tensor}
$F$ implies $\sfname{InfTens}$.
\formallabel{GlobalTheory.T.sum_size_archimedean_tensor_bound}
\formallabel{GlobalTheory.SumSize.archimedean_tensor_size_bound_of_local}
\formallabel{GlobalTheory.ArchTens.infTensAt_of_places}
\formallabel{GlobalTheory.ArchTens.sum_placeSz_le}
\end{theorem}

\begin{proof}
For a local place $w$ of the grid at $\mu_0$ let $s(w)=e(w)\log\Nm(w\cap\OO_F)$; then $\sum_ws(w)\le\sigma(G)$, because the
local places outside $W_0$ have coset exponent $0$ and $W_0$ is the image of $S_0$ under $v\mapsto w(v)$, with
$e(w(v))=\mathrm{depth}(v)$ and $w(v)\cap\OO_F=v$. Consider, for a constant $c$, the statements: (Flip$_c$) for every embedding
$\sigma'$, $|\sigma'T_{\mathrm{flip}}(\mu_0,\Phi,w,z)|\le e^{cm^2N^2s(w)}$; (Dat$_c$) the same for the data
$\mathrm{dat}_w(\beta)(X,Y)$; (Coef$_c$) $|\mathrm{coef}_w(\beta)|\le e^{cm^2N^2s(w)}$. For large data: Flip$_c$ holds for
every $c\ge48$ (\lref{GlobalTheory.T.local_character_power_conductor_bound}); Flip$_{c/2}$ implies Dat$_c$ for $c\ge c_2$
(\lref{GlobalTheory.T.archimedean_tensor_factor_unfolding}); and Coef$_c$ holds for $c\ge c_3$
(\lref{GlobalTheory.T.archimedean_tensor_volume_bound}). For $C\ge\max(4\cdot48,2c_2,2c_3,0)$, Flip$_{C/4}$ holds, hence
Dat$_{C/2}$, and Coef$_{C/2}$ holds. Multiplying over the
local places, the product of the coefficients and the product of the data have absolute values at most
$\exp(\tfrac C2m^2N^2\sum_ws(w))$ each, so their product is at most $\exp(Cm^2N^2\sigma(G))$.
\end{proof}

\begin{theorem}[Counting data, cosets and terms]\label{pf:fit:sumcount}
$F$ and $\sfname{DepthK}$ imply $\sfname{InfCount}$.
\formallabel{GlobalTheory.T.sum_count_bound_archimedean}
\formallabel{GlobalTheory.SumCount.archimedean_sum_count_bound}
\formallabel{GlobalTheory.SumCount.infCountK_of}
\end{theorem}

\begin{proof}
$\sfname{InfCount}$ is the conjunction of three $\fitBigC$ statements, and the thresholds are the maxima of theirs: the number of
data of the grid is at most $e^{ma_1}$ (\lref{GlobalTheory.T.datumCountBound_largeData}); the support of every datum has at
most $e^{ma_1}$ cosets (\lref{GlobalTheory.T.theta_kCoset_count_formula}, a counting formula for the cosets, from $F$); and
every coset expansion of depth $e(\beta)$ has at most $e^{ma_1}$ terms
(\lref{GlobalTheory.T.pureTensorCountBound_of_cosetDepthBound}, from the depth bound $\sfname{DepthK}$).
\end{proof}

\begin{theorem}[Counting classes]\label{pf:fit:classcount}
$F$ implies $\sfname{InfClass}$.
\formallabel{GlobalTheory.T.class_count_bound_archimedean}
\formallabel{GlobalTheory.ClassCount.archimedean_class_count_bound}
\end{theorem}

\begin{proof}
For large data the size rows give $N,f\le b$, $12\le b$, $1\le\mathrm{rd}_E$ and $\log\mathrm{rd}_E\le\log b$, so
$h(n)=(12f^2n^4\mathrm{rd}_E^2)^{6fn^2}$ satisfies $\log h(N)\le48fN^2\log b$. The norm of the modulus of the grid is at most
$e^{2\sigma(G)}$ (\lref{GlobalTheory.T.absNorm_gridModulus_le_exp_depthLogNormSum}), and for every definite node $X$ of the
grid (of rank at most $N$) and every subgroup $K$ containing the principal congruence subgroup of this modulus, the number of
classes at level $K$ is at most $h(\mathrm{rk}X)\,\Nm(\mathfrak m)^{(\mathrm{rk}X)^2}$
(\lref{GlobalTheory.T.definiteVertexClasses_card_le_principalCongruence}). Since $h$ is increasing, all these numbers are at most
$M=h(N)e^{2\sigma(G)N^2}\le e^{K_0}$, $K_0=48fN^2\log b+2N^2\sigma(G)$. The levels of the classes of the grid and of a coset at
the $m$ top nodes and the $m+1$ row nodes contain the principal congruence subgroup
(\lref{GlobalTheory.T.principalCongruence_le_topLevel_classSum} and \lref{GlobalTheory.T.classSum_vertexLevels_contain_principalCongruence}),
so the numbers of classes are at most $M^{2m+1}\le e^{(2m+1)K_0}$. Finally $(2m+1)K_0\le3mK_0\le m\,a_1(C)$ for $C\ge144$, because
$144fN^2\log b\le Cf N b\log b$ (as $N\le b$) and $6N^2\sigma(G)\le CmN^2\sigma(G)$.
\end{proof}

\begin{theorem}[Row maps at the infinite places]\label{pf:fit:opposite}
$F$, $\sfname{Inp}(D)$ for all $D$, the dictionary of local bounds and $\sfname{DepthK}$ imply $\sfname{InfRow}$.
\formallabel{GlobalTheory.T.sum_size_bound_opposite_conjugation}
\formallabel{GlobalTheory.SumSize.OppConj.archimedean_sum_size_opposite_conjugate}
\formallabel{GlobalTheory.SumSize.OppConj.infRowK_of}
\formallabel{GlobalTheory.SumSize.OppConj.infRowAt_of}
\end{theorem}

\begin{proof}
For a linear map $k$ between spaces of vectors with $K$-structure let $\bar k(v)=\overline{k(\bar v)}$. By
\lref{GlobalTheory.T.conj_mean_kernel_pairing} (from $F$, $\sfname{Inp}$, the dictionary and $\sfname{DepthK}$), for large
data there is a function $g$ on the infinite places with $\fitmi g\le a_1$ such that at every place $p$
the top maps multiply the norms by at most $e^{g(p)}$ and the \emph{conjugated} row maps $\bar k$ multiply the norms by at most
$e^{g(p)}$. Let $\bar p$ be the place of the embedding $p\circ c$, $c$ the complex conjugation of $K$; $p\mapsto\bar p$ is a
bijection, so $g'(p)=g(\bar p)$ satisfies $\fitmi{g'}=\fitmi g\le a_1$. If $v=\sigma(u)$ has norm at most $B$ at $p$, then $\bar v$
has norm at most $B$ at $\bar p$ (the norms are computed on the coefficient vectors, and $\sigma\circ c=\bar\sigma$); applying the
bound for $\bar k$ at $\bar p$ and conjugating back, $k(v)$ has norm at most $e^{g(\bar p)}B=e^{g'(p)}B$ at $p$. This is
$\sfname{InfRow}$ with the function $g'$.
\end{proof}

% ---------------------------------------------------------------------------
\subsection{Slopes of test lattices}\label{sec:pf:fit:testlat}

We prove the slope bound for test lattices used in \cref{pf:fit:leafmm}. Let $K$ be a totally complex number field with an
involution $c$ which is complex conjugation at a place $v_0$, $I$ a finite set (of embedding types), $n\ge1$, and
$\lambda=(\lambda_\tau)_{\tau\in I}$ dominant weights of rank $n$ with \emph{spread} at most $R\ge1$, i.e.\
$\lambda_\tau(i)-\lambda_\tau(j)\le R$ for all $\tau,i,j$. Let $G_\tau\in M_n(K)$ be definite Gram matrices, $\mathcal C$ a non-empty
set and $M(\gamma,\tau)\subseteq K^n$ ($\gamma\in\mathcal C$) $\OO_K$-lattices which are self-dual for $G_\tau$. In the space
$V_\lambda$ of the representations of highest weights $\lambda$ there are the Weyl-type lattices
$\Delta_\lambda(M(\gamma))\subseteq\nabla_\lambda(M(\gamma))$ \cite{LeanSource2026}. For an enumeration
$e\colon\{1,\dots,N\}\to\mathcal C\times(\text{basis of }V_\lambda)$ and lattices $T(\gamma)\subseteq V_\lambda$, the \emph{sum
lattice} $\Sigma(T)\subseteq K^N$ consists of the vectors whose $\gamma$-components lie in $T(\gamma)$, and $G_\Sigma$ is the
Gram matrix of the sum of the pairings induced by the $G_\tau$. A lattice $L\subseteq K^N$ with a metric $X$ satisfies the
\emph{test lattice condition} with bound $B$ if $L$ is finitely generated and non-zero, $X(v)$ is invertible at every infinite place
$v$, and $\fitmumax(L;X)\le B$. Put $B(f,n,R)=4fnR\log(nR)$.

\begin{theorem}[Slopes of sums of Weyl lattices]\label{pf:fit:testlat}
In this situation, with $f=|I|$ and $B=B(f,n,R)$:
\begin{enumerate}[label=(\alph*)]
\item if $X$ is a metric of $G_\Sigma$, then $\Sigma(\nabla_\lambda(M))$ and $\Sigma(\Delta_\lambda(M))$ satisfy the test lattice
condition with bound $B$;
\item if $X$ is a metric of ${}^tG_\Sigma^{-1}$, then the dual of $\Sigma(\Delta_\lambda(M))$ satisfies the test lattice condition
with bound $B$;
\item the same conclusions hold for the lattices obtained from self-dual lattices $L(\gamma)$ of a hermitian matrix $J$ over a CM
field $E$ (positive definite at all embeddings) through embeddings $\epsilon_\tau\colon E\to K$ compatible with the involutions,
with the Gram matrices $(\epsilon_\tau(J))_\tau$ and their duals.
\end{enumerate}
\formallabel{GlobalTheory.TestLat.slope_holds}
\formallabel{GlobalTheory.TestLat.slopeK_holds}
\formallabel{GlobalTheory.TestLat.slopeE_holds}
\formallabel{GlobalTheory.T.maxSlope_sum_weyl_and_dualWeyl_lattices_le}
\formallabel{GlobalTheory.TestLat.testLatticeCondition_sumLat_nabK_and_delK}
\end{theorem}

\begin{proof}
(a) There is an integer $s\ge1$ with $\log s\le B$ such that $s$ times the Weyl pairing of two vectors of
$\nabla_\lambda(M(\gamma))$ is integral for every self-dual $M(\gamma)$: the normalising constant of the pairing satisfies the
bound (\lref{GlobalTheory.T.pairing_normalising_constant_log_le_bound}), and the integrality is
\lref{GlobalTheory.T.extNorm_mul_weylPairing_integral_selfDual}. Let $T(\gamma)$ be lattices with
$T(\gamma)\subseteq\nabla_\lambda(M(\gamma))$; the lattices $\nabla_\lambda(M(\gamma))$ and $\Delta_\lambda(M(\gamma))$ are lattices
and $\Delta_\lambda\subseteq\nabla_\lambda$ (\lref{GlobalTheory.T.nablaDeltaOperator_lattice_preserved}). The sum lattice
$\Sigma(T)$ is a lattice of $K^N$, in particular finitely generated and non-zero
(\lref{GlobalTheory.T.orthogonalSumCoefficientLattice_isLattice_and_dual}), and $G_\Sigma$ is definite with values
$x^*G_\Sigma y=\sum_\gamma\langle x_\gamma,y_\gamma\rangle_\gamma$ (\lref{GlobalTheory.T.sumPairingGramMatrix_values_and_definite}).
Hence $s\,x^*G_\Sigma y\in\OO_K$ for $x,y\in\Sigma(T)$. A non-zero lattice on which $s$ times a definite hermitian form with metric
$X$ is integral has maximal slope at most $\log s$ (\lref{GlobalTheory.T.maxSlope_le_log_of_integral_hermitian_form}). So
$\fitmumax(\Sigma(T);X)\le\log s\le B$; apply this to $T=\nabla_\lambda(M)$ and to $T=\Delta_\lambda(M)$.

(b) is \lref{GlobalTheory.T.maxSlope_dual_sum_weyl_lattices_le} (the dual of $\Sigma(\Delta_\lambda(M))$ for the dual metric).

(c) The matrices $\epsilon_\tau(J)$ and their duals are definite Gram matrices, and the transported lattices are self-dual for
them (\lref{GlobalTheory.T.defGram_selfDual_after_field_embedding}); the lattices of (c) are the lattices of (a) and (b)
for these data, so (c) follows from (a) and (b).
\end{proof}

% ---------------------------------------------------------------------------
\subsection{Remarks on the formal statements}\label{sec:pf:fit:remarks}

\begin{remark}[Unused hypotheses]\label{pf:fit:rem:unused}
\begin{enumerate}[label=(\alph*)]
\item The hypothesis $\sfname{DIVinN}(\mathcal F)$ is not used: it occurs in $\sfname{ConstantFunctionFits}$ and is
carried through the leaf statements, but none of the leaves, nor $\sfname{ExistsFunctionIn}$, invokes it.
\item The bound $|\mathrm{logDer}(1)|\le x^\eps$ ($\sfname{LogDer}$) is one of the conditions of the parameter choices of
the families; the hypothesis $F$ of its formal statement is not used.
\item In the formal proof of the kernel leaf the hypothesis $\sfname{LinK}$ is not used (it is needed only for the unique map
leaf).
\item The formal statement of the archimedean leaf takes $\sfname{ArchSt}$ and $\sfname{BotSmall}$ as hypotheses but does
not use them; the bound is obtained from the grid statement \cref{pf:fit:archgrid}, whose proof uses the archimedean bound for
norms of pairs proved independently (\cref{pf:fit:archnorm}).
\item In \cref{pf:fit:sqtop} the hypothesis of Howe duality at the real place is also a component of $F$.
\item The formal statement of the ratio formula \cref{pf:fit:petratio}(a) takes $\sfname{LNe}$ as a hypothesis but does not use it;
the formula holds for all data satisfying the standing conditions.
\end{enumerate}
\end{remark}

\begin{remark}[The constants]\label{pf:fit:rem:const}
In \cref{pf:fit:flb} the constant $C$ of the leaf conditions and the constant $C_0$ of the numbers $N_Z(C_0,C_H)$ are
different: the entry $hh=x^{C_H}b$ of the numbers does not depend on $C_0$, so $\fitFUN$ holds for all values of the constant
of the numbers, which is the form required in $\sfname{ConstantFunctionFits}$. All numerical constants in the leaf conditions
are explicit ($10^5$ in $\fitpol$, $1.3\cdot10^{10}$ in \cref{pf:fit:leafnum}, $42$, $45$, $80$); the requirement
$C_H\ge16(\lceil|C_t|\rceil+10)+1$ is part of the largeness of the families in the sense of \cref{def:families:closed}.
\end{remark}

\begin{remark}[Where the hypotheses enter]\label{pf:fit:rem:where}
$\sfname{LambdaLowerBound}_2$ enters only through the leaf \ref{pf:fit:F:IPb}, that is, through lower bounds for the scalars
$\lambda_{j,\mu_0,p,v}$ at the base type, which by \cref{pf:fit:ip} control the height of $\Psi$ at the infinite places.
$\sfname{FunctionalsSeparate}$ enters only through $\sfname{InjK}$, that is, through the injectivity of $\Psi$ and the kernel
of $V_{\mu_0}$. Of the components of $\sfname{JoinOut}$, $\sfname{ToricDenominator}$ enters through $\sfname{ThbSt}$ and
$\sfname{RhoLow}$, $\sfname{Transport}$ through $\sfname{TorEx}$ and $\sfname{BetaArch}$, $\sfname{TwistedTransport}$ through
the toric estimate $\sfname{ThA}$, $\sfname{ZetaDenominators}$ through $\sfname{DenEx}$ and $\sfname{DenSum}$, and
$\sfname{Refl}$ through $\sfname{OneSqSt}$.
\end{remark}

\begin{remark}[Name]\label{pf:fit:rem:name}
Despite the names of the statements (``Fourier coefficient fit'', ``constant function fits''), the conclusion is an Arakelov
upper bound for the maximal slope of a hermitian lattice; the ``fit'' refers to the function datum attached to a grid whose
values on the base type are the class sums which, by $\sfname{FormulaK}$, coincide up to a constant and roots of unity with
theta pairings.
\end{remark}

